% concatenated from V3.tex
\documentclass[reqno]{amsart}
\usepackage{amssymb,amsmath,amsthm,mathtools,mathrsfs,xcolor}
\usepackage[english]{babel}
\usepackage{cite}
%\usepackage{showlabels}
\RequirePackage[colorlinks,citecolor=blue,urlcolor=blue]{hyperref}

\DeclareMathOperator{\Var}{Var}
\DeclareMathOperator{\Cov}{Cov}
\DeclareMathOperator{\diam}{diam}
\DeclareMathOperator{\dist}{dist}
\newcommand{\R}{\mathbb R}
\newcommand{\N}{\mathbb N}
\newcommand{\E}{\mathbb E}
\renewcommand{\P}{\mathbb P}
\renewcommand{\d}{\mathrm d}
\newcommand{\e}{\mathrm e}
\newcommand{\1}{\mathbf 1}
\newcommand{\cH}{\mathcal H}
\newcommand{\cB}{\mathcal B}
\newcommand{\inner}[2]{\langle #1,#2\rangle}
\renewcommand{\leq}{\leqslant}
\renewcommand{\geq}{\geqslant}

\numberwithin{equation}{section}
\newtheorem{theorem}{Theorem}[section]
\newtheorem{proposition}[theorem]{Proposition}
\newtheorem{lemma}[theorem]{Lemma}
\newtheorem{corollary}[theorem]{Corollary}
\theoremstyle{definition}
\newtheorem{definition}[theorem]{Definition}
\newtheorem{remark}[theorem]{Remark}

\title[Stochastic heat equation]{Sharp regularity and small ball probabilities for the stochastic heat equation on bounded domains}

\author[J. Hu and C. Y. Lee]{Jingwu Hu \and Cheuk Yin Lee}
\address{Department of Statistics and Probability, Michigan State University, East Lansing, MI 48824, United States}
\email{hujingwu@msu.edu}
\address{School of Science and Engineering, The Chinese University of Hong Kong (Shenzhen), Longgang, Shenzhen, Guangdong 518172, China}
\email{leecheukyin@cuhk.edu.cn}

\subjclass[2020]{60H15; 60G15; 60G17; 60G60}
\keywords{Stochastic heat equation; Gaussian noise; H\"older regularity; strong local nondeterminism; modulus of continuity; law of the iterated logarithm; small ball probability}

\begin{document}

\begin{abstract}
We consider the stochastic heat equation
\[
	\partial_t u(t,x) = \Delta u(t,x) + \dot{W}_\alpha(t,x)
\]
on a bounded Lipschitz domain with zero Dirichlet boundary condition and zero initial condition, where $\dot{W}_\alpha$ is a Gaussian noise that is white in time and whose spatial covariance is the kernel of $(-\Delta)^{-\alpha}$ with $\alpha>0$. We prove that a unique pointwise defined mild solution exists if and only if $\alpha>d/2-1$. In this case, if in addition the domain is $C^2$, we also establish spatial and temporal H\"older regularity of the solution.
When $d/2-1<\alpha<d/2$, we show that the H\"older exponents are optimal and obtain exact local and uniform moduli of continuity, a Chung-type law of the iterated logarithm, and sharp small ball probability estimates for the solution.
%we establish a global canonical-metric upper bound and sharp two-sided estimates on compact interior sets, boundary-degenerate strong local nondeterminism, exact local and compact-rectangle uniform moduli of continuity, a Chung-type law of the iterated logarithm, and matching small ball probability estimates. 
\end{abstract}

\maketitle
%\tableofcontents

\section{Introduction}

Let \(D\subset\R^d\) be a bounded domain. 
As usual, a domain is understood to be a nonempty open connected set. 
We fix a parameter $\alpha > 0$ and consider the stochastic heat equation with Dirichlet boundary condition:
\begin{equation}\label{eq:SHE}
	\begin{cases}
		\partial_t u(t,x)=\Delta u(t,x)+\dot W_\alpha(t,x),
    	& t>0,\ x\in D,\\
		u(t,x)=0, & t>0,\ x\in\partial D,\\
		u(0,x)=0, & x\in D.
	\end{cases}
\end{equation}
Formally, the noise $\dot W_\alpha$ is given by $\dot W_\alpha = (-\Delta)^{-\alpha/2} \dot W$, where $\dot W$ is a space-time white noise on $[0,\infty) \times D$ and $-\Delta$ denotes the positive Dirichlet Laplacian on $D$.
More precisely, $\dot W_\alpha$ is defined as a centered generalized Gaussian field on a complete probability space $(\Omega, \mathcal{F}, \P)$, indexed by $(t,x) \in [0,\infty) \times D$, with covariance
\begin{align}\label{W:corr:G}
    \E[\dot W_\alpha(t,x) \dot W_\alpha(s,y)] = \delta_0(t-s) G_{\alpha}(x,y),
\end{align}
where $G_{\alpha}$ is the kernel of the operator $(-\Delta)^{-\alpha}$.
As in standard theory of stochastic partial differential equations (SPDEs), the solution to \eqref{eq:SHE} is interpreted as mild solution \cite{Walsh,Dalang,DPZ}; see \eqref{eq:mild} below.

% The noise \(W_\alpha\) is white in time and spatially correlated through the Dirichlet Laplacian. Let \(A=-\Delta_D\), and let \(\{(\lambda_n,f_n)\}_{n\ge1}\) be an orthonormal Dirichlet eigenbasis. We use the convention
% \begin{equation}\label{eq:G}
%     G_{2\alpha}(x,y)
%     =
%     \sum_{n=1}^\infty \lambda_n^{-\alpha}f_n(x)f_n(y),
% \end{equation}
% where the series is interpreted as the distributional kernel of \(A^{-\alpha}\). Whenever \(G_{2\alpha}\) has a function representative, this agrees with the corresponding pointwise kernel. Formally,
% \begin{equation}\label{eq:noise}
%     \E[\dot W_\alpha(t,x)\dot W_\alpha(s,y)]
%     =
%     \delta_0(t-s)G_{2\alpha}(x,y)
% \end{equation}
% and the identity is understood in the sense of distributions.

% As is customary, we interpret \eqref{eq:SHE} in its mild form; see \cite{Walsh,Dalang,K09,DK,DPZ}. The Dirichlet heat kernel and the spectral calculus of \(A\) will be used to construct the solution and study its sample functions; see \cite{Davies,McLean} for standard background.



If $\dot{W}_\alpha$ is space-time white noise (i.e., $\alpha=0$), then
\eqref{eq:SHE} has a pointwise solution only when $d=1$.
In this case, sample path properties of the linear and nonlinear stochastic heat equation on a bounded interval have been studied in \cite{HL26}.
In this paper, we will focus on the linear case for $\alpha>0$.
%
%
%
%
%Sharp H\"older regularities are known for linear and nonlinear stochastic heat equations and similar SPDEs on $\R_+ \times \R^d$ driven by Gaussian noise that is either white in space or has spatial covariance given by the Riesz kernel $|x-y|^{-\beta}$, where $0<\beta<d$; see \cite{SS02,CD14,C17,HSWX20,GSWX25, KS23} and the references therein. 
%For linear and nonlinear stochastic heat equations and fractional diffusion equations on $\R_+ \times \R^d$, sharp H\"older regularities have been obtained, e.g., in \cite{Walsh,SS02,CD14,C17,HSWX20,GSWX25, KS23}, where the Gaussian noise is typically either white in space or has spatial covariance given by the Riesz kernel $|x-y|^{-\beta}$, where $0<\beta<d$.
%
When $D=\R^d$, it is common to consider spatially-homogeneous Gaussian noise which is white in space or has spatial covariance given by the Riesz kernel $|x-y|^{-\beta}$, where $0<\beta<d$ (see \cite{Dalang}).
Although the Riesz kernel may also be used on bounded domains \cite{CCL24}, it is often more natural to use spatial covariances that are similar to the one in \eqref{W:corr:G} on general bounded domains, such as $d$-dimensional torus \cite{CCV}, Riemannian manifolds \cite{CO25}, and metric measure spaces \cite{BCHOTW}.
In fact, when $0<\alpha<d/2$, the covariance $G_\alpha(x,y)$ behaves like the Riesz kernel with $\beta = d-2\alpha$ (see Lemma \ref{lem:G:bd} below).
Since the noise $\dot{W}_\alpha(t,x)$ is by definition a generalized Gaussian field and may not be a pointwise function, an important question is to determine when there exists a pointwise solution and determine regularity of the solution.

There is a large literature on H\"older regularity of linear and nonlinear stochastic heat equations and related SPDEs on $\R_+ \times \R^d$; see, e.g., \cite{Walsh,CD14,SS00,SS02,BC18,HL19,BJQ16,BQS19,KS23,C17,L17,DS26}.
In particular, \cite{KS23} studied optimal H\"older regularity of SPDEs on $\R_+ \times \R^d$ with additive Gaussian noise, including the stochastic heat equation.
H\"older and Sobolev regularity of stochastic heat equations and parabolic SPDEs on bounded domains have also been studied in \cite{CCV,DKZ,SV,NV09,DPZ,DS26}, but there are not many explicit results about optimality of H\"older regularity on general domains.
%, but optimality of the H\"older exponents was not established for bounded domains.
%H\"older regularities for the stochastic heat equation with multiplicative noise on the $d$-dimensional torus were studied in \cite{CCV}, 


The goal of this paper is to study sharp H\"older regularity of the solution to \eqref{eq:SHE} on bounded domains and establish further sample path properties including exact moduli of continuity, Chung's law of the iterated logarithm, and sharp small ball probability estimates.

Sample path properties of Gaussian random fields and SPDEs have been studied extensively. Exact local and uniform moduli of continuity and Chung's laws of the iterated logarithm for Gaussian processes and Gaussian random fields have been established in \cite{MWX13,LX23,LZ26,LX10}, while similar sample path results for stochastic heat equations and other related SPDEs can be found in \cite{FKM15,HSWX20,Chen23,GSWX25,WX24,HL26,LX19,CLX26}.
%Exact local and uniform moduli of continuity for Gaussian random fields have been studied in \cite{MWX13, LX23, MR, LZ26}.
%For SPDEs, related results can be found in \cite{FKM15, WX24, GSWX25, HL26, HSWX20, LX19}.
In many cases, the local modulus functions are given by the law of the iterated logarithm (LIL) and the uniform modulus functions are similar to L\'evy's modulus of continuity for Brownian motion. 
These moduli of continuity specify the upper envelope ($\limsup$) for the increments at small scales.
The lower envelope ($\liminf$) for the increments is usually characterized by Chung's LIL, originally due to Chung \cite{Chung} for Brownian motion.
It is a well-known fact that Chung's LIL is closely related to small ball probabilities \cite{LS01}.
%Chung-type LILs have been studied in \cite{LS01,LX23} for Gaussian processes and Gaussian random fields, and
We refer to \cite{AJM21,C24,C25,C26,KKM24,H24,FJK23,HL26} for small ball probability results for stochastic heat equations and related SPDEs.


%Our first goal is to identify the condition for existence of a pointwise random-field solution. For \(\alpha>0\), this condition is
%\begin{equation}\label{eq:Dalang}
%    \alpha>\frac d2-1.
%\end{equation}
%We then study the fine sample path properties of the solution in the regime
%\begin{equation}\label{eq:alpha}
%    \max\left\{0,\frac d2-1\right\}<\alpha<\frac d2.
%\end{equation}
%The condition \(\alpha>0\) allows us to use the heat-kernel representation of the fractional Green kernel. Thus, the one-dimensional space-time white-noise case \(\alpha=0\) is a limiting case and is not covered by our argument. Set

%For the exact sample path results, we also use the coordinate-anisotropic metric

%The two metrics are comparable, but the limiting constants in our exact sample path results are expressed in terms of \(\Delta\), as required by the Lee--Xiao theorems. The critical temporal and spatial H\"older exponents are
%\(\theta\) and \(2\theta\), respectively. At the formal endpoint \(\alpha=0\) in dimension one, these exponents reduce to \(1/4\) and \(1/2\). The corresponding metric dimension is


%Sample path regularities of Gaussian random fields and SPDEs have been studied extensively. Exact moduli of continuity and Chung's laws of the iterated logarithm for anisotropic Gaussian fields were developed in \cite{MWX13,LX23}, while sharp regularity results for stochastic heat equations and related SPDEs were established by \cite{FKM15,HSWX20,Chen23,LX25,GSWX25}. Strong local nondeterminism is one of the main tools behind these results; see \cite{Berman,Berman87,Pitt,Xiao08,CLX26}.
%small ball probabilities for Gaussian processes are surveyed in \cite{LS01}. For stochastic heat equations and SPDEs, small ball estimates and related exceptional flat-point problems have been studied in \cite{AJM21,Chen24,Chen25,KKM24}. 

%Our small ball estimate combines canonical metric estimates, strong local nondeterminism, Anderson's inequality \cite{Anderson}, and the Talagrand entropy lower bound for Gaussian processes \cite[Lemma~4.1]{LX23}. It is stated for cylinders \(I=[a,T]\times A\), where \(A\subset D\) is open and need not be bounded away from \(\partial D\). The proof combines an outer entropy estimate with an interior packing argument.

%The closest whole-space results for the present noise class are those of Guo, Song, Wang, and Xiao \cite{GSWX25}, who study sample paths and small ball probabilities for stochastic fractional diffusion equations, and Chen, Lee, and Xia \cite{CLX26}, who establish strong local nondeterminism for time-fractional diffusion equations. 

%Recently, \cite{GSWX25} and \cite{CLX26} studied sample path properties and small ball probabilities for stochastic fractional diffusion equations on $\R^d$ driven by space-time fractional colored noise.


%The bounded-domain Dirichlet problem considered here has a different spectral covariance, is not translation invariant, and develops a genuine boundary degeneracy. Our contribution is to quantify that degeneracy through the first Dirichlet eigenfunction, to prove the required spectral-duality estimates on a general bounded \(C^2\)-domain, and to obtain full-domain small ball estimates by combining global entropy bounds with an interior approximation argument.

%This paper is also closely related to the Gaussian analysis in \cite{HL26}, where nonlinear parabolic SPDEs and the open KPZ equation on bounded intervals are studied. Here we consider the additive Gaussian equation on a general bounded domain. We prove the required analytic estimates directly and then apply the Gaussian-field results of Lee and Xiao \cite{LX23}. The exact uniform modulus is obtained on compact coordinate rectangles. Near the Dirichlet boundary, the variance and the SLND lower bound both degenerate.
%; the corresponding full-domain exact uniform modulus remains open.

\subsection{Main results}

%We first state the existence and covariance result.
We first present a necessary and sufficient condition for the existence of a unique random-field solution to \eqref{eq:SHE}. 
Such condition is also known as Dalang's condition \cite{Dalang}.
According to standard SPDE theory \cite{Walsh, Dalang, DPZ}, \eqref{eq:SHE} has a unique pointwise mild solution given by
\begin{equation}\label{eq:mild}
    u(t,x)=\int_0^t\int_D P_{t-r}(x,y)\,W_\alpha(\d r,\d y)
\end{equation}
if and only if the above Wiener integral process is well defined pointwise, where $P$ denotes the Dirichlet heat kernel.

\begin{theorem}[Dalang condition]
\label{thm:exist}
Suppose $D \subset \R^d$ is a bounded Lipschitz domain.
Then \eqref{eq:mild} is pointwise defined for every $(t,x) \in (0,\infty)\times D$ if and only if
\begin{align}\label{Dalang}
	\alpha>d/2-1.
\end{align}
\end{theorem}


%\begin{theorem}[Existence and covariance representation]\label{thm:exist}
%% Let \(D\subset\R^d\) be bounded and let \(\alpha>0\). 
%The stochastic integral 
%\begin{equation}\label{eq:mild}
%    u(t,x)=\int_0^t\int_D P_{t-r}(x,y)\,W_\alpha(\d r,\d y)
%\end{equation} 
%is well defined for every \(t>0\) and \(x\in D\) if and only if \(\alpha>d/2-1\). In this case \(u\) is a centered Gaussian random field with covariance
%\begin{align}\label{eq:cov}
%    \E[u(t,x)u(s,y)]
%    &=
%    \int_0^{t\wedge s}
%    \int_D\int_D
%    P_{t-r}(x,z)P_{s-r}(y,w)G_{2\alpha}(z,w)\,\d z\,\d w\,\d r  \\
%    &=
%    \sum_{n=1}^\infty
%    \lambda_n^{-\alpha}f_n(x)f_n(y)
%    \int_0^{t\wedge s}
%    \e^{-\lambda_n(t-r)}\e^{-\lambda_n(s-r)}\,\d r .
%    \notag
%\end{align}
%The first line of \eqref{eq:cov} is understood as the quadratic-form pairing
%associated with \(A^{-\alpha}\) whenever \(G_{2\alpha}\) does not have a
%function representative.
%Equivalently,
%\begin{equation}\label{eq:spectral}
%    u(t,x)=
%    \sum_{n=1}^\infty
%    \lambda_n^{-\alpha/2}
%    \left(\int_0^t \e^{-\lambda_n(t-r)}\,\d\beta_n(r)\right)f_n(x),
%\end{equation}
%where \(\{\beta_n\}_{n\ge1}\) are independent standard Brownian motions.
%\end{theorem}

%For the remaining results, suppose that \(D\) is a bounded \(C^2\)-domain and that \eqref{eq:alpha} holds. The first Dirichlet eigenfunction \(f_1\) is chosen positive and \(L^2(D)\)-normalized; the principal eigenvalue \(\lambda_1>0\) is simple because \(D\) is connected. All pathwise statements refer to the separable jointly continuous version constructed in Proposition~\ref{prop:continuous}. Since \(0<\theta<1/2\), both \(\rho\) and \(\Delta\) are metrics. Recall that
%\begin{equation}\label{eq:Qsum}
%    Q=\frac1\theta+\frac d{2\theta}
%    =
%    \frac{2(d+2)}{2\alpha+2-d}.
%\end{equation}

We obtain H\"older regularity for the solution when in addition $D$ is $C^2$.

\begin{theorem}[H\"older regularity]\label{thm:reg}
Suppose $D \subset \R^d$ is a bounded $C^2$ domain and \eqref{Dalang} holds.
Then $\{u(t,x)\}_{t \ge 0, x \in D}$ has a version that is locally H\"older continuous of order $\beta_0$ in $t$ and order $\beta_1$ in $x$ for any $\beta_0 \in (0, \frac{\theta\wedge 1}{2})$
and $\beta_1 \in (0,\theta\wedge 1)$, where
\begin{align}\label{theta}
	\theta=\alpha-d/2+1.
\end{align}
%If $\alpha>d/2$, then $u(t,x)$ is locally H\"older of order $\frac12 -$ in $t$ and differentiable in $x$.
\end{theorem}

Next, we establish the optimality of the H\"older exponents and sharp moduli of continuity under the condition
\begin{align}\label{alpha:Holder:range}
	d/2-1<\alpha<d/2,
\end{align}
which is equivalent to $0<\theta<1$.
In this case, we define the metric $\rho$ on $[0,\infty) \times D$ by
\begin{equation}\label{eq:rho}
    \rho((t,x),(s,y))
    =|t-s|^{\theta/2}+\sum_{i=1}^d |x_i-y_i|^{\theta},
\end{equation}
for any $(t,x) = (t,x_1,\dots, x_d), (s,y) = (s,y_1,\dots, y_d) \in [0,\infty)\times D$.
Note that $\rho$ is equivalent to the metric $|t-s|^{\theta/2} + |x-y|^\theta$.
The next theorem gives sharp upper and lower bounds for the variance of increments including exact dependence near the boundary and at small $t$, and shows that the solution satisfies the strong local nondeterminism property under condition \eqref{alpha:Holder:range}.
This implies that the H\"older exponents in Theorem \ref{thm:reg} are optimal under this condition.
We say that a centered Gaussian random field $\{ X(s) \}_{s \in I \subset \R^N}$ satisfies strong local nondeterminism (SLND) if there exists $C>0$ such that
\begin{align}\label{SLND}
	\Var(X(s)\mid X(s_1),\dots, X(s_n)) \ge C \min_{1\le i\le n} \Var(X(s)-X(s_i)),
\end{align}
uniformly for all $n \in \N_+$ and $s,s_1,\dots, s_n \in I$ \cite{Berman, Pitt, CD82, MP87, Xiao08}.
SLND has been a useful tool for studying various sample properties of Gaussian processes and Gaussian random fields; see \cite{X09,Xiao08} and the references therein.
One of the main contributions of this paper is showing the SLND property for the solution.

\begin{theorem}\label{thm:optimal}
Suppose $D \subset \R^d$ is a bounded $C^2$ domain and \eqref{alpha:Holder:range} holds.
Then for any $T>0$,
%\begin{align}\label{var:u-u}
%	\Var(u(t,x)-u(s,y)) \asymp \left[\rho((t,x),(s,y)) \wedge (t\vee s)^{\theta/2} \wedge (d(x)\vee d(y))^\theta\right]^2
%\end{align}
\begin{align}\label{var:u-u}
	\Var(u(t,x)-u(s,y)) \asymp \left[\rho((t,x),(s,y)) \wedge \left((\sqrt{t} \wedge d(x))^\theta \vee (\sqrt{s} \wedge d(y))^\theta\right)\right]^2
\end{align}
uniformly for all $t,s \in [0,T]$ and $x,y \in D$, where $\theta$ is given by \eqref{theta}, $\rho$ is given by \eqref{eq:rho}, and $d(x)$ is the distance-to-boundary defined by
\begin{align}\label{d(x)}
	d(x):= \mathrm{dist}(x,\partial D).
\end{align}
In particular, setting $s=0,y=x$ yields $\Var(u(t,x)) \asymp t^{\theta} \wedge (d(x))^{2\theta}$ on $[0,T]\times D$.
Moreover, for any $T>\delta>0$ and compact subset $D' \subset D$, $\{u(t,x)\}_{(t,x) \in [\delta,T]\times D'}$ satisfies SLND in the sense of \eqref{SLND} with $C$ depending on $\delta,T,D'$.
\end{theorem}


%\begin{theorem}[Strong local nondeterminism]\label{thm:SLND}
%Assume that \(D\subset\R^d\) is a bounded \(C^2\)-domain and that \eqref{eq:alpha} holds. Then for every \(T>0\),
%\begin{equation}\label{eq:SLND}
%    \Var(u(t,x)-u(s,y))
%    \le C\rho((t,x),(s,y))^2
%\end{equation}
%uniformly for \((t,x),(s,y)\in(0,T]\times D\). In particular, for every \(a>0\) and \(\delta>0\),
%\[
%    \Var(u(t,x)-u(s,y))\asymp \rho((t,x),(s,y))^2,
%    \qquad (t,x),(s,y)\in[a,T]\times D_\delta ,
%\]
%whenever \(D_\delta\ne\varnothing\).
%The two-sided estimate therefore holds on compact interior sets. Moreover, there exists
%\(c=c(D,d,\alpha,T)>0\) such that for all \(n\ge1\),
%\[
%    \Var\left(
%        u(t,x)\mid u(t_1,x_1),\ldots,u(t_n,x_n)
%    \right)
%    \ge
%    c\left[
%        \min_{1\le j\le n}\rho((t,x),(t_j,x_j))^2
%        \wedge t^{\theta}
%        \wedge f_1(x)^{2\theta}
%    \right]
%\]
%uniformly for \((t,x),(t_j,x_j)\in(0,T]\times D\).
%\end{theorem}


%Define
%\begin{equation}\label{eq:Delta}
%    \Delta((t,x),(s,y))
%    =|t-s|^{\theta/2}+\sum_{i=1}^d|x_i-y_i|^{\theta}.
%\end{equation}

As a consequence of the optimal regularity and the SLND property, we obtain sharp local and uniform moduli of continuity, Chung's law of the iterated logarithm, and sharp small ball probability estimates for the solution.
To present these results, let us define
\begin{align}
	&B_\rho((t,x),r) := \{ (s,y) \in [0,\infty) \times D : \rho((t,x),(s,y)) \le r \},\\
	&B^*_\rho((t,x),r) := \{ (s,y) \in [0,\infty) \times D : 0< \rho((t,x),(s,y)) \le r\}
\end{align}
for any $(t,x) \in [0,\infty) \times D$ and $r>0$.

\begin{theorem}[Exact local modulus of continuity]\label{thm:moc}
Suppose \(D\) is a bounded \(C^2\) domain and
\eqref{alpha:Holder:range} holds. 
Then, for every
\(z_0=(t_0,x_0)\in(0,\infty)\times D\), there exists a nonrandom number 
\(K_0 \in(0,\infty)\) depending on $z_0$ such that
\[
    \lim_{\varepsilon\to0^+}
    \sup_{z\in B_\rho^*(z_0,\varepsilon)}
    \frac{|u(z)-u(z_0)|}
    {\rho(z,z_0)\sqrt{\log\log(1/\rho(z,z_0))}}
    =
    K_0
    \qquad\text{a.s.}
\]
\end{theorem}

\begin{theorem}[Exact uniform modulus of continuity]\label{thm:umoc}
Suppose \(D\) is a bounded \(C^2\) domain and
\eqref{alpha:Holder:range} holds. 
%For every nondegenerate closed coordinate rectangle \(R\Subset(0,\infty)\times D\), 
Then, for any $T>0$, there exists a nonrandom number \(K_1 \in(0,\infty)\) depending on $T$ such that
\begin{align}\label{umoc}
    \lim_{\varepsilon\to0^+}
    \sup_{\substack{z,z'\in [0,T]\times D\\0<\rho(z,z')\le\varepsilon}}
    \frac{|u(z)-u(z')|}
    {\rho(z,z')\sqrt{\log(1/\rho(z,z'))}}
    =
    K_1
    \qquad\text{a.s.}
\end{align}
\end{theorem}


\begin{theorem}[Chung's law of the iterated logarithm]\label{thm:lil}
Suppose \(D\) is a bounded \(C^2\) domain and
\eqref{alpha:Holder:range} holds. 
Then, for every \(z_0=(t_0,x_0)\in(0,\infty)\times D\), there exists a nonrandom number \(K_2 \in(0,\infty)\) depending on $z_0$ such that
\[
    \liminf_{\varepsilon\to0^+}
    \frac{
        \sup_{z\in B_\rho(z_0,\varepsilon)}
        |u(z)-u(z_0)|
    }{
        \varepsilon(\log\log(1/\varepsilon))^{-1/Q}
    }
    =
    K_2
    \qquad\text{a.s.,}
\]
where
\begin{equation}\label{eq:Q}
    Q=\frac{2}{\theta}+ \sum_{i=1}^d\frac{1}{\theta}
    =
    \frac{2+d}{\alpha-d/2+1}.
\end{equation}
\end{theorem}


\begin{theorem}[small ball probability]\label{thm:sbp}
Suppose \(D\) is a bounded \(C^2\) domain and \eqref{alpha:Holder:range} holds. Then, for any $T>0$, there exist \(K_3,K_4>0\) depending on $T$ such that
\[
    \exp\left(-K_3\varepsilon^{-Q}\right)
    \le
    \P\left\{
        \sup_{(t,x)\in [0,T]\times D}|u(t,x)|\le\varepsilon
    \right\}
    \le
    \exp\left(-K_4\varepsilon^{-Q}\right),
\]
uniformly for all \(\varepsilon \in (0, 1]\), where $Q$ is given by \eqref{eq:Q}.
%In particular, \(A=D\) is allowed, and the exponent \(Q\) is sharp. The restriction \(a>0\) is needed in the interior packing argument.
\end{theorem}

Matching small ball probability bounds for nonlinear stochastic heat equations driven by space-time white noise on a bounded interval have been established in \cite{AJM21}.
The case of nonlinear stochastic heat equations on torus and compact Riemannian manifolds driven by temporally-white, spatially-colored Gaussian noise have also been studied in \cite{C24,C25,C26}.
However, in \cite{C24,C25,C26}, the exponents in the upper and lower bounds do not match.
In Theorem \ref{thm:sbp} above, we are able to obtain matching bounds for the linear stochastic heat equation \eqref{eq:SHE} driven by colored noise of a similar type on bounded $C^2$ domains, thanks to our SLND property.

Theorems \ref{thm:moc}--\ref{thm:sbp} are established using the framework of \cite{LX23}, which relies on sharp variance bounds for the increments and the SLND property.
Sharp variance bounds in the case of bounded domains are more involved compared to the case of bounded intervals in \cite{HL26}.
In particular, the Dirichlet eigenfunctions are not uniformly bounded unless $d=1$ and uniform bounds of the Dirichlet eigenvalues and eigenfunctions are not good enough to derive those sharp variance bounds in higher dimensions.
It turns out that sharp variance estimates require additional tools such as Gaussian-type heat kernel estimates.
Also, the domain is required to be sufficiently smooth (with $C^2$ boundary) in order to apply gradient estimates for the heat kernel to obtain sharp spatial regularity.
Similarly to \cite{HL26}, the SLND property is established using the orthonormal basis of eigenfunctions and a scaling argument involving suitable test functions.

Since the temporal process $\{u(t,x_0)\}_{t \in [\delta,T]}$ and the spatial process $\{ u(t_0,x) \}_{x \in D'}$ also satisfy SLND, the following results can be obtained using the same proofs for the above main theorems.

\begin{corollary}
Suppose $D\subset \R^d$ is a bounded $C^2$ domain and \eqref{alpha:Holder:range} holds.
Then, for any fixed $T \ge t_0>0$ and $x_0 \in D$, there exist constants $C_0,C_0',C_1,C_1',C_2,C_2' \in (0,\infty)$ such that
\begin{gather*}
	\lim_{\varepsilon\to0^+} \sup_{t \ge 0:0< |t-t_0| \le \varepsilon} \frac{|u(t,x_0)-u(t_0,x_0)|}{|t-t_0|^{\theta/2} \sqrt{\log\log(1/|t-t_0|)}} = C_0 \quad \text{a.s.},\\
	\lim_{\varepsilon\to0^+} \sup_{x \in D: 0< |x-x_0| \le \varepsilon} \frac{|u(t_0,x)-u(t_0,x_0)|}{|x-x_0|^\theta \sqrt{\log\log(1/|x-x_0|)}} = C_0' \quad \text{a.s.},\\
	\lim_{\varepsilon\to0^+} \sup_{t,t' \in [0,T]: 0<|t-t'| \le \varepsilon} \frac{|u(t,x_0)-u(t',x_0)|}{|t-t'|^{\theta/2} \sqrt{\log(1/|t-t'|)}} = C_1 \quad \text{a.s.},\\
	\lim_{\varepsilon\to0^+} \sup_{x,x'\in D: 0<|x-x'| \le \varepsilon} \frac{|u(t_0,x)-u(t_0,x')|}{|x-x'|^\theta \sqrt{\log(1/|x-x'|)}} = C_1' \quad \text{a.s.},\\
	\liminf_{\varepsilon\to0^+} \frac{\sup_{t: |t-t_0| \le \varepsilon}|u(t,x_0)-u(t_0,x_0)|}{\varepsilon^{\theta/2} (\log\log(1/\varepsilon))^{-\theta/2}} = C_2 \quad \text{a.s.},\\
	\liminf_{\varepsilon\to0^+} \frac{\sup_{x \in D: |x-x_0| \le \varepsilon}|u(t_0,x)-u(t_0,x_0)|}{\varepsilon^\theta (\log\log(1/\varepsilon))^{-\theta/d}} = C_2' \quad \text{a.s.}
\end{gather*}
Moreover, there exist $C_3,C_3',C_4,C_4' \in (0,\infty)$ such that
\begin{gather*}
	\exp(-C_3 \varepsilon^{-2/\theta}) \le \P\left\{ \sup_{t \in [0,T]} |u(t,x_0)| \le \varepsilon \right\} \le \exp(-C_4 \varepsilon^{-2/\theta}),\\
	\exp(-C_3' \varepsilon^{-d/\theta}) \le \P\left\{ \sup_{x \in D} |u(t_0,x)| \le \varepsilon \right\} \le \exp(-C_4' \varepsilon^{-d/\theta}),
\end{gather*}
uniformly for all $\varepsilon \in (0,1]$.
\end{corollary}

%The boundary factor in Theorem~\ref{thm:optimal} has a direct geometric meaning.
%On a bounded connected \(C^2\)-domain, the positive first Dirichlet
%eigenfunction satisfies \(f_1(x)\asymp d(x)\). Consequently the same distance
%scale that forces the variance to vanish at the boundary also caps the SLND
%lower bound. On compact interior rectangles this cap is uniformly positive
%and one recovers the usual nondegenerate SLND regime; globally, the displayed
%factor records precisely how that regime deteriorates near \(\partial D\).


The rest of the paper is organized as follows. 
In Section \ref{s:pre}, we recall some basic facts about the Dirichlet heat kernel and fractional Laplacian, and give the precise definition of the Gaussian noise and corresponding stochastic integrals.
In Section \ref{s:exist}, we prove Theorem \ref{thm:exist}, derive variance bounds, and give a series representation for the solution.
In Section \ref{sec:regularity}, we prove Theorem \ref{thm:reg} and obtain temporal and spatial regularity of the solution.
In Section \ref{s:optimal}, we show that the solution satisfies SLND and prove Theorem \ref{thm:optimal}.
In Section \ref{s:proofs}, we prove Theorems \ref{thm:moc}--\ref{thm:sbp}.
%In Section 2, we construct the solution on a bounded domain and prove the covariance representation. Section 3 is devoted to increment estimates, boundary behavior, strong local nondeterminism, and exact moduli of continuity on bounded connected \(C^2\)-domains. In Section 4, we establish the small ball probability estimates. The analytic and spectral estimates used in the proofs are collected in the appendix.

\subsection{Notation}
We clarify some notation that will be used throughout the paper:
\(\mathbb N_+=\{1,2,\ldots\}\); \(\mathbb N_0=\{0,1,2,\ldots\}\);
\(\R_+=(0,\infty)\);
%For \(x\in D\), set
%\[
%    d(x)=\dist(x,\partial D),\qquad
%    D_\delta=\{x\in D:d(x)\ge\delta\},\qquad
%    \partial_\delta D=\{x\in D:d(x)<\delta\}.
%\]
%Let \(\mathcal T=[0,\infty)\times\overline D\). For
%\(\sigma\in\{\rho,\Delta\}\), \(z_0\in\mathcal T\), and \(r>0\), write
%\[
%    B_\sigma(z_0,r)
%    =\{z\in\mathcal T:\sigma(z,z_0)\le r\},
%    \qquad
%    B_\sigma^*(z_0,r)
%    =B_\sigma(z_0,r)\setminus\{z_0\}.
%\]
For a set \(A\), \(\#A\) denotes its cardinality and \(\1_A\) its indicator
function; \(a\wedge b=\min\{a,b\}\),
\(a\vee b=\max\{a,b\}\); $\log_+(x) = \log(e \vee x)$, where $\log$ denotes natural logarithm; \(c,C\) denote finite positive constants
whose values may change from line to line; For two functions \(f\)
and \(g\), ``\(f\lesssim g\)'' means there is $C>0$ such that \(f(x)\le Cg(x)\) for all $x$; ``\(f\asymp g\)'' means
\(f\lesssim g\) and \(g\lesssim f\); ``\(f(x)=O(g(x))\)'' means \(|f|\lesssim |g|\); For any $k \ge 1$, $\|X\|_k = (\E|X|^k)^{1/k}$ denotes the $L^k(\Omega)$-norm of a random variable $X$; $B(x,r)$ and $\overline{B}(x,r)$ respectively denote open and closed Euclidean balls centered at $x$ with radius $r$; $|x|$ denotes the Euclidean norm of $x$.

\section{Preliminaries}\label{s:pre}

\subsection{Dirichlet Heat kernel}

Let $D\subset \R^d$ be a bounded Lipschitz domain.
According to standard spectral theory \cite{Davies,McLean}, the Dirichlet Laplacian $-\Delta$ has a discrete spectrum, with strictly positive eigenvalues $0<\lambda_1 < \lambda_2 \le \lambda_3 \le \dots$ and an orthonormal basis of eigenfunctions $\{f_n\}_{n=1}^\infty$ for $L^2(D)$, with respect to the inner product $\langle \phi\,,\psi\rangle_{L^2(D)} = \int_D \phi(x) \psi(x)\, \d x$, such that $-\Delta f_n = \lambda_n f_n$ in $D$ and $f_n=0$ on $\partial D$.
Moreover, $f_1$ can be chosen to be positive \cite[Proposition 1.4.3]{Davies} and $f_n \in H^1_0(D) \cap C^\infty(D)$ \cite[Theorem 9.31]{Brezis}.
The Dirichlet heat kernel is given by
\begin{align}\label{P}
    P_t(x,y) = \sum_{n=1}^\infty \e^{-\lambda_n t} f_n(x) f_n(y) \qquad \forall t > 0, x, y \in D.
\end{align}
Due to Weyl's law $\lambda_n \asymp n^{2/d}$ (see \cite[Corollary 3.15]{FLW}) and $\sup_{x \in D}|f_n(x)|\lesssim \lambda_n^{d/4}$ (see \cite[Proposition 5]{F95}), for any $t>0$ the sum in \eqref{P} is uniformly absolutely convergent for $x,y\in D$.
Also, it follows from \eqref{P} that $P_t$ is symmetric, i.e., $P_t(x,y) = P_t(y,x)$ for all $t>0$ and $x,y\in D$, and $P_t$ has the semigroup property:
\begin{align}\label{semigroup}
    \int_D P_t(x,y) P_s(y,z) \, \d y = P_{t+s}(x,z) \quad \forall t,s>0, x,z \in D.
\end{align}
%Let \(A=-\Delta_D\) be the positive Dirichlet Laplacian. Its eigenvalues are
%\[
%    0<\lambda_1<\lambda_2\le\lambda_3\le\cdots,
%\]
%and we choose a real orthonormal basis \(\{f_n\}_{n\ge1}\) of \(L^2(D)\) such that \(Af_n=\lambda_n f_n\). We use the standard spectral theory of the Dirichlet Laplacian and its heat kernel; see \cite{Davies,McLean}. 
%The Dirichlet heat kernel is given by
%\begin{equation}\label{eq:heat}
%    P_t(x,y)=\sum_{n=1}^\infty \e^{-\lambda_n t}f_n(x)f_n(y),
%    \qquad t>0,\ x,y\in D.
%\end{equation}
%It is symmetric and satisfies the semigroup identity
%\begin{align}\label{eq:semigroup}
%    \int_D P_t(x,z)P_s(z,y)\,\d z=P_{t+s}(x,y).
%\end{align}
Recall a two-sided heat kernel estimate \cite{Riahi}:
\begin{lemma}\label{lem:P}
    If $D \subset \R^d$ is a bounded Lipschitz domain, then there exist constants $c_1,c_2,C_1,C_2>0$ and $0<\mu \le 1 \le \nu$ such that
    \begin{align}\begin{split}\label{P:bd}
        &c_1\left( 1\wedge \frac{f_1(x)}{1 \wedge t^{\mu/2}} \right)\left( 1\wedge \frac{f_1(y)}{1 \wedge t^{\mu/2}} \right) \frac{\e^{-\lambda_1 t}}{1\wedge t^{d/2}}
        \exp\left(-\frac{|x-y|^2}{C_1t}\right)
        \le P_t(x,y) \\
        &
        \le c_2 \left( 1\wedge \frac{f_1(x)}{1 \wedge t^{\nu/2}} \right)\left( 1\wedge \frac{f_1(y)}{1 \wedge t^{\nu/2}} \right) \frac{\e^{-\lambda_1 t}}{1\wedge t^{d/2}}
        \exp\left(-\frac{|x-y|^2}{C_2t}\right)
    \end{split}\end{align}
    uniformly for all $t>0$ and $x,y \in D$. If in addition $D$ is $C^{1,\gamma}$ for some $\gamma>0$, then \eqref{P:bd} holds with $\mu=\nu=1$ and $f_1(x) \asymp d(x)$ for all $x \in D$, where $d(x)$ is the distance-to-boundary given by \eqref{d(x)}.
\end{lemma}

When in addition $D$ is a $C^2$ domain, we also have the following gradient estimate for the heat kernel:

\begin{lemma}\label{lem:grad:P}
If $D \subset \R^d$ is a bounded $C^2$ domain, then there exists $C>0$ such that
\begin{align*}
	|\nabla_x P_t(x,y)| \le \frac{e^{-\frac12 \lambda_1 t}}{t^{(d+1)/2}} \exp\left(-\frac{|x-y|^2}{Ct}\right),
\end{align*}
uniformly for all $t>0$ and $x,y \in D$.
\end{lemma}

\begin{proof}
By Theorem 2.1 of \cite{Z06}, when $D$ is a bounded $C^2$ domain, there exists $c>0$ such that
\begin{align*}
	|\nabla_x \log P_t(x,y)| \le 
	\begin{cases}
	\frac{c}{d(x)} & \text{if $d(x) \le \sqrt{t}$,}\\
	\frac{c}{\sqrt{t}} \left( 1+ \frac{|x-y|}{\sqrt{t}}\right) & \text{if $d(x)>\sqrt{t}$,}
	\end{cases}
\end{align*}
uniformly for all $t \in (0,1]$ and $x,y \in D$.
It follows that
\begin{align*}
	|\nabla_x P_t(x,y)| \le 
	\begin{cases}
	\frac{c}{d(x)}P_t(x,y) & \text{if $d(x) \le \sqrt{t}$,}\\
	\frac{c}{\sqrt{t}} \left( 1+ \frac{|x-y|}{\sqrt{t}}\right)P_t(x,y) & \text{if $d(x)>\sqrt{t}$,}
	\end{cases}
\end{align*}
uniformly for all $t \in (0,1]$ and $x \in D$.
Using Lemma \ref{lem:P} in both cases and using also the inequality $z e^{-z^2} \lesssim e^{-z^2/2}$ for $z \ge 0$ in the  $d(x)>\sqrt{t}$ case, we deduce that
\begin{align}\label{grad:log:P}
	|\nabla_x P_t(x,y)| \lesssim \frac{e^{-\lambda_1 t}}{t^{(d+1)/2}} \exp\left( - \frac{|x-y|^2}{2C_2 t} \right),
\end{align}
uniformly for all $t \in (0,1]$ and $x,y \in D$.
It remains to show the same estimate for $t\ge 1$.
In this case, we use the semigroup property \eqref{semigroup} to write
\[
	\nabla_x P_t(x,y) = \int_D \nabla_x P_{1/2}(x,z) P_{t-1/2}(z,y) \, \d z
\]
for all $t \ge 1$ and $x,y \in D$, where the exchange of $\nabla_x$ and the integral can be justified using the dominated convergence theorem and the gradient estimate for $t=1/2$ established above.
Then, by \eqref{grad:log:P} and Lemma \ref{lem:P}, we have
\begin{align*}
	|\nabla_x P_t(x,y)| &\lesssim e^{-\lambda_1(t-1/2)} \int_D  \exp\left( - \frac{|x-z|^2}{C_2} \right) \, \d z\\
	& \lesssim e^{-\lambda_1 t} \lesssim \frac{e^{-\frac12\lambda_1 t}}{t^{(d+1)/2}} \exp\left( - \frac{|x-y|^2}{2C_2t} \right),
\end{align*}
uniformly for all $t \ge 1$ and $x,y \in D$, where the last inequality holds because $e^{-\lambda_1 t} \lesssim e^{-\frac12 \lambda_1 t}\,t^{-(d+1)/2}$ for $t \ge 1$ and the exponent factor is bounded below for all $t \ge 1$ and $x,y \in D$.
\end{proof}

\subsection{Fractional Laplacian}

For any $\alpha>0$, the fractional Laplacian $(-\Delta)^{-\alpha}$ is the bounded linear operator on $L^2(D)$ defined by
\begin{align}\label{frac:Lapl}
    (-\Delta)^{-\alpha} \phi = \sum_{n=1}^\infty \lambda_n^{-\alpha}\langle \phi,f_n \rangle_{L^2(D)}\, f_n \qquad \forall \phi \in L^2(D).
\end{align}
Following \cite{BC23, BCHOTW}, we define the Schwartz space on $D$ by
\begin{align}\label{S(D)}
	\mathcal{S}(D) = \left\{ \phi \in C_0(D) : \forall k \ge 0, \lim_{n\to \infty} n^k\left|\langle \phi, f_n \rangle_{L^2(D)}\right| = 0 \right\}.
\end{align}

\begin{lemma}
Let $D \subset \R^d$ be a bounded Lipschitz domain.
    For any $\alpha>0$, the operator $(-\Delta)^{-\alpha}$ has a positive, positive definite kernel given by
    \begin{align}\label{G}
        G_{\alpha}(x,y) = \frac{1}{\Gamma(\alpha)} \int_0^\infty r^{\alpha-1} P_r(x,y) \, \d r \quad \forall x, y \in D,
    \end{align}
    in the sense that for any $\phi \in \mathcal{S}(D)$,
    \begin{align}\label{frac:Lapl:ker}
        (-\Delta)^{-\alpha} \phi(x) = \int_D G_\alpha(x,y) \phi(y) \, \d y, \quad \forall x \in D.
    \end{align}
    In particular, for any $\phi, \psi \in \mathcal{S}(D)$, we have
    \begin{align}\label{G:id}
        % \int_D (-\Delta)^{-\beta} \varphi(x) (-\Delta)^{-\beta} \psi(x) \, \d x 
        \big\langle (-\Delta)^{-\alpha/2}\phi \,, (-\Delta)^{-\alpha/2}\psi \big\rangle_{L^2(D)} = \int_D \int_D \phi(y) G_{\alpha}(y,z) \psi(z) \, \d y\, \d z.
    \end{align}   
\end{lemma}

\begin{proof}
Lemma \ref{lem:P} implies that the integral in \eqref{G} is convergent and $G_\alpha(x,y)>0$ for any distinct $x,y \in D$.
The relation \eqref{frac:Lapl:ker} follows from \eqref{G}, \eqref{P}, and the identity 
$\lambda^{-\alpha} = \frac{1}{\Gamma(\alpha)}\int_0^\infty r^{\alpha-1} e^{-\lambda r}\, \d r$ for $\alpha,\lambda>0$.
As the kernel of the positive operator $(-\Delta)^{-\alpha}$, $G_\alpha(x,y)$ is positive definite.
Finally, \eqref{G:id} follows from \eqref{frac:Lapl} and the orthonormality of $\{f_n\}_{n\ge1}$.
%\eqref{G}, \eqref{frac:Lapl:ker} and the semigroup property \eqref{semigroup}.
\end{proof}

\begin{lemma}\label{lem:G:bd}
Let $D \subset \R^d$ be a bounded Lipschitz domain.
For any $\alpha>0$, there exist $C,C'>0$ such that
\[
	C f_1(x)f_1(y) K_{d,\alpha}(|x-y|) \le  G_\alpha(x,y) \le C' K_{d,\alpha}(|x-y|)
\]
for all $x,y \in D$, where $f_1$ is the first Dirichlet eigenfunction and
\[
	K_{d,\alpha}(|x-y|) = \begin{cases}
	|x-y|^{2\alpha-d}& \text{if $0<\alpha<d/2$,}\\
	\log_+(1/|x-y|) & \text{if $\alpha=d/2$,}\\
	1 & \text{if $\alpha>d/2$.}
	\end{cases}
\]
\end{lemma}

\begin{proof}
Let $M$ be the diameter of $D$.
By \eqref{G} and the upper bound in Lemma \ref{lem:P},
\begin{align*}
	&G_\alpha(x,y) 
	\lesssim \int_0^\infty \frac{r^{\alpha-1} e^{-\lambda_1 r}}{1\wedge r^{d/2}} e^{-\frac{|x-y|^2}{C_2 r}}\, \d r\\
	& \lesssim \int_0^{|x-y|^2} r^{\alpha-d/2-1} e^{-\frac{|x-y|^2}{C_2r}} \, \d r + \int_{|x-y|^2}^{2M^2} r^{\alpha-d/2-1} \, \d r + \int_{2M^2}^\infty r^{\alpha-1} e^{-\lambda_1 r} \, \d r.
\end{align*}
For the first term, we may use the fact that $\sup_{0<z \le 1} z^{\alpha-d/2-1}e^{-1/z} < \infty$. Then, it follows that
\[
	G_\alpha(x,y) \lesssim |x-y|^{2\alpha-d} + K_{d,\alpha}(|x-y|) + 1 \lesssim K_{d,\alpha}(|x-y|).
\]
To show the lower bound, we use the lower bound in Lemma \ref{lem:P} and deduce as follows:
\begin{align*}
	G_\alpha(x,y) &\gtrsim (1\wedge f_1(x))(1\wedge f_1(y)) \int_0^\infty \frac{r^{\alpha-1} e^{-\lambda_1 r}}{1\wedge r^{d/2}} e^{-\frac{|x-y|^2}{C_1 r}}\, \d r\\
	& \gtrsim f_1(x)f_1(y) \int_{|x-y|^2}^{2M^2} r^{\alpha-d/2-1} \, \d r \gtrsim f_1(x) f_1(y) K_{d,\alpha}(|x-y|).
\end{align*}
This completes the proof.
\end{proof}

%\begin{lemma}[Fractional Green kernel]\label{lem:Green}
%For \(\beta>0\), define \(A^{-\beta/2}\) by spectral calculus. In the strong operator topology on \(L^2(D)\),
%\begin{equation}\label{eq:G-int}
%    A^{-\beta/2}
%    =
%    \frac1{\Gamma(\beta/2)}
%    \int_0^\infty r^{\beta/2-1}\e^{-rA}\,\d r .
%\end{equation}
%Equivalently, for \(\varphi,\psi\in L^2(D)\),
%\begin{align}\label{eq:G-form}
%    \inner{\varphi}{A^{-\beta/2}\psi}_{L^2(D)}
%    &=\sum_{n=1}^\infty \lambda_n^{-\beta/2}
%      \inner{\varphi}{f_n}\inner{\psi}{f_n} \\
%    &=\frac1{\Gamma(\beta/2)}\int_0^\infty
%      r^{\beta/2-1}\inner{\varphi}{P_r\psi}_{L^2(D)}\,\d r .
%      \notag
%\end{align}
%We denote by \(G_\beta\) the distributional kernel of \(A^{-\beta/2}\). Whenever the kernel is represented by a locally integrable function, \eqref{eq:G-form} agrees with the corresponding double integral; away from the diagonal it also agrees with
%\[
%    \frac1{\Gamma(\beta/2)}
%    \int_0^\infty r^{\beta/2-1}P_r(x,y)\,\d r .
%\]
%Thus \(G_{2\alpha}\) is the distributional kernel of \(A^{-\alpha}\), and the formal spectral expansion
%\[
%    W_\alpha(\d r,\d y)
%    =\sum_{n\ge1}\lambda_n^{-\alpha/2}f_n(y)\,\d\beta_n(r)
%\]
%has covariance operator \(A^{-\alpha}\). In the regime of this paper \(\alpha>0\), so no renormalization of the operator-valued integral is needed.
%\end{lemma}

%\begin{proof}
%For each eigenvalue \(\lambda_n>0\),
%\[
%    \lambda_n^{-\beta/2}
%    =\frac1{\Gamma(\beta/2)}
%      \int_0^\infty r^{\beta/2-1}\e^{-\lambda_n r}\,\d r .
%\]
%The spectral theorem applies this scalar identity to \(A\) and yields \eqref{eq:G-int}. Pairing with \(\varphi,\psi\in L^2(D)\) gives \eqref{eq:G-form}; for \(\varphi=\psi\) the interchange follows from the positive spectral measure, and the general case follows by polarization. 
%Testing \eqref{eq:G-form} on \(C_c^\infty(D)\otimes C_c^\infty(D)\) identifies the distributional kernel.
%\end{proof}

%Let \(\cH_\alpha(D)\) be the completion of \(L^2(D)\) under
%\begin{equation}\label{eq:H}
%    \inner{\varphi}{\psi}_{\cH_\alpha(D)}
%    =
%    \sum_{n=1}^\infty
%    \lambda_n^{-\alpha}
%    \inner{\varphi}{f_n}_{L^2(D)}
%    \inner{\psi}{f_n}_{L^2(D)} .
%\end{equation}
%This form is positive definite because every \(\lambda_n^{-\alpha}\) is strictly positive and \(A^{-\alpha}\) is injective. For smooth compactly supported functions (or whenever the integral is absolutely convergent), \eqref{eq:H} equals
%\[
%    \int_D\int_D
%    \varphi(x)G_{2\alpha}(x,y)\psi(y)\,\d x\,\d y;
%\]
%for general elements of the completion the identity is understood in the quadratic-form sense by density.
%The noise \(W_\alpha\) is the centered Gaussian isonormal process over
%\[
%    L^2(\R_+;\cH_\alpha(D)).
%\]
%This is the deterministic Hilbert-space realization of the stochastic integral used below. It agrees with Walsh integration when a corresponding worthy martingale-measure realization is available; see \cite{Walsh,DK,DPZ}.
%The noise \(W_\alpha\) is white in time and spatially colored through the spectral calculus of the Dirichlet Laplacian. This bounded-domain perspective is closely related to the colored-noise framework for parabolic SPDEs with Dirichlet or Neumann boundary conditions studied in \cite{CCL23}. It is also parallel to the intrinsic spectral construction of colored Gaussian noises on compact manifolds in \cite{CO25}.
%Thus, for deterministic \(F,G\in L^2(\R_+;\cH_\alpha(D))\),
%\begin{equation}\label{eq:isometry}
%    \E[W_\alpha(F)W_\alpha(G)]
%    =
%    \int_0^\infty
%    \inner{F(r,\cdot)}{G(r,\cdot)}_{\cH_\alpha(D)}\,\d r .
%\end{equation}
%
%\subsection{The additive Gaussian solution}
%
%For \(t>0\) and \(x\in D\), set
%\begin{equation}\label{eq:p}
%    p_{t,x}(r,y)=P_{t-r}(x,y)\1_{(0,t)}(r).
%\end{equation}
%The mild solution \eqref{eq:mild} is \(u(t,x)=W_\alpha(p_{t,x})\) whenever \(p_{t,x}\in L^2(\R_+;\cH_\alpha(D))\). This is the usual mild formulation for stochastic heat equations; see \cite{Walsh,K09,DK}.
%
%\begin{lemma}[Dalang condition]\label{lem:Dalang}
%Let \(\alpha>0\). For every \(t>0\) and \(x\in D\),
%\(p_{t,x}\in L^2(\R_+;\cH_\alpha(D))\) if and only if
%\(\alpha>d/2-1\).
%\end{lemma}
%
%\begin{proof}
%By \eqref{eq:isometry}, \eqref{eq:G-int}, and the semigroup property,
%\begin{align*}
%    \|p_{t,x}\|_{L^2(\R_+;\cH_\alpha(D))}^2
%    &=
%    \int_0^t
%    \int_D\int_D
%    P_{t-r}(x,z)G_{2\alpha}(z,w)P_{t-r}(x,w)\,\d z\,\d w\,\d r     \\
%    &=
%    \frac1{\Gamma(\alpha)}
%    \int_0^t\int_0^\infty
%    q^{\alpha-1}P_{2(t-r)+q}(x,x)\,\d q\,\d r .
%\end{align*}
%Let \(s=t-r\) and \(\tau=2s+q\). Split the region into \(\tau\le 1\) and \(\tau>1\). For the large-\(\tau\) part, the spectral gap of the bounded Dirichlet domain gives \(P_\tau(x,x)\le C_x e^{-\lambda_1\tau/2}\), making the integral finite for every \(\alpha>0\). For the small-\(\tau\) part, domain domination yields \(P_\tau^D(x,x)\le (4\pi\tau)^{-d/2}\). The integral is therefore controlled by
%\[
%    \int_0^t\int_0^1 q^{\alpha-1}(2s+q)^{-d/2}\,\d q\,\d s,
%    \qquad s=t-r.
%\]
%Split the region near the origin into \(0<q\le s\) and \(s<q<1\). On \(q\le s\),
%\[
%    \int_0^s q^{\alpha-1}(2s+q)^{-d/2}\,\d q
%    \lesssim
%    s^{-d/2}\int_0^s q^{\alpha-1}\,\d q
%    \lesssim
%    s^{\alpha-d/2}.
%\]
%On \(s<q<1\), the integral is bounded by
%\[
%    \int_s^1 q^{\alpha-1-d/2}\,\d q,
%\]
%which is \(O(s^{\alpha-d/2})\) if \(\alpha<d/2\), \(O(\log(1/s))\) if \(\alpha=d/2\), and \(O(1)\) if \(\alpha>d/2\). Therefore the remaining \(s\)-integral is finite exactly when
%\[
%    \alpha-\frac d2>-1,
%    \qquad\text{that is,}\qquad
%    \alpha>\frac d2-1 .
%\]
%The upper bound thus follows from the global Euclidean heat-kernel upper bound. For necessity, fix \(x\in D\) and choose \(r_0<d(x)\). By domain monotonicity, \(P_\tau^D(x,x)\ge P_\tau^{B(x,r_0)}(x,x)\). 
%The Dirichlet heat kernel of a ball satisfies \(P_\tau^{B(x,r_0)}(x,x)\ge c\tau^{-d/2}\) for all sufficiently small \(\tau\). Restricting the double integral to \(\tau\le \tau_0\) shows that it diverges whenever \(\alpha\le d/2-1\). This is the bounded-domain analogue of the usual Dalang-type integrability condition for random-field solutions; see \cite{Walsh,DK,HSWX20}.
%The assumption \(\alpha>0\) is used here only to invoke \eqref{eq:G-int}; within this positive-order covariance class the criterion is the usual Dalang condition.
%\end{proof}
%
%\begin{proof}[Proof of Theorem~\ref{thm:exist}]
%Lemma~\ref{lem:Dalang} gives the existence criterion and the definition \(u(t,x)=W_\alpha(p_{t,x})\). Since \(W_\alpha\) is isonormal, \(u\) is a centered Gaussian field. Let
%\[
%    u_N(t,x)=\sum_{n=1}^N\lambda_n^{-\alpha/2}f_n(x)
%       \int_0^t\e^{-\lambda_n(t-r)}\,\d\beta_n(r).
%\]
%The isometry and Lemma~\ref{lem:Dalang} give
%\[
%    \E|u(t,x)-u_N(t,x)|^2
%    =\sum_{n>N}\lambda_n^{-\alpha}f_n(x)^2
%      \frac{1-\e^{-2\lambda_nt}}{2\lambda_n}
%    \longrightarrow0.
%\]
%Thus the series \eqref{eq:spectral} converges in \(L^2(\Omega)\) at every \((t,x)\). Computing covariances for finite partial sums and passing to the limit by Cauchy--Schwarz gives both identities in \eqref{eq:cov}; in particular, the displayed covariance series converges at every pair of parameter points.
%\end{proof}

\subsection{Gaussian noise and stochastic integrals}

Let $\alpha > 0$.
The Gaussian noise $W_\alpha$ in \eqref{eq:SHE} is defined as follows.
Define the Hilbert space $\mathcal{H}^\alpha(D)$ as the completion of $\mathcal{S}(D)$ with respect to the inner product
\begin{align}\begin{split}\label{ip:alpha}
    &\langle \phi\,, \psi \rangle_{\alpha}
    := \sum_{n=1}^\infty \lambda_n^{-\alpha} \langle \phi\,, f_n \rangle_{L^2(D)} \langle \psi\,, f_n \rangle_{L^2(D)}\\
    &= \big\langle (-\Delta)^{-\alpha/2} \phi\,, (-\Delta)^{-\alpha/2} \psi \big\rangle_{L^2(D)}
    = \int_D \int_D \phi(y) G_{\alpha}(y,z) \psi(z) \, \d y \, \d z,
    % \qquad \forall \varphi, \psi \in L^2(D).
\end{split}\end{align}
where the last two identities are due to \eqref{frac:Lapl} and \eqref{G:id}.
Next, define the Hilbert space
\begin{equation*}\label{eq:Hilbertspace}
    \mathcal{H}_{\alpha}(D) = L^2(\R_+\,; \mathcal{H}^{\alpha}(D)) = \left\{ f: \R_+ \to \mathcal{H}^\alpha(D) : \int_0^\infty \|f(s\,,\cdot)\|_{\alpha}^2\, \d s < \infty \right\}.
\end{equation*}
The noise $W_\alpha =\{ W_{\alpha}(\phi)\,; \phi \in \mathcal{H}_{\alpha}(D) \}$ is defined as a centered isonormal Gaussian process on a complete probability space $(\Omega, \mathcal{F}, \P)$ with covariance
\begin{align}\begin{split}\label{W:corr:WfWg}
    &\E[W_{\alpha}(f)W_{\alpha}(g)] = \int_0^\infty \langle f(s,\cdot), g(s,\cdot) \rangle_{\alpha} \, \d s\\
    &= \frac{1}{\Gamma(\alpha)} \int_0^\infty \d s \int_0^\infty \d r \, r^{\alpha-1} \int_D \int_D \d y\, \d z \, f(s,y) P_r(y,z) g(s,z)
\end{split}\end{align}
for any $f, g \in \mathcal{H}_\alpha(D)$, where the last equality follows from \eqref{G} and \eqref{ip:alpha}.
For any $t>0$ and $x \in D$, define $p_{t,x}$ by
\begin{align}\label{p}
    p_{t,x}(s,y) = P_{t-s}(x,y) \1_{(0,t)}(s) \qquad \forall s \in (0,t), y \in D.
\end{align}
Hence, \eqref{eq:mild} can be interpreted as $u(t,x) = W_\alpha(p_{t,x})$ whenever it is well defined.

\begin{lemma}\label{lem:p:S(D)}
For any $t>s>0$ and $x \in D$, $p_{t,x}(s,\cdot) \in \mathcal{S}(D)$.
\end{lemma}

\begin{proof}
Fix $t>s>0$ and $x \in D$.
By \eqref{P} and orthonormality of $\{f_n\}_{n \ge 1}$, for any $k \ge 0$,
\begin{align*}
	n^k \left|\langle p_{t,x}(s,\cdot), f_n \rangle_{L^2(D)}\right|
	& = n^k \left| \langle P_{t-s}(x,\cdot), f_n\rangle_{L^2(D)}\right|
	= n^k |e^{-\lambda_n(t-s)} f_n(x)|\\
	&\lesssim n^k e^{-\lambda_n(t-s)} \lambda_n^{d/4} \to 0 \quad \text{as $n \to \infty$,}
\end{align*}
where we have used the fact that $\sup_{x \in D}|f_n(x)|\lesssim \lambda_n^{d/4}$ (see \cite[Proposition 5]{F95}) to obtain the last inequality and Weyl's law $\lambda_n \asymp n^{2/d}$ (see \cite[Corollary 3.15]{FLW}) to conclude the limit as $n \to \infty$.
This shows that $p_{t,x}(s,\cdot) \in \mathcal{S}(D)$.
\end{proof}

\section{Existence and variance bounds}\label{s:exist}

In this section, we establish a necessary and sufficient condition for the existence of solution and derive variance bounds.
We start with the following lemma.

\begin{lemma}\label{lem:r-integral}
Fix $\alpha>0$ and $\lambda>0$.
Then, there exists $C>0$ such that
\[
	\int_0^\infty \frac{r^{\alpha-1} e^{-\lambda(2s+r)}}{1 \wedge (2s+r)^{d/2}} \, \d r
	\le \begin{cases}
	C e^{-2\lambda s}(1+ s^{\alpha-d/2}) & \text{if $\alpha<d/2$,}\\
	C e^{-2\lambda s}\log(e + 1/s) & \text{if $\alpha=d/2$,}\\
	C e^{-2\lambda s} & \text{if $\alpha>d/2$,}
	\end{cases}
\]
uniformly for all $s>0$.
\end{lemma}

\begin{proof}
For $0<s<1/2$, we split the integral and estimate as follows:
\[
	\int_0^{2s} \frac{r^{\alpha-1}\d r}{(2s)^{d/2}}  + \int_{2s}^1 r^{\alpha-d/2-1} \, \d r + \int_1^\infty r^{\alpha-1} e^{-\lambda(2s+r)} \, \d r
	\lesssim
	\begin{cases}
	s^{\alpha-d/2} & \text{if $\alpha<d/2$,}\\
	\log(1/s) & \text{if $\alpha=d/2$,}\\
	1 & \text{if $\alpha>d/2$.}
	\end{cases}
\]
For $s\ge 1/2$, we bound the integral by
\[
	\int_0^\infty r^{\alpha-1} e^{-\lambda(2s+r)} \, \d r \propto e^{-2\lambda s}.
\]
Combine the two cases to obtain the asserted bounds.
\end{proof}

\subsection{Proof of Theorem \ref{thm:exist}}

\begin{proof}
We shall prove that $p_{t,x} \in \mathcal{H}_\alpha(D)$ for every $(t,x) \in (0\,,\infty) \times D$ if and only if \eqref{Dalang} holds. That is,
\[
	\int_0^\infty \|p_{t,x}(s,\cdot)\|_\alpha^2 \, \d s < \infty 
	\quad \Leftrightarrow \quad 
	\alpha>d/2-1.
\]
By Lemma \ref{lem:p:S(D)}, for any $t>s>0$ and $x \in D$, $p_{t,x}(s,\cdot) \in \mathcal{S}(D)$, so we may use \eqref{ip:alpha}, \eqref{G}, symmetry of $P_t$, and the semigroup property \eqref{semigroup} to write
\begin{align}
	&\int_0^\infty \|p_{t,x}(s,\cdot)\|_\alpha^2 \, \d s
        = \int_0^t \int_D \int_D P_{t-s}(x,y) G_\alpha(y,z) P_{t-s}(x,z) \, \d y \, \d z \, \d s\notag \\
        & = \frac{1}{\Gamma(\alpha)} \int_0^t \d s \int_0^\infty \d r \, r^{\alpha-1} \int_D \int_D \d y\, \d z\, P_s(x,y) P_r(y,z) P_s(z,x)\notag \\
        & = \frac{1}{\Gamma(\alpha)} \int_0^t \d s \int_0^\infty \d r \, r^{\alpha-1} P_{2s+r}(x,x).
        \label{int:p:norm-sq}
\end{align}
Lemmas \ref{lem:P} and \ref{lem:r-integral} imply that
\begin{align}
	\int_{1/4}^t \d s \int_0^\infty \d r \, r^{\alpha-1} P_{2s+r}(x,x) \lesssim \int_{1/4}^t \d s \int_0^\infty \d r \, \frac{r^{\alpha-1} e^{-\lambda_1(2s+r)}}{1 \wedge (2s+r)^{d/2}} =O(1)
\end{align}
uniformly for all $t \ge 1/4$.
Hence, it suffices to prove that
\begin{align}\label{I:finite}
	I=I(t,x) < \infty \quad \forall (t,x)\in (0,1/4)\times D
	\quad \Leftrightarrow \quad
	\alpha>d/2-1,
\end{align}
where
\[
	I=I(t,x):=\int_0^t \d s \int_0^\infty \d r \, r^{\alpha-1} P_{2s+r}(x,x).
\]
%First, suppose that $\alpha>d/2-1$.
%We write $I = I_1 + I_2 + I_3$, where
%\begin{align*}
%        &I_1 = \int_0^t \d s \int_0^{2s} \d r \, r^{\alpha-1} P_{2s+r}(x,x),\\
%        &I_2 = \int_0^t \d s \int_{2s}^1 \d r \, r^{\alpha-1} P_{2s+r}(x,x),\\
%        &I_3 = \int_0^t \d s \int_1^\infty \d r \, r^{\alpha-1} P_{2s+r}(x,x).
%\end{align*}
%Then, by Lemma \ref{lem:P},
%\begin{align*}
%        I_1 \lesssim \int_0^t \d s \int_0^{2s} \d r \, \frac{r^{\alpha-1}\e^{-\lambda_1(2s+r)}}{1 \wedge (2s+r)^{d/2}} \lesssim \int_0^t \d s \, s^{-d/2} \int_0^{2s} r^{\alpha-1}  \lesssim t^{\alpha-d/2+1},
%\end{align*}
%uniformly for all $t \in (0,1/4)$ and $x \in D$, where the last inequality holds due to the assumption that $\alpha>d/2-1$.
%Similarly, we have
%\begin{align*}
%        I_2 \lesssim \int_0^t \d s \int_{2s}^1 \d r \, r^{\alpha-d/2-1} \lesssim 
%        \begin{cases}
%        \int_0^t 1\, \d s = t & \text{if $\alpha-d/2>0$,}\\
%        \int_0^t s^{\alpha-d/2}\, \d s \lesssim t^{\alpha-d/2+1} & \text{if $\alpha-d/2<0$,}
%	\end{cases}
%\end{align*}
%uniformly for all $t \in (0,1/4)$, $x \in D$, and
%\begin{align*}
%        I_3 \lesssim \int_0^t \d s \int_1^\infty \d r \, r^{\alpha-1} \e^{-\lambda_1(2s+r)} \lesssim t,
%\end{align*}
%uniformly for all $t \in (0,1/4)$ and $x \in D$.
%Combining the estimates yields $I<\infty$.
%Moreover, combining the estimates yields $\Var(W_\alpha(p_{t,x})) \propto I(t\,,x) \lesssim t^{\alpha-d/2+1}+t$ uniformly for all $t \in (0\,,1/4)$ and $x \in D$.
First, suppose that $\alpha>d/2-1$.
By Lemmas \ref{lem:P} and \ref{lem:r-integral}, 
\[
	I \lesssim \int_0^t \d s \int_0^\infty \d r \, \frac{r^{\alpha-1} e^{-\lambda_1(2s+r)}}{1 \wedge (2s+r)^{d/2}} \lesssim \begin{cases}
	t^{\alpha-d/2 + 1} & \text{if $d/2-1<\alpha<d/2$,}\\
	t \log_+(1/t) & \text{if $\alpha=d/2$,}\\
	t & \text{if $\alpha>d/2$,}
	\end{cases}
\]
uniformly for all $t \in (0,1/4)$. 
This proves that $I<\infty$.

Conversely, suppose that $I<\infty$. We may use the lower bound in Lemma \ref{lem:P} to deduce that
\begin{align*}
	I &\gtrsim (1\wedge f_1(x))^2 \int_0^t \d s \int_0^{2s} \d r \, \frac{r^{\alpha-1}\e^{-\lambda_1(2s+r)}}{(2s+r)^{d/2}}\\
        &\gtrsim |f_1(x)|^2\, \e^{-\lambda_1} \int_0^t \d s\, s^{-d/2} \int_0^{2s} \d r \, r^{\alpha-1}
        \gtrsim |f_1(x)|^2\, \e^{-\lambda_1}\int_0^t s^{\alpha-d/2} \, \d s,
\end{align*}
uniformly for all $t \in (0,1/4)$ and $x \in D$.
Since $I<\infty$, the last integral $\int_0^t s^{\alpha-d/2} \, \d s$ is finite, which implies that $\alpha>d/2-1$.
This proves \eqref{I:finite} and completes the proof of Theorem \ref{thm:exist}.
\end{proof}

\subsection{Variance bounds}

The proof of Theorem \ref{thm:exist} implies the following estimate:

\begin{lemma}\label{lem:var:u:1}
    Suppose $D \subset \R^d$ is a bounded Lipschitz domain. If $\alpha>d/2-1$, then 
\[
	\Var(u(t,x)) \lesssim 1 \wedge (\rho_0(t))^2
\]
uniformly for all $t>0$ and $x \in D$, where
\begin{align}\label{rho_0}
	\rho_0(t)
	:= \begin{cases}
	t^{(2\alpha-d+2)/4} & \text{if $d/2-1<\alpha<d/2$,}\\
	\sqrt{t \log_+(1/t)}& \text{if $\alpha = d/2$,}\\
	\sqrt{t} & \text{if $\alpha > d/2$.}
	\end{cases}
\end{align}
\end{lemma}

%\begin{lemma}\label{lem:var:u:lb}
%Suppose $D \subset \R^d$ is a bounded Lipschitz domain. If $\alpha>d/2-1$, then for each $T>0$, there exists $C>0$ depending on $T$ such that
%\[
%	\Var(u(t,x)) \ge C |f_1(x)|^2\, t^{\alpha-d/2+1}
%\]
%uniformly for all $t \in [0,T]$ and $x \in D$.
%\end{lemma}

Next, we establish the dependence of $\Var(u(t,x))$ on the distance-to-boundary $d(x)$ defined by \eqref{d(x)}.

\begin{lemma}\label{lem:var:u:2}
Suppose $D \subset \R^d$ is a bounded $C^{1,\gamma}$ domain for some $\gamma>0$. If $\alpha>d/2-1$, then 
\[
	\Var(u(t,x)) \lesssim (d_0(x))^2
\]
uniformly for all $t>0$ and $x \in D$, where
\begin{align}\label{d_0}
d_0(x):=\begin{cases}
(d(x))^{(2\alpha-d+2)/2}  & \text{if $d/2-1<\alpha<d/2$,}\\
d(x)\sqrt{\log_+(1/d(x))} & \text{if $\alpha = d/2$,}\\
d(x)& \text{if $\alpha > d/2$.}
\end{cases}
\end{align}
\end{lemma}

\begin{proof}
Note that the case that $d(x) \ge 1$ is trivial because Lemma \ref{lem:var:u:1} implies that
\[
	\Var(u(t,x)) \lesssim 1 \lesssim (d_0(x))^2,
\]
uniformly for all $t>0$ and $x \in D$ with $d(x) \ge 1$.
Next, we consider the case that $d(x)<1$.
By \eqref{W:corr:WfWg}, \eqref{int:p:norm-sq}, the change of variable $r \to 2s+r$ and Fubini's theorem,
\begin{align*}
	\Var(u(t,x)) 
	&\propto \int_0^t \d s \int_0^\infty \d r \, r^{\alpha-1} P_{2s+r}(x,x)\\
	&= \int_0^t \d s \int_{2s}^\infty \d r \, (r-2s)^{\alpha-1} P_r(x,x)\\
	&= \int_0^\infty \left[\int_0^{t \wedge (r/2)} (r-2s)^{\alpha-1}  \d s \right] P_r(x,x) \, \d r
	\lesssim \int_0^\infty r^{\alpha} P_r(x,x) \, \d r,
\end{align*}
uniformly for all $t>0$ and $x \in D$.
Then, we may use Lemma \ref{lem:P} and the inequality $e^{-\lambda_1 r}/(1 \wedge r^{d/2}) \lesssim r^{-d/2}e^{-\frac12 \lambda_1 r}$ for all $r>0$ to deduce the following:
\begin{align*}
	&\Var(u(t,x))
	\lesssim \int_0^\infty \left( 1 \wedge \frac{d(x)}{1\wedge \sqrt{r}} \right)^2 r^{\alpha-d/2} e^{-\frac12\lambda_1 r} \, \d r\\
	& \lesssim \int_0^{d^2(x)} r^{\alpha-d/2} \, \d r + d^2(x) \int_{d^2(x)}^1 r^{\alpha-d/2-1} \, \d r + d^2(x)\int_1^\infty r^{\alpha-d/2} e^{-\frac12\lambda_1 r} \, \d r\\
	&\lesssim (d_0(x))^2
\end{align*}
uniformly for all $t>0$ and $x \in D$ with $d(x)<1$.
This completes the proof.
\end{proof}

Combining Lemmas \ref{lem:var:u:1} and \ref{lem:var:u:2} yields the following:

\begin{proposition}\label{prop:var:u}
Suppose $D \subset \R^d$ is a bounded $C^{1,\gamma}$ domain for some $\gamma>0$.
If $\alpha>d/2-1$, then
\[
	\Var(u(t,x)) \lesssim (d_0(x) \wedge \rho_0(t))^2
\]
uniformly for all $t>0$ and $x \in D$, where $\rho_0$ and $d_0$ are given by \eqref{rho_0} and \eqref{d_0}, respectively.
\end{proposition}

%\begin{proof}
%    By \eqref{ip:alpha}, \eqref{G:id} and the symmetry and semigroup property of the heat kernel $P$,
%    \begin{align*}
%        &\int_0^\infty \|p_{t,x}(s\,,\cdot)\|_\alpha^2 \, \d s
%        = \int_0^t \int_D \int_D P_{t-s}(x\,,y) G_\alpha(y\,,z) P_{t-s}(x\,,z) \, \d y \, \d z \, \d s\\
%        & = \frac{1}{\Gamma(\alpha)} \int_0^t \d s \int_0^\infty \d r \, r^{\alpha-1} \int_D \int_D \d y\, \d z\, P_s(x\,,y) P_r(y\,,z) P_s(z\,,x)\\
%        & = \frac{1}{\Gamma(\alpha)} \int_0^t \d s \int_0^\infty \d r \, r^{\alpha} P_{2s+r}(x\,,x). 
%    \end{align*}
%    By Lemma \ref{lem:P}, for all $t\ge 1/4$ and $x \in D$, we have
%    \begin{align}
%        \int_{1/4}^t \d s \int_0^\infty \d r \, r^{\alpha-1} P_{2s+r}(x\,,x) \lesssim \int_{1/4}^t \d s \int_0^\infty \d r \, r^{\alpha-1}\e^{-\lambda_1(2s+r)} =O(1).
%    \end{align}
%    Hence, $p_{t,x} \in \mathcal{H}_\alpha(D)$ for all $(t\,,x) \in (0\,,\infty) \times D$ if and only if 
%    \begin{align}\label{I:finite}
%        I=I(t\,,x):=\int_0^t \d s \int_0^\infty \d r \, r^{\alpha-1} P_{2s+r}(x\,,x) < \infty \quad \forall t\in (0\,,1/4), x \in D.
%    \end{align}
%    In particular, we may write $I = I_1 + I_2 + I_3$, where
%    \begin{align*}
%        &I_1 = \int_0^t \d s \int_0^{2s} \d r \, r^{\alpha-1} P_{2s+r}(x\,,x),\\
%        &I_2 = \int_0^t \d s \int_{2s}^1 \d r \, r^{\alpha-1} P_{2s+r}(x\,,x),\\
%        &I_3 = \int_0^t \d s \int_1^\infty \d r \, r^{\alpha-1} P_{2s+r}(x\,,x).
%    \end{align*}
%
%    Suppose $\alpha>d/2-1$.
%    Then, by Lemma \ref{lem:P},
%    \begin{align*}
%        I_1 \lesssim \int_0^t \d s \int_0^{2s} \d r \, \frac{r^{\alpha-1}\e^{-\lambda_1(2s+r)}}{1 \wedge (2s+r)^{d/2}} \lesssim \int_0^t \d s \, s^{-d/2} \int_0^{2s} r^{\alpha-1}  \lesssim t^{\alpha-d/2+1},
%    \end{align*}
%    uniformly for all $t \in (0\,,1/4)$ and $x \in D$, where the last inequality holds due to the assumption that $\alpha>d/2-1$.
%    Similarly, we have
%    \begin{align*}
%        I_2 \lesssim \int_0^t \d s \int_{2s}^1 \d r \, r^{\alpha-d/2-1} \lesssim 
%        \begin{cases}
%        \int_0^t 1\, \d s = t & \text{if $\alpha-d/2>0$,}\\
%        \int_0^t s^{\alpha-d/2}\, \d s \lesssim t^{\alpha-d/2+1} & \text{if $\alpha-d/2<0$,}
%        \end{cases}
%    \end{align*}
%    uniformly for all $t \in (0\,,1/4)$, $x \in D$, and
%    \begin{align*}
%        I_3 \lesssim \int_0^t \d s \int_1^\infty \d r \, r^{\alpha-1} \e^{-\lambda_1(2s+r)} \lesssim t,
%    \end{align*}
%    uniformly for all $t \in (0\,,1/4)$ and $x \in D$.
%    Moreover, combining the estimates yields $\Var(W_\alpha(p_{t,x})) \propto I(t\,,x) \lesssim t^{\alpha-d/2+1}+t$ uniformly for all $t \in (0\,,1/4)$ and $x \in D$.
%
%    Conversely, suppose \eqref{I:finite} holds. By Lemma \ref{lem:P}, we have
%    \begin{align*}
%        I_1 &\gtrsim (1\wedge \phi_1(x))^2 \int_0^t \d s \int_0^{2s} \d r \, \frac{r^{\alpha-1}\e^{-\lambda_1(2s+r)}}{(2s+r)^{d/2}}\\
%        &\gtrsim (1\wedge \phi_1(x))^2 \e^{-\lambda_1} \int_0^t \d s\, s^{-d/2} \int_0^{2s} \d r \, r^{\alpha-1}
%        \gtrsim (1\wedge \phi_1(x))^2 \e^{-\lambda_1}\int_0^t s^{\alpha-d/2} \, \d s,
%    \end{align*}
%    uniformly for all $t \in (0\,,1/4)$ and $x \in D$.
%    Since $I<\infty$, it follows that the last integral is finite, which implies that $\alpha>d/2-1$.
%\end{proof}

\subsection{A series representation}

\begin{proposition}\label{prop:spec:rep}
Suppose $D \subset \R^d$ is a bounded Lipschitz domain. 
Then, under condition \eqref{Dalang}, the solution has the representation
\begin{equation}\label{eq:spectral}
    u(t,x)=
    \sum_{n=1}^\infty
    \lambda_n^{-\alpha/2}
    \left(\int_0^t \e^{-\lambda_n(t-r)}\,\d B_n(r)\right)f_n(x),
\end{equation}
where \(\{B_n(t)\}_{n\ge1}\) are independent standard Brownian motions.
% and the series converges in $L^2(\Omega)$.
\end{proposition}

\begin{proof}
Note that both sides of \eqref{eq:spectral} are centered Gaussian processes, so it suffices to show that they have the same covariance function.
The covariance of $u(t,x)$ can be computed as follows, using \eqref{eq:mild}, the first identities in both \eqref{ip:alpha} and \eqref{W:corr:WfWg}, and \eqref{P}:
\begin{align*}
	\E[u(t,x)u(s,y)]
	&=\int_0^{t \wedge s} \langle P_{t-r}(x, \cdot) , P_{s-r}(y,\cdot) \rangle_\alpha\, \d r\\
%    &=
%    \int_0^{t\wedge s}
%    \int_D\int_D
%    P_{t-r}(x,z)P_{s-r}(y,w)G_{2\alpha}(z,w)\,\d z\,\d w\,\d r  \\
	&=\sum_{n=1}^\infty\lambda_n^{-\alpha}f_n(x)f_n(y)\int_0^{t\wedge s} \e^{-\lambda_n(t-r)}\e^{-\lambda_n(s-r)}\,\d r.
\end{align*}
The last expression coincides with the covariance of the right-hand side of \eqref{eq:spectral} by independence of $B_n$ and Wiener isometry.
Hence, the Gaussian processes on both sides of \eqref{eq:spectral} have the same law.
\end{proof}

\section{H\"older regularity}\label{sec:regularity}

In this section, we study temporal and spatial H\"older regularity of the solution. 
We first establish some lemmas.

\begin{lemma}\label{lem:spectral:1}
For any $a \in \R$, there exists $C>0$ such that
\begin{align*}
	\sum_{n \ge 1: \lambda_n \le R} \lambda_n^{-a} |f_n(x)|^2 \le 
	\begin{cases}
	C(1+R^{d/2-a}) & \text{if $a<d/2$,}\\
	C\log_+ R & \text{if $a=d/2$,}\\
	C & \text{if $a>d/2$,}
	\end{cases}
\end{align*}
uniformly for all $x \in D$ and $R > 0$.
\end{lemma}

\begin{proof}
First consider $a=0$. By Lemma \ref{lem:P},
\[
	\sum_{n: \lambda_n \le R} |f_n(x)|^2 \le e \sum_{n: \lambda_n \le R} e^{-\lambda_n/R}|f_n(x)|^2
	\le e \, P_{1/R}(x,x) \lesssim R^{d/2}.
\]
Next, we let $k_0 \in \mathbb{Z}$ be such that $2^{k_0}< \lambda_1\le 2^{k_0+1}$.
Then, for any $a \in \R$, a dyadic decomposition, together with the preceding, implies that
\begin{align*}
	\sum_{n \ge 1: \lambda_n \le R} \lambda_n^{-a} |f_n(x)|^2 
	&\le \sum_{k\ge k_0: 2^k \le R} 2^{-a k}\sum_{n\ge 1: 2^{k}< \lambda_n \le 2^{k+1}} |f_n(x)|^2\\
	&\lesssim \sum_{k \ge k_0: 2^k \le R} 2^{(d/2-a)k}
	\lesssim 
	\begin{cases}
	1+R^{d/2-a} & \text{if $a<d/2$,}\\
	\log_+ R & \text{if $a=d/2$,}\\
	1 & \text{if $a>d/2$.}
	\end{cases}
\end{align*}
This completes the proof.
\end{proof}

\begin{lemma}\label{lem:spectral:2}
For any $a > d/2$, there exists $C>0$ such that
\begin{align*}
	\sum_{n \ge 1: \lambda_n \ge R} \lambda_n^{-a} |f_n(x)|^2 \le C R^{d/2-a}
\end{align*}
uniformly for all $x \in D$ and $R>0$.
\end{lemma}

\begin{proof}
Let $k_1 \in \mathbb{Z}$ be such that $2^{k_1} \le R < 2^{k_1+1}$. Then by Lemma \ref{lem:spectral:1},
\begin{align*}
	\sum_{n \ge 1: \lambda_n \ge R} \lambda_n^{-a} |f_n(x)|^2 
	&\le \sum_{k= k_1}^\infty 2^{-a k}\sum_{n\ge 1: 2^{k}< \lambda_n \le 2^{k+1}} |f_n(x)|^2\\
	&\lesssim \sum_{k=k_1}^\infty 2^{(d/2-a)k}
	\lesssim 2^{(d/2-a)k_1}
	\lesssim R^{d/2-a}.
\end{align*}
This completes the proof.
\end{proof}

\begin{lemma}\label{lem:spectral:3}
Suppose \(D\subset\R^d\) is a bounded \(C^2\) domain.
Then, for any $a<d/2+1$, there exists $C>0$ such that
\[
    \sum_{n \ge 1: \lambda_n\le R}\lambda_n^{-a}|\nabla f_n(x)|^2
    \le CR^{d/2+1-a}
\]
uniformly for all $x \in D$ and $R > 0$.
\end{lemma}

\begin{proof}
By \eqref{P} and Lemma \ref{lem:grad:P}, for any $R > 0$,
\[
\begin{aligned}
    \sum_{n \ge 1}\e^{-2\lambda_n/R}|\nabla f_n(x)|^2
    &=
    \int_D|\nabla_x P_{R^{-1}}(x,y)|^2\,\d y        \\
    &\lesssim
    R^{d+1}
    \int_D
    \exp\left(-\frac{2|x-y|^2}{C/R}\right)\,\d y    
    \lesssim R^{d/2+1}.
\end{aligned}
\]
Therefore
\[
    \sum_{n \ge 1: \lambda_n\le R}|\nabla f_n(x)|^2
    \le
    \e^2\sum_{n \ge 1} \e^{-2\lambda_n/R}|\nabla f_n(x)|^2
    \lesssim R^{d/2+1}.
\]
Let $k_0 \in \mathbb{Z}$ be such that $2^{k_0} < \lambda_1 \le 2^{k_0+1}$.
Then, for any $a<d/2+1$, applying the same dyadic decomposition as in the proof of Lemma \ref{lem:spectral:1} gives
\[
\begin{aligned}
    \sum_{n \ge 1: \lambda_n\le R}
    \lambda_n^{-a}|\nabla f_n(x)|^2
    &\le \sum_{k \ge k_0: 2^k \le R} 2^{-ak} \sum_{n \ge 1: 2^k < \lambda_n \le 2^{k+1}} |\nabla f_n(x)|^2\\
    & \lesssim \sum_{k \ge k_0: 2^k \le R} 2^{(d/2+1-a)k}  \lesssim R^{d/2+1-a}.
\end{aligned}
\]
This completes the proof of Lemma \ref{lem:spectral:3}.
\end{proof}


\subsection{Temporal regularity}

%We first establish temporal and spatial increment estimates for the solution.

\begin{proposition}\label{prop:u-u:t}
Suppose $D \subset \R^d$ is a bounded Lipschitz domain and $\alpha>d/2-1$.
Then, for any $T>0$, there exists $C>0$ such that
\[
	\Var(u(t',x)-u(t,x)) \le C \rho_0^2(|t'-t|)
%	\begin{cases}
%	C|t'-t|^{\alpha-d/2+1} & \text{if $d/2-1<\alpha<d/2$,}\\
%	C|t'-t| \log_+(1/|t'-t|) & \text{if $\alpha=d/2$,}\\
%	C|t'-t| & \text{if $\alpha>d/2$,}
%	\end{cases}
\]
uniformly for all $t,t' \in [0,T]$ and $x \in D$, where $\rho_0$ is given by \eqref{rho_0}, i.e.,
\[
	\rho_0^2(r) = \begin{cases}
	r^{\alpha-d/2+1} & \text{if $d/2-1<\alpha<d/2$,}\\
	r \log_+(1/r) & \text{if $\alpha=d/2$,}\\
	r & \text{if $\alpha>d/2$.}
	\end{cases}
\]
\end{proposition}

\begin{proof}
For $0\le t < t' \le T$ and $x \in D$, we may use \eqref{W:corr:WfWg} and the semigroup property \eqref{semigroup} to write $\Var(u(t',x)-u(t,x))\lesssim I_1+I_2$, where
\begin{align*}
	& I_1:=\int_t^{t'}\d s\int_0^\infty\d r\,
    r^{\alpha-1}P_{2(t'-s)+r}(x,x) = \int_0^{t'-t} \d s \int_0^\infty \d r \, r^{\alpha-1} P_{2s+r}(x,x),\\
	&I_2:=\int_0^t \d s \int_0^\infty \d r\, r^{\alpha-1} \big[P_{2s+2(t'-t)+r}(x,x)-2P_{2s+(t'-t)+r}(x,x) 
        +P_{2s+r}(x,x) \big].
\end{align*}
By Lemma \ref{lem:var:u:1},
\begin{align}\label{var:u-u:t:I1}
	I_1 = \Var(u(t'-t, x)) \lesssim 
%	(\rho_0(t'-t))^2
	\begin{cases}
	|t'-t|^{\alpha-d/2+1} & \text{if $d/2-1<\alpha<d/2$,}\\
	|t'-t| \log_+(1/|t'-t|) & \text{if $\alpha=d/2$,}\\
	|t'-t| & \text{if $\alpha>d/2$.}
	\end{cases}
\end{align}
Next, we may use \eqref{P} and Fubini's theorem to write
\begin{align*}
	I_2 &= \int_0^t \d s \int_0^\infty \d r \, r^{\alpha-1} \sum_{n=1}^\infty e^{-\lambda_n(2s+r)} \left( 1 - e^{-\lambda_n(t'-t)} \right)^2 |f_n(x)|^2\\
	& = \sum_{n=1}^\infty \int_0^\infty \d r \, r^{\alpha-1} e^{-\lambda_n r} \left( \frac{1-e^{-2\lambda_n t}}{2\lambda_n} \right) \left( 1 - e^{-\lambda_n(t'-t)} \right)^2 |f_n(x)|^2.
%	& \lesssim \sum_{n=1}^\infty \lambda_n^{-\alpha-1} 
\end{align*}
Then, the identity $\int_0^\infty r^{\alpha-1} e^{-\lambda r} \, \d r = \Gamma(\alpha) \lambda^{-\alpha}$ for $\alpha,\lambda>0$ and the inequality $1-e^{-z} \le 1 \wedge z$ for $z>0$ imply that
\begin{align*}
	I_2 &\lesssim \sum_{n=1}^\infty \lambda_n^{-\alpha-1} \left( 1 \wedge \lambda_n^2|t'-t|^2 \right) |f_n(x)|^2\\
	&\lesssim |t'-t|^2 \sum_{n \ge 1: \lambda_n \le 1/|t'-t|} \lambda_n^{-\alpha+1} |f_n(x)|^2 + \sum_{n \ge 1: \lambda_n > 1/|t'-t|} \lambda_n^{-\alpha-1} |f_n(x)|^2\\
	& \lesssim \begin{cases}
	|t'-t|^2 |t'-t|^{-d/2+\alpha-1} + |t'-t|^{-d/2+\alpha+1} & \text{if $d/2-1<\alpha<d/2+1$,}\\
	|t'-t|^2 \log_+(1/|t'-t|) + |t'-t|^{-d/2+\alpha+1} & \text{if $\alpha=d/2+1$,}\\
	|t'-t|^2  + |t'-t|^{-d/2+\alpha+1} & \text{if $\alpha>d/2+1$.}
	\end{cases}
\end{align*}
It follows that
\begin{align}\label{var:u-u:t:I2}
	I_2 \lesssim 
	\begin{cases}
	|t'-t|^{\alpha-d/2+1} & \text{if $d/2-1<\alpha<d/2+1$,}\\
	|t'-t|^2 \log_+(1/|t'-t|) & \text{if $\alpha=d/2+1$,}\\
	|t'-t|^2 & \text{if $\alpha>d/2+1$.}
	\end{cases}
\end{align}
Combining \eqref{var:u-u:t:I1} and \eqref{var:u-u:t:I2} yields the desired estimate.
\end{proof}
%
%\begin{lemma}\label{lem:time}
%Assume \eqref{eq:alpha}. For every \(T>0\),
%\[
%    \E|u(t,x)-u(s,x)|^2\le C|t-s|^{2\theta},
%    \qquad s,t\in[0,T],\ x\in D .
%\]
%\end{lemma}
%
%\begin{proof}
%If \(s=0<t\), the term \(I_2\) below is absent and the estimate of \(I_1\) gives \(\E|u(t,x)|^2\le Ct^{2\theta}\). By symmetry it therefore remains to consider \(0<s<t\le T\). Put \(\delta=t-s\). By the isometry, symmetry, and the semigroup property,
%\[
%    \Var(u(t,x)-u(s,x))\lesssim I_1+I_2,
%\]
%where
%\[
%    I_1
%    :=
%    \int_s^t\int_0^\infty
%    q^{\alpha-1}P_{2(t-r)+q}(x,x)\,\d q\,\d r
%\]
%and
%\[
%\begin{aligned}
%    I_2
%    :=
%    \int_0^s\int_0^\infty q^{\alpha-1}
%    &\bigl[
%        P_{2(s-r)+2\delta+q}(x,x)
%        -2P_{2(s-r)+\delta+q}(x,x)  \\
%    &\hspace{3.2cm}
%        +P_{2(s-r)+q}(x,x)
%    \bigr]\,\d q\,\d r .
%\end{aligned}
%\]
%For \(I_1\), using the diagonal heat-kernel bound and changing variables
%\(\eta=t-r\), \(q=\eta w\), we obtain
%\[
%\begin{aligned}
%    I_1
%    &\lesssim
%    \int_0^\delta\int_0^\infty
%    q^{\alpha-1}(2\eta+q)^{-d/2}\,\d q\,\d\eta       \\
%    &=
%    \left(\int_0^\infty
%    \frac{w^{\alpha-1}}{(2+w)^{d/2}}\,\d w\right)
%    \int_0^\delta \eta^{\alpha-d/2}\,\d\eta
%    \lesssim
%    \delta^{\alpha+1-d/2}.
%\end{aligned}
%\]
%The \(w\)-integral is finite at the origin because \(\alpha>0\), and finite at infinity because \(\alpha<d/2\); the outer \(\eta\)-integral is finite because \(\alpha>d/2-1\). 
%For \(I_2\), let \(h(\tau)=P_\tau(x,x)\). The spectral expansion gives
%\[
%    h(a+2\delta)-2h(a+\delta)+h(a)
%    =
%    \sum_{n\ge1}\e^{-\lambda_n a}
%    (1-\e^{-\lambda_n\delta})^2f_n(x)^2
%    \le h(a)-h(a+\delta).
%\]
%Moreover, the signed derivative identity is
%\[
%    -h'(\tau)
%    =\sum_{n\ge1}\lambda_n\e^{-\lambda_n\tau}f_n(x)^2
%    =
%    \int_D |\nabla_yP_{\tau/2}(x,y)|^2\,\d y
%    \lesssim \tau^{-d/2-1},
%\]
%while \(h(\tau)\lesssim \tau^{-d/2}\). The identity follows first for \(\tau\ge\tau_0>0\) by spectral convergence and the Dirichlet-form identity, and then for every \(\tau>0\). 
%All constants here are uniform in \(x\in D\); for large \(\tau\), the exponential factors in Lemma~\ref{lem:heat} are absorbed into the displayed polynomial bounds. Hence
%\[
%    0\le h(\tau)-h(\tau+\delta)
%    \lesssim
%    \tau^{-d/2}\left(1\wedge\frac{\delta}{\tau}\right).
%\]
%With \(\eta=s-r\) and \(\tau=2\eta+q\),
%\[
%\begin{aligned}
%    I_2
%    &\lesssim
%    \int_0^s\int_0^\infty
%    q^{\alpha-1}\bigl(h(2(s-r)+q)-h(2(s-r)+q+\delta)\bigr)\,\d q\,\d r\\
%    &=
%    \int_0^s\int_{2\eta}^\infty
%    (\tau-2\eta)^{\alpha-1}
%    \bigl(h(\tau)-h(\tau+\delta)\bigr)\,\d\tau\,\d\eta\\
%    &=
%    \int_0^\infty
%    \left[
%        \int_0^{s\wedge \tau/2}(\tau-2\eta)^{\alpha-1}\,\d\eta
%    \right]
%    \bigl(h(\tau)-h(\tau+\delta)\bigr)\,\d\tau\\
%    &\lesssim
%    \int_0^\infty
%    \tau^\alpha\bigl(h(\tau)-h(\tau+\delta)\bigr)\,\d\tau    \\
%    &\lesssim
%    \int_0^\delta \tau^{\alpha-d/2}\,\d\tau
%    +
%    \delta\int_\delta^\infty \tau^{\alpha-d/2-1}\,\d\tau
%    \lesssim
%    \delta^{\alpha+1-d/2}.
%\end{aligned}
%\]
%The bracket in the change of variables is bounded by \(\int_0^{s\wedge\tau/2}(\tau-2\eta)^{\alpha-1}\,\d\eta \le \tau^\alpha/(2\alpha)\), again using \(\alpha>0\).
%Since \(\alpha+1-d/2=2\theta\), the proof is complete.
%\end{proof}

\subsection{Spatial regularity}

In order to establish spatial regularity, we will 
%assume more regularity for the domain and need to use 
need the following lemma, whose proof will be given after the proof of Proposition \ref{prop:u-u:x}.

\begin{lemma}\label{lem:curve}
If $D\subset \R^d$ is a bounded $C^k$ domain for some $k \in \N_+$, then there exist $\varepsilon_0 \in (0,1)$ and $C_0>0$ such that the following property holds: For every pair $x,y \in D$ with $|x-y|\le \varepsilon_0$, there is a $C^k$ curve $\gamma:[0,1] \to D$ such that $\gamma(0)=x$, $\gamma(1)=y$, and
\[
	|\gamma'(\tau)| \le C_0|x-y| \qquad \forall \tau \in [0,1].
\]
\end{lemma}

Assuming Lemma \ref{lem:curve}, we can prove spatial regularity of the solution.

\begin{proposition}\label{prop:u-u:x}
Suppose $D \subset \R^d$ is a bounded $C^2$ domain and $\alpha>d/2-1$.
Then, for any $T>0$, there exists $C>0$ such that
\begin{align}\label{var:u-u:x}
	\Var(u(t,x')-u(t,x)) \le C \rho_1^2(|x'-x|),
\end{align}
uniformly for all $t \in [0,T]$ and $x,x'\in D$, where $\rho_1$ is given by
\[
	\rho_1^2(r) :=\begin{cases}
	r^{2\alpha-d+2} & \text{if $d/2-1<\alpha<d/2$,}\\
	r^2 \log_+(1/r) & \text{if $\alpha=d/2$,}\\
	r^2 & \text{if $\alpha>d/2$,}
	\end{cases}
\]
\end{proposition}

\begin{proof}
Let $\varepsilon_0 \in (0,1)$ and $C_0>0$ be the constants given by Lemma \ref{lem:curve}.
According to Lemma \ref{lem:var:u:1}, $\Var(u(t,x))$ is uniformly bounded, so it suffices to prove \eqref{var:u-u:x} for $|x'-x| \le \varepsilon_0$.
For any $t \in [0,T]$ and $x,x' \in D$ with $|x'-x|\le \varepsilon_0$, we may use \eqref{W:corr:WfWg}, the semigroup property \eqref{semigroup} and \eqref{P} to see that
\begin{align*}
	&\Var(u(t,x')-u(t,x))\\
	&\propto \int_0^t \d s \int_0^\infty \d r \, r^{\alpha-1} \left[ P_{2s+r}(x',x') - 2P_{2s+r}(x',x)+ P_{2s+r}(x,x) \right]\\
	&=\int_0^t \d s \int_0^\infty \d r \, r^{\alpha-1} \sum_{n=1}^\infty e^{-\lambda_n (2s+r)} |f_n(x')-f_n(x)|^2\\
	&=\int_0^t \d s \int_0^\infty \d r \, r^{\alpha-1} \int_D \d y\, |P_{s+r/2}(x',y)-P_{s+r/2}(x,y)|^2.
\end{align*}
On the one hand, we may use $(a+b)^2 \le 2(a^2+b^2)$ for $a,b \ge 0$ and Lemma \ref{lem:P} to deduce that
\begin{align*}
	&\int_D |P_{s+r/2}(x',y)-P_{s+r/2}(x,y)|^2 \, \d y\\
	& \lesssim \int_D \frac{e^{-2\lambda_1 (s+r/2)}}{1 \wedge (s+r/2)^d} \, e^{- \frac{2|x-y|^2}{C_2(s+r/2)}} \, \d y + \int_D \frac{e^{-2\lambda_1 (s+r/2)}}{1\wedge (s+r/2)^d} e^{- \frac{2|x'-y|^2}{C_2(s+r/2)}}  \, \d y\\
	& \lesssim \int_{\R^d} \frac{e^{-\lambda_1 (s+r/2)}}{(s+r/2)^d} \, e^{- \frac{2|x-y|^2}{C_2(s+r/2)}} \, \d y + \int_{\R^d} \frac{e^{-\lambda_1 (s+r/2)}}{ (s+r/2)^d} e^{- \frac{2|x'-y|^2}{C_2(s+r/2)}}  \, \d y
	\lesssim \frac{e^{-\frac12\lambda_1(2s+r)}}{(2s+r)^{d/2}}
\end{align*}
uniformly for all $s \in (0,T]$, $r > 0$ and $x,x'\in D$.
On the other hand, by Lemma \ref{lem:curve}, for any $x,x'\in D$, we can find a $C^2$ curve $\gamma:[0,1] \to D$ going from $x$ to $x'$ such that $|\gamma'(\tau)| \le C_0|x'-x|$ for any $\tau \in [0,1]$.
Then, by the fundamental theorem of calculus, Cauchy-Schwarz inequality and Lemma \ref{lem:grad:P}, there exists $C>0$ such that
\begin{align*}
	&\int_D |P_{s+r/2}(x',y)-P_{s+r/2}(x,y)|^2 \, \d y
	= \int_D\left| \int_0^1 \nabla_x P_{s+r/2}(\gamma(\tau),y) \cdot \gamma'(\tau) \, \d \tau \right|^2 \,  \d y \\
	& \lesssim \int_D \d y \left( \int_0^1 \d \tau\, |\gamma'(\tau)|^2 \right) \left( \int_0^1 \d \tau  \, |\nabla_x P_{s+r/2}(\gamma(\tau),y)|^2 \right)\\
	& \lesssim |x'-x|^2 \int_0^1 \d \tau \int_{\R^d} \d y \, \frac{e^{-\lambda_1(s+r/2)}}{(s+r/2)^{d+1}} e^{- \frac{2|\gamma(\tau)-y|^2}{C(s+r/2)}}
	\lesssim |x'-x|^2 \frac{e^{-\frac12 \lambda_1 (2s+r)}}{(2s+r)^{d/2+1}}
\end{align*}
uniformly for all $s \in (0,T]$, $r>0$ and $x,x' \in D$ with $|x'-x| \le \varepsilon_0$.
Combining these two bounds, we have
\begin{align*}
	\int_D |P_{s+r/2}(x',y)-P_{s+r/2}(x,y)|^2 \, \d y
	\lesssim
	 \frac{e^{-\frac12 \lambda_1 (2s+r)}}{(2s+r)^{d/2}} \left( 1 \wedge \frac{|x'-x|^2}{2s+r} \right),
\end{align*}
uniformly for all $s \in (0,T]$, $r>0$ and $x,x' \in D$ with $|x'-x| \le \varepsilon_0$. 
It follows that
\begin{align*}
	\Var(u(t,x')-u(t,x))
	& \lesssim \int_0^t \d s \int_0^\infty \d r \,\frac{r^{\alpha-1} e^{-\frac12\lambda_1(2s+r)}}{(2s+r)^{d/2}} \left( 1 \wedge \frac{|x'-x|^2}{2s+r} \right)\\
	&=\int_0^t \d s \int_{2s}^\infty \d r \,\frac{(r-2s)^{\alpha-1} e^{-\frac12\lambda_1r}}{r^{d/2}} \left( 1 \wedge \frac{|x'-x|^2}{r} \right)\\
	&=\int_0^\infty \d r \left[\int_0^{t \wedge (r/2)} (r-2s)^{\alpha-1} \, \d s \right] \frac{e^{-\frac12 \lambda_1r}}{r^{d/2}} \left( 1 \wedge \frac{|x'-x|^2}{r}\right)\\
	& \lesssim \int_0^\infty \d r \, r^{\alpha-d/2} e^{-\frac12 \lambda_1 r} \left( 1 \wedge \frac{|x'-x|^2}{r}\right).
\end{align*}
Next, we split the integral and estimate as follows:
\begin{align*}
	\Var(u(t,x')-u(t,x))
	& \lesssim \int_0^{|x'-x|^2} r^{\alpha-d/2} \, \d r + |x'-x|^2 \int_{|x'-x|^2}^1 r^{\alpha-d/2-1} \, \d r\\
	& \qquad  + |x'-x|^2 \int_1^\infty r^{\alpha-d/2-1} e^{-\frac12 \lambda_1 r} \, \d r\\
	& \lesssim \begin{cases}
	|x'-x|^{2\alpha-d+2} & \text{if $d/2-1<\alpha<d/2$,}\\
	|x'-x|^2 \log_+(1/|x'-x|) & \text{if $\alpha=d/2$,}\\
	|x'-x|^2 & \text{if $\alpha>d/2$.}
	\end{cases}
\end{align*}
This completes the proof of Proposition \ref{prop:u-u:x}, assuming Lemma \ref{lem:curve}.
\end{proof}

\begin{proof}[Proof of Lemma \ref{lem:curve}]
Since $D$ is a $C^k$ domain, for every $q \in \partial D$, we can find a radius $\varepsilon \in (0,1)$, a $C^k$ function $\phi:\R^{d-1} \to \R$, and an orthonormal coordinate system $(z,v)_q$ for $\R^d$, with $z \in \R^{d-1}, v \in \R$ and with origin $(0,0)_q=q$ such that
\begin{align*}
	&D \cap B(q,\varepsilon) = B(q,\varepsilon) \cap \{ (z,v+\phi(z))_q : z \in \R^{d-1}, v>0 \},\\
	&\partial D \cap B(q,\varepsilon) = B(q,\varepsilon) \cap \{ (z, \phi(z))_q : z \in \R^{d-1} \}.
\end{align*}
Since $D$ is bounded, its boundary $\partial D$ is compact, so we can find $q_1,\dots,q_n \in \partial D$ such that $\partial D \subset \bigcup_{i=1}^n B(q_i, \varepsilon_i/12)$, where $\varepsilon_1,\dots,\varepsilon_n \in (0,1)$ are the corresponding radii, $\phi_1,\dots, \phi_n :\R^{d-1} \to \R$ are the corresponding $C^1$ functions, and $(z,v)_i$ are the corresponding coordinate systems.

Since each $\phi_i$ is $C^k$ for some $k \ge 1$,
\begin{align}\label{curve:L}
	L_i := \sup_{|z| \le \varepsilon_i} |\nabla \phi_i(z)| < \infty.
\end{align}
We may take 
\begin{align}\label{curve:eps_0}
	C_0 := 2\left[1+\max_{1 \le i \le n}L_i^2\right]
	\quad \text{and}\quad
	\varepsilon_0 := \min_{1\le i\le n} \left[ \frac{\varepsilon_i}{12} \wedge \frac{\varepsilon_i}{4C_0} \right].
\end{align}
Consider any pair $x, y \in D$ with $|x-y| \le \varepsilon_0$.
Since $\overline{B}(x,\varepsilon_0) \cap \bigcup_{i=1}^n B(q_i, \varepsilon_i/12)$ is either empty or nonempty, and since $\varepsilon_0 \le \varepsilon_i/12$, either one of the following holds:
\begin{enumerate}
\item $\overline{B}(x, \varepsilon_0) \subset D$, or
\item $\overline{B}(x, \varepsilon_0) \subset B(q_i, \varepsilon_i/4)$ for some $i \in \{1,\dots, n\}$.
\end{enumerate}

In case (1), we may take $\gamma$ as the straight line from $x$ to $y$, and get $|\gamma'(\tau)| \le |x-y|$ for all $\tau \in [0,1]$.
In case (2), we may write 
\[
	x = F(z_1,v_1) \quad \text{and} \quad 
	y = F(z_2,v_2)
\]
for some $z_1, z_2 \in \R^{d-1}$ and $v_1,v_2>0$, where $F: \R^{d-1} \times \R \to \R^d$ is given by
\[
	F(z,v) = (z, v+ \phi_i(z))_i.
\]
Define $\gamma: [0,1] \to \R^d$ by
\[
	\gamma(\tau) = (\tau z_2+(1-\tau)z_1 \ ,\  \tau v_2+(1-\tau)v_1 + \phi_i(\tau z_2+(1-\tau)z_1))_i.
\]
Then, $\gamma$ is a $C^k$ curve such that $\gamma(0) = x$, $\gamma(1)=y$, and
\begin{align*}
	\gamma'(\tau) &= ( z_2-z_1\ , \ v_2-v_1 + \nabla \phi_i(\tau z_2+(1-\tau)z_1) \cdot (z_2-z_1) )_i.
\end{align*}
By \eqref{curve:L}, we have
\begin{align}\begin{split}\label{gamma':bd}
	|\gamma'(\tau)| &\le \sqrt{ |z_1-z_2|^2 + 2|v_1-v_2|^2+2L_i^2 |z_1-z_2|^2}\\
	& \le \sqrt{2(1+L_i^2)}\, |(z_1,v_1)-(z_2,v_2)|
\end{split}\end{align}
for all $\tau \in [0,1]$.
Note that $F$ is invertible with $F^{-1}(a,b) = (a,b-\phi_i(a))$ and
\begin{align*}
	|F^{-1}(a_1,b_1) - F^{-1}(a_2,b_2)|
	&\le \sqrt{2(1+L_i^2)}\, |(a_1,b_1)-(a_2,b_2)|
\end{align*}
for all $(a_1,b_1), (a_2,b_2) \in \R^{d-1} \times \R$.
It follows that
\[
	|(z_1,v_1)-(z_2,v_2)| = |F^{-1}(x)-F^{-1}(y)| \le \sqrt{2(1+L_i^2)} \, |x-y|.
\]
The preceding together with \eqref{gamma':bd} and the choice of $C_0$ in \eqref{curve:eps_0} implies that
\begin{align}\label{gamma':bd2}
	|\gamma'(\tau)| \le 2(1+L_i^2) |x-y| \le C_0|x-y| \quad \text{for all $\tau \in [0,1]$.}
\end{align}
This yields the desired estimate.
It remains to verify that the curve $\gamma$ lies in $D$.
Indeed, \eqref{gamma':bd2}, $|x-y| \le \varepsilon_0$, and $\overline{B}(x,\varepsilon_0) \subset B(q_i,\varepsilon_i/4)$ imply that for every $\tau \in [0,1]$, 
\begin{align*}
	|\gamma(\tau)-q_i| &\le |\gamma(\tau)-x| + |x-q_i| = \left|\int_0^\tau \gamma'(r) \, \d r\right| + |x-q_i|\\
	&\le C_0 \varepsilon_0 + \frac{\varepsilon_i}{4}
	\le C_0\left(\frac{\varepsilon_i}{4C_0}\right) + \frac{\varepsilon_i}{4} = \frac{\varepsilon_i}{2},
\end{align*}
where the last inequality follows from the choice of $\varepsilon_0$ in \eqref{curve:eps_0}. This shows that $\gamma(\tau) \in \overline{B}(q_i,\varepsilon_i/2) \cap \{ (z,v+\phi_i(z))_i : z \in \R^{d-1}, v>0 \} \subset D \cap \overline{B}(q_i,\varepsilon_i/2)$ for all $\tau \in [0,1]$.
The proof of Lemma \ref{lem:curve} is thus complete.
\end{proof}

%\begin{lemma}[Spatial increment upper bound]\label{lem:space}
%Assume \eqref{eq:alpha}. For every \(T>0\),
%\[
%    \E|u(t,x)-u(t,y)|^2\le C|x-y|^{4\theta},
%    \qquad t\in[0,T],\ x,y\in D .
%\]
%\end{lemma}
%
%\begin{proof}
%Let \(h=|x-y|\). The case \(h=0\) is trivial. By the isometry and the Green-kernel representation,
%\[
%\begin{aligned}
%    \Var(u(t,x)-u(t,y))
%    &\lesssim
%    \int_0^t\int_0^\infty q^{\alpha-1}
%    \inner{F_{r,h}}{P_qF_{r,h}}_{L^2(D)}\,\d q\,\d r,
%\end{aligned}
%\]
%where \(F_{r,h}(z)=P_r(x,z)-P_r(y,z)\).
%
%If \(0<r\le h^2\), then
%\[
%    \|F_{r,h}\|_1^2\le C,
%    \qquad
%    \|F_{r,h}\|_2^2\le C r^{-d/2}.
%\]
%Indeed, sub-Markovianity gives \(\|F_{r,h}\|_1\le2\), while the semigroup identity and the diagonal heat-kernel bound give
%\[
%    \|F_{r,h}\|_2^2
%    \le2P_{2r}(x,x)+2P_{2r}(y,y)\le Cr^{-d/2}.
%\]
%If \(h^2<r\le T\), quasiconvexity of bounded \(C^2\)-domains allows us to join \(x\) and \(y\) by an arclength curve \(\gamma\) of length at most \(Ch\); see \cite{BB12}. The fundamental theorem of calculus and the gradient heat-kernel estimate give
%\[
%    \|F_{r,h}\|_1^2\le C h^2r^{-1},
%    \qquad
%    \|F_{r,h}\|_2^2\le C h^2r^{-d/2-1}.
%\]
%To see these bounds, integrate \(\nabla_\xi P_r(\xi,z)\) along \(\gamma\), use Minkowski's inequality, and apply
%\[
%    \sup_{\xi\in D}\int_D|\nabla_\xi P_r(\xi,z)|\,\d z\le Cr^{-1/2},
%    \qquad
%    \sup_{\xi\in D}\int_D|\nabla_\xi P_r(\xi,z)|^2\,\d z
%       \le Cr^{-d/2-1},
%\]
%which follow by integrating the gradient estimate in Lemma~\ref{lem:heat}. Combining these estimates with Lemma~\ref{lem:qform}, we first get, for \(r\le h^2\),
%\[
%    \int_0^\infty q^{\alpha-1}
%    \inner{F_{r,h}}{P_qF_{r,h}}\,\d q
%    \lesssim
%    \int_0^\infty q^{\alpha-1}(r^{-d/2}\wedge q^{-d/2})\,\d q
%    \lesssim r^{\alpha-d/2}.
%\]
%Thus the contribution of \(0<r\le t\wedge h^2\) is at most
%\(Ch^{2\alpha+2-d}\). For \(r>h^2\),
%\[
%    \int_0^\infty q^{\alpha-1}
%    \inner{F_{r,h}}{P_qF_{r,h}}\,\d q
%    \lesssim
%    h^2 r^{-1-d/2+\alpha},
%\]
%and hence, when \(h^2<t\),
%\[
%    \int_{h^2}^{t} h^2r^{-1-d/2+\alpha}\,\d r
%    \lesssim h^{2\alpha+2-d},
%\]
%because \(\alpha<d/2\). Therefore
%\[
%    \Var(u(t,x)-u(t,y))    \lesssim h^{2\alpha+2-d}=h^{4\theta}.
%\]
%The \(q\)-integrals are finite because the standing regime \eqref{eq:alpha} has \(\alpha>0\) at zero and \(\alpha<d/2\) at infinity; the subsequent \(r\)-integrals use \(\alpha>d/2-1\) and \(\alpha<d/2\), respectively.
%\end{proof}


%\begin{lemma}\label{lem:boundary}
%Assume that \(D\) is a bounded \(C^2\)-domain and that \eqref{eq:alpha} holds. For every \(T>0\), there exists \(C>0\) such that
%\[
%    \sup_{0\le t\le T}\E|u(t,x)|^2
%    \le C d(x)^{4\theta},
%    \qquad x\in D .
%\]
%Since \(f_1(x)\asymp d(x)\) on bounded \(C^2\)-domains, equivalently
%\[
%    \sup_{0\le t\le T}\E|u(t,x)|^2
%    \le C f_1(x)^{4\theta}.
%\]
%Consequently, for every \(0<a<T<\infty\),
%\[
%    \sup_{\substack{t\in[a,T]\\x\in\partial_\delta D}}
%    \E|u(t,x)|^2
%    \le C\delta^{4\theta}
%    \xrightarrow[\delta\downarrow0]{}0 .
%\]
%\end{lemma}
%
%\begin{proof}
%It is enough to consider \(t>0\). By the stochastic isometry, the Green-kernel representation, and the semigroup property,
%\[
%    \E|u(t,x)|^2
%    =
%    \frac1{\Gamma(\alpha)}
%    \int_0^t\int_0^\infty
%    q^{\alpha-1}P_{2(t-r)+q}(x,x)\,\d q\,\d r .
%\]
%Put \(r_x=d(x)\). With \(\eta=t-r\) and \(\tau=2\eta+q\),
%\[
%\begin{aligned}
%    \E|u(t,x)|^2
%    &\lesssim
%    \int_0^\infty
%    \left[
%        \int_0^{t\wedge \tau/2}(\tau-2\eta)^{\alpha-1}\,\d\eta
%    \right]
%    P_\tau(x,x)\,\d\tau                                      \\
%    &\lesssim
%    \int_0^\infty
%    \tau^\alpha P_\tau(x,x)\,\d\tau .
%\end{aligned}
%\]
%The same display without the boundary factor, together with \(\alpha>d/2-1\) and the large-time exponential decay of the heat kernel, shows that
%\[
%    \sup_{\substack{0\le t\le T\\x\in D}}\E|u(t,x)|^2<\infty .
%\]
%Hence the desired estimate is immediate when \(r_x\ge1\). We may assume \(0<r_x<1\). Using the boundary-weighted Dirichlet heat-kernel estimate from Lemma~\ref{lem:heat}, with the large-time exponential factor, we get
%\[
%    P_\tau(x,x)
%    \lesssim
%    \left(1\wedge\frac{r_x}{\sqrt \tau}\right)^2
%    \tau^{-d/2}\e^{-c(\tau-1)_+}.
%\]
%Therefore
%\[
%\begin{aligned}
%    \E|u(t,x)|^2
%    &\lesssim
%    \int_0^{r_x^2}\tau^{\alpha-d/2}\,\d\tau
%    +
%    r_x^2\int_{r_x^2}^{1}\tau^{\alpha-d/2-1}\,\d\tau       \\
%    &\qquad
%    +
%    r_x^2\int_1^\infty\tau^{\alpha-d/2-1}\e^{-c(\tau-1)}\,\d\tau .
%\end{aligned}
%\]
%The first two terms are bounded by \(Cr_x^{2\alpha+2-d}\), while the last is bounded by \(Cr_x^2\). Since \(2\alpha+2-d=4\theta\in(0,2)\) in \eqref{eq:alpha} and \(r_x\le\diam(D)\),
%\[
%    r_x^2
%    =
%    r_x^{4\theta}r_x^{2-4\theta}
%    \le
%    \diam(D)^{2-4\theta}r_x^{4\theta}.
%\]
%This proves the \(d(x)^{4\theta}\) bound. The comparison with \(f_1(x)^{4\theta}\) follows from the first-eigenfunction boundary comparison \cite[Remark~1.1(b)]{Z02}.
%\end{proof}

%\begin{proposition}\label{prop:continuous}
%Extend the field to \([0,\infty)\times\overline D\) by setting \(u(0,x)=0\) and \(u(t,x)=0\) for \(x\in\partial D\). Under the standing assumptions, it has a separable modification that is jointly continuous on \([0,\infty)\times\overline D\). For every \(T>0\),
%\[
%    d_u(z,z')\le C_T\rho(z,z'),
%    \qquad z,z'\in[0,T]\times\overline D.
%\]
%All suprema in the sequel are taken for this version and are measurable.
%\end{proposition}
%
%\begin{proof}
%For two interior points the displayed estimate follows by decomposing the increment into a temporal and a spatial part and applying Lemmas~\ref{lem:time} and \ref{lem:space}. If one spatial point lies on \(\partial D\), Lemma~\ref{lem:boundary} and \(d(y)\le |x-y|\) give the same bound. If one time is zero, the endpoint case of Lemma~\ref{lem:time} applies. Thus the estimate holds on the metric completion.
%
%The \(\rho\)-covering numbers of \([0,T]\times\overline D\) are bounded by \(Cr^{-Q}\). Since \(d_u\le C_T\rho\), the Dudley entropy integral
%\[
%    \int_0^1\sqrt{\log N_{d_u}([0,T]\times\overline D,r)}\,\d r
%\]
%is finite. The Gaussian continuity theorem therefore gives a separable, uniformly continuous modification on each such compact parameter set. Apply this construction on \([0,m]\times\overline D\), \(m\in\mathbb N_+\), and use indistinguishability on overlaps (first on a fixed countable dense set, then by continuity) to obtain one jointly continuous version on the full parameter space.
%\end{proof}
%
%\begin{remark}[On the full-domain uniform modulus]\label{rem:boundary}
%Lemma~\ref{lem:boundary} is the boundary-layer estimate obtained from the Dirichlet heat-kernel decay. It does not imply the stronger negligibility statement
%\[
%    \lim_{\delta\downarrow0}\limsup_{\varepsilon\downarrow0}
%    \sup_{\substack{z,z'\in[a,T]\times\partial_\delta D\\
%    0<\rho(z,z')\le\varepsilon}}
%    \frac{|u(z)-u(z')|}
%    {\rho(z,z')\sqrt{\log(1/\rho(z,z'))}}
%    =0 .
%\]
%Indeed, for every fixed \(\delta>0\), the layer \(\partial_\delta D\) contains points at positive distance from \(\partial D\), and at scales much smaller than that distance the local increments are governed by the same interior canonical metric. Thus the compact-interior exact modulus below cannot be upgraded to \([a,T]\times D\) by simply replacing \(D_\delta\) with \(D\). Such an upgrade would require a scale-dependent decomposition and control of pairs whose boundary distances are comparable to their increment scale.
%\end{remark}

%\subsection{Canonical metric and variance estimates}
%
%Canonical metric estimates for stochastic heat equations with spatially correlated noise are closely related to the sharp regularity theory developed in \cite{HSWX20,Chen23,LX25}; the bounded-domain estimates below also use the Dirichlet heat-kernel bounds from \cite{Davies,Z06,S04,Daners}.
%
%\begin{proposition}[Sharp variance estimates]\label{prop:metric}
%Assume that \(D\) is a bounded \(C^2\)-domain and that
%\eqref{eq:alpha} holds. For every \(T>0\),
%\[
%    d_u((t,x),(s,y))^2
%    \le C\rho((t,x),(s,y))^2
%\]
%uniformly for \((t,x),(s,y)\in(0,T]\times D\). In particular, for all \(0<a<T\) and \(\delta>0\) such that \(D_\delta\ne\varnothing\),
%\[
%    d_u((t,x),(s,y))\asymp \rho((t,x),(s,y)),
%    \qquad (t,x),(s,y)\in[a,T]\times D_\delta .
%\]
%The lower bound will only be used on compact interior sets, where the initial-time and boundary factors in the SLND estimate are uniformly bounded away from zero.
%\end{proposition}
%
%\begin{proof}
%The upper bound follows by decomposing the increment into temporal and spatial parts and applying Lemmas~\ref{lem:time} and \ref{lem:space}. For the interior lower bound, fix \((t,x),(s,y)\in[a,T]\times D_\delta\). Proposition~\ref{prop:SLND}, applied with one comparison point, gives
%\[
%    \Var(u(t,x)\mid u(s,y))
%    \ge
%    c\left(\rho((t,x),(s,y))^2\wedge a^{2\theta}
%        \wedge \inf_{z\in D_\delta}f_1(z)^{4\theta}\right).
%\]
%Since \(D\) is \(C^2\), \(f_1(z)\asymp d(z)\) by \cite[Remark~1.1(b)]{Z02}, and the last infimum is positive. Also \(\rho\) is bounded on \([a,T]\times D_\delta\). Hence
%\[
%    \Var(u(t,x)-u(s,y))
%    \ge \Var(u(t,x)\mid u(s,y))
%    \ge c_{a,T,\delta}\rho((t,x),(s,y))^2 .
%\]
%Here the first inequality is the Hilbert-space projection identity
%\[
%    \Var(u(t,x)\mid u(s,y))
%    =\inf_{b\in\R}\Var(u(t,x)-b u(s,y)),
%\]
%with \(b=1\) as an admissible choice. If \(R_I\) is the \(\rho\)-diameter of the cylinder and \(m=\min\{a^{2\theta},\inf_{D_\delta}f_1^{4\theta}\}\), then explicitly \(\rho^2\wedge m\ge\min\{1,m/R_I^2\}\rho^2\).
%This yields the asserted two-sided canonical-metric estimate on compact interior sets.
%\end{proof}

\subsection{Proof of Theorem \ref{thm:reg}}

\begin{proof}
Fix $T>0$.
Under \eqref{Dalang}, i.e., $\alpha>d/2-1$,
Propositions \ref{prop:u-u:t} and \ref{prop:u-u:x} imply that
\begin{align*}
	\|u(t,x)-u(t',x')\|_2
	&\le \|u(t,x)-u(t',x)\|_2 + \|u(t',x)-u(t',x')\|_2\\
	& \lesssim \tilde\rho((t,x),(t',x')),
\end{align*}
uniformly for all $(t,x),(t',x') \in [0,T]\times D$, where 
\[
	\tilde\rho((t,x),(t',x')) :=
	\begin{cases}
	|t-t'|^{\theta/2} + |x-x'|^\theta & \text{if $\frac{d}{2}-1 < \alpha < \frac{d}{2}$,}\smallskip\\
	\sqrt{|t-t'| \log_+\frac{1}{|t-t'|}} + |x-x'|\sqrt{\log_+\frac{1}{|x-x'|}} & \text{if $\alpha=d/2$,}\smallskip\\
	\sqrt{|t-t'|} + |x-x'| & \text{if $\alpha>d/2$,}
	\end{cases}
\]
and $\theta$ is given by \eqref{theta}.
Since $u$ is Gaussian, it follows that for any $k \ge 2$, there exists $C_{k,T} >0$ such that
\[
	\|u(t,x)-u(t',x')\|_k \le C_{k,T} \,\tilde\rho((t,x),(t',x')),
\]
uniformly for all $(t,x),(t',x') \in [0,T]\times D$.
Hence, by Kolmogorov's continuity theorem, $u$ is a.s.~locally H\"older continuous of order $\beta_0$ in $t$ and of order $\beta_1$ in $x$, for any $\beta_0 \in (0,(\theta\wedge 1)/2)$ and $\beta_1 \in (0,\theta\wedge 1)$.
\end{proof}

\section{Optimal regularity}\label{s:optimal}

\subsection{Strong local nondeterminism}

%Strong local nondeterminism is a central tool in the sample path analysis of Gaussian random fields; see \cite{Berman,Berman87,Pitt,Xiao08}. 
%The localized spectral-duality argument used here is adapted to the bounded-domain heat equation and is compatible with the anisotropic framework of \cite{LX23,LX25}; it is also closely related to the linear analysis in \cite{HL26}.

\begin{proposition}\label{prop:SLND}
Suppose \(D\subset\R^d\) is a bounded \(C^2\) domain and \eqref{alpha:Holder:range} holds. Then for every \(T>0\), there exists \(C>0\) such that
\[
    \Var\left(
        u(t,x)\mid u(t_1,x_1),\ldots,u(t_n,x_n)
    \right)
    \ge
    C\left[
        \min_{1\le j\le n}
        \rho^2((t,x),(t_j,x_j))
        \wedge t^{\theta}
        \wedge (d(x))^{2\theta}
    \right]
\]
uniformly for all \(n\ge1\) and all \((t,x),(t_1,x_1),\dots,(t_n,x_n)\in[0,T]\times D\).
\end{proposition}

\begin{proof}
For a centered Gaussian vector $(X,Y_1,\dots, Y_n)$, the conditional variance $\Var(X\mid Y_1,\dots,Y_n)$ is the squared \(L^2(\Omega)\)-distance from $X$ to the linear span of the conditioned variables, i.e.,
\[
	\Var(X\mid Y_1,\dots,Y_n) =\inf_{a_1,\dots, a_n \in \R}\E\Bigg[\bigg(X-\sum_{j=1}^n a_jY_j\bigg)^2\Bigg].
\]
Thus it is enough to prove that there exists $C>0$ such that
\begin{align}\begin{split}\label{SLND:claim}
	&\E\Bigg[\bigg(u(t,x)-\sum_{j=1}^n a_j u(t_j,x_j)\bigg)^2\Bigg]\\
	& \ge C \left[ \min_{1\le j\le n}
        \rho^2((t,x),(t_j,x_j))
        \wedge t^{\theta}
        \wedge (d(x))^{2\theta}\right],
\end{split}\end{align}
uniformly for all $n \ge 1$, $(t,x),(t_1,x_1),\dots, (t_n,x_n) \in [0,T] \times D$, and $a_1,\dots, a_n \in \R$.
To this end, we first recall that
\[
	u(t,x)-\sum_{j=1}^n a_j u(t_j,x_j) = W_\alpha\bigg( p_{t,x} - \sum_{j=1}^n a_j p_{t_j,x_j}\bigg).
\]
By Theorem \ref{thm:exist}, the right-hand side above is a well-defined Wiener integral under condition \eqref{Dalang}.
Hence, we may use the first identities of \eqref{ip:alpha} and \eqref{W:corr:WfWg}, and \eqref{P} to deduce that 
\begin{align}\label{eq:Fourier}
&\E\Bigg[\bigg(u(t,x)-\sum_{j=1}^n a_j u(t_j,x_j)\bigg)^2\Bigg]\\
&= \sum_{k=1}^\infty \lambda_k^{-\alpha} \int_\R \bigg| e^{-\lambda_k (t-r)} f_k(x) {\bf 1}_{(0,t)}(r) - \sum_{j=1}^n a_j e^{-\lambda_k (t_j-r)} f_k(x_j) {\bf 1}_{(0,t_j)}(r) \bigg|^2 \d r \notag\\
&=
    \frac1{2\pi}
    \sum_{k=1}^\infty \lambda_k^{-\alpha}
    \int_\R
    \frac{\left|
        (e^{-i\tau t}-e^{-\lambda_k t})f_k(x)
        -\sum_{j=1}^n a_j
        (e^{-i\tau t_j}-e^{-\lambda_k t_j})f_k(x_j)
    \right|^2}
    {\lambda_k^2+\tau^2}\,\d\tau, \notag
\end{align}
where the last equality follows from Plancherel's theorem and the identity
\[
    \int_0^t \e^{-\lambda(t-r)}\e^{-i\tau r}\,\d r
    =
    \frac{\e^{-i\tau t}-\e^{-\lambda t}}{\lambda-i\tau}.
\]
Let \(R_D=\diam(D)\) and \(A_D=R_D\wedge1\wedge\sqrt T\). 
Define
\begin{align}\label{SLND:r}
    r
    :=
    A_D
    \left[
        \min_{1\le j\le n}
        \left(
            \sqrt{\frac{|t-t_j|}{T}}
            \vee
            \frac{|x-x_j|}{R_D}
        \right)
        \wedge
        \sqrt{\frac{t}{T}}
        \wedge
        \frac{d(x)}{R_D}
    \right].
\end{align}
Choose and fix two test functions
\[
    \varphi\in C_c^\infty((-1/2,1/2)),
    \qquad
    \psi\in C_c^\infty(B(0,1/2)),
\]
with \(\varphi(0)=\psi(0)=1\).
If \(r=0\), the desired lower bound in \eqref{SLND:claim} is trivial. So, we may assume that \(r>0\), and define
\[
    \varphi_{r^2}(u):=r^{-2}\varphi(u/r^2)
    \quad \text{and} \quad
    \psi_{x,r}(y):=r^{-d}\psi\left(\frac{y-x}{r}\right).
\]
Following \cite[Section~3]{HL26}, we define
\begin{align}\label{SLND:I}
    I&:=
    \sum_{k=1}^\infty \inner{\psi_{x,r}}{f_k}_{L^2(D)} \times\\
    & \quad \times
    \int_\R
    \left[
        (\e^{-i\tau t}-\e^{-\lambda_k t})f_k(x)
        -
        \sum_{j=1}^n a_j
        (\e^{-i\tau t_j}-\e^{-\lambda_k t_j})f_k(x_j)
    \right]
    \e^{i\tau t}\widehat\varphi_{r^2}(\tau)\,\d\tau,\notag
\end{align}
where we use the Fourier convention
\[
    \widehat g(\tau)=\int_\R\e^{-i\tau v}g(v)\,\d v,
    \qquad
    g(v)=\frac1{2\pi}\int_\R\e^{i\tau v}\widehat g(\tau)\,\d\tau.
\]
In particular, \(\widehat\varphi_{r^2}(\tau)=\widehat\varphi(r^2\tau)\), and Fourier inversion gives
\begin{align*}
    I
    =2\pi\sum_{k=1}^\infty
    \Bigg[&\bigl(\varphi_{r^2}(0)
        -\e^{-\lambda_k t}\varphi_{r^2}(t)\bigr)f_k(x)\\
    &-\sum_{j=1}^n a_j
       \bigl(\varphi_{r^2}(t-t_j)
        -\e^{-\lambda_k t_j}\varphi_{r^2}(t)\bigr)f_k(x_j)
    \Bigg]
    \inner{\psi_{x,r}}{f_k}_{L^2(D)}.
\end{align*}
Note that \(A_D\le\sqrt T\), thus we have \(r\le\sqrt t\). 
Because \(\operatorname{supp}\varphi\subset[-1/2,1/2]\), the inequality \(r\le\sqrt t\) implies $\varphi_{r^2}(t)=0$.
Hence,
\begin{align*}
	I
	=2\pi\sum_{k=1}^\infty
	\Bigg[\varphi_{r^2}(0) f_k(x)
	-\sum_{j=1}^n a_j\varphi_{r^2}(t-t_j)f_k(x_j)\Bigg]
	\inner{\psi_{x,r}}{f_k}_{L^2(D)}.
\end{align*}
Since \(\psi_{x,r}\in C_c^\infty(D)\), we may use integration by parts, $\sup_{x \in D}|f_n(x)|\lesssim \lambda_n^{d/4}$ (see \cite[Proposition 5]{F95}) and Weyl's law $\lambda_n \asymp n^{2/d}$ (see \cite[Corollary 3.15]{FLW}) to show that for any $z \in D$,
\begin{align*}
	&\sum_{n=1}^\infty \left|\langle \psi_{x,r} , f_n \rangle_{L^2(D)} f_n(z)\right|
	= \sum_{n=1}^\infty \lambda_n^{-d} \left|\langle \psi_{x,r} , (-\Delta)^d f_n \rangle_{L^2(D)}\right| |f_n(z)| \\
	&= \sum_{n=1}^\infty \lambda_n^{-d} \left| \langle(-\Delta)^d \psi_{x,r}, f_n  \rangle_{L^2(D)}\right| |f_n(z)|
	\lesssim \|(-\Delta)^d \psi_{x,r}\|_{L^2(D)} \sum_{n=1}^\infty n^{-3/2} < \infty.
\end{align*}
The preceding and the fact that each $f_n$ is continuous on $D$ (see \cite[Theorem 9.31]{Brezis}) then imply pointwise convergence of the following series:
\[
    \sum_{k=1}^\infty
    \inner{\psi_{x,r}}{f_k}_{L^2(D)}f_k(z)
    =\psi_{x,r}(z) \quad \text{for all $z\in D$.}
\]
It follows that
\begin{align}\label{SLND:I:eval}
    I=2\pi
    \Bigg[\varphi_{r^2}(0)\psi_{x,r}(x)-\sum_{j=1}^n a_j \varphi_{r^2}(t-t_j)\psi_{x,r}(x_j)\Bigg].
\end{align}
Next, we evaluate $I$ by checking the support of the test functions.
%Since \(A_D\le\sqrt T\), we have \(r\le\sqrt t\), hence \(t\ge r^2\). 
%Because \(\operatorname{supp}\varphi\subset[-1/2,1/2]\), the inequality \(t\ge r^2\) implies
%\begin{align}\label{phi=0}
%	\varphi_{r^2}(t)=0.
%\end{align}
%We next check the support separation for every comparison point. 
Recall \eqref{d(x)}.
Since \(A_D\le R_D\), we have \(r\le d(x)\), so $\psi_{x,r}(\cdot)$ is supported in \(B(x,r/2)\subset D\). 
Fix \(j\in\{1,\ldots,n\}\) and consider the product $\varphi_{r^2}(t-t_j)\psi_{x,r}(x_j)$.
If \(\varphi_{r^2}(t-t_j)=0\), then this product is 0. Otherwise \(|t-t_j|\le r^2/2\). Hence
\[
    \sqrt{\frac{|t-t_j|}{T}}\le \frac{r}{\sqrt{2T}}.
\]
By the definition of \(r\), the maximum
\[
    \sqrt{\frac{|t-t_j|}{T}}\vee\frac{|x-x_j|}{R_D}
\]
is at least \(r/A_D\).
But the first term $\sqrt{|t-t_j|/T}$ in the maximum is at most \(r/\sqrt{2T}<r/A_D\) because \(A_D\le\sqrt T\), so the maximum must be attained by the second term $|x-x_j|/R_D$. Hence,
\[
    r\le A_D\frac{|x-x_j|}{R_D}\le |x-x_j|.
\]
Since \(\operatorname{supp}\psi_{x,r}\subset B(x,r/2)\), this gives \(\psi_{x,r}(x_j)=0\). Therefore, in any case, we have
\begin{equation}\label{eq:vanish}
    \varphi_{r^2}(t-t_j)\psi_{x,r}(x_j)=0
    \quad \text{for each $j=1,\dots, n$.}
\end{equation}
%Following the localized Fourier-pairing argument in
%\cite[Section~3]{HL26}, define
%\[
%\begin{aligned}
%    I
%    :=
%    \sum_{k=1}^\infty
%    \int_\R
%    &\left[
%        (\e^{-i\tau t}-\e^{-\lambda_k t})f_k(x)
%        -
%        \sum_{j=1}^n a_j
%        (\e^{-i\tau t_j}-\e^{-\lambda_k t_j})f_k(x_j)
%    \right]  \\
%    &\hspace{2.5cm}\times
%    \e^{i\tau t}\widehat\varphi_{r^2}(\tau)
%    \inner{\psi_{x,r}}{f_k}_{L^2(D)}\,\d\tau .
%\end{aligned}
%\]
%The estimate for \(J\) below shows that this is a well-defined pairing in the
%weighted \(\ell^2\times L^2\) space. 
%Fourier inversion gives
%It follows that
%\[
%\begin{aligned}
%    I
%    =2\pi\sum_{k=1}^\infty
%    \Bigg[&\bigl(\varphi_{r^2}(0)
%        -\e^{-\lambda_k t}\varphi_{r^2}(t)\bigr)f_k(x)\\
%    &-\sum_{j=1}^n a_j
%       \bigl(\varphi_{r^2}(t-t_j)
%        -\e^{-\lambda_k t_j}\varphi_{r^2}(t)\bigr)f_k(x_j)
%    \Bigg]
%    \inner{\psi_{x,r}}{f_k}_{L^2(D)} .
%\end{aligned}
%\]
%Since \(\psi_{x,r}\in C_c^\infty(D)\), we have the following identity (see XXX):
%\[
%    \sum_{k=1}^\infty
%    f_k(z)\inner{\psi_{x,r}}{f_k}_{L^2(D)}
%    =\psi_{x,r}(z) \quad \text{for all $z\in D$.}
%\]
%To justify this identity, choose an integer \(m>d/4\). Since
%\(\psi_{x,r}\in\operatorname{Dom}(A^m)\),
%\[
%    \inner{\psi_{x,r}}{f_k}_{L^2(D)}
%    =\lambda_k^{-m}
%      \inner{A^m\psi_{x,r}}{f_k}_{L^2(D)}.
%\]
%Moreover, the heat-kernel representation and the spectral gap imply
%\[
%    \sum_{k=1}^\infty\lambda_k^{-2m}f_k(q)^2
%    =
%    \frac1{\Gamma(2m)}
%    \int_0^\infty s^{2m-1}P_s(q,q)\,\d s
%    <\infty .
%\]
%Cauchy--Schwarz therefore gives absolute convergence. For \(\eta>0\), the
%corresponding damped sum equals
%\[
%    \sum_{k=1}^\infty \e^{-\eta\lambda_k}f_k(q)
%      \inner{\psi_{x,r}}{f_k}_{L^2(D)}
%    =P_\eta\psi_{x,r}(q).
%\]
%The absolute convergence permits passage to the limit as \(\eta\downarrow0\),
%while the heat semigroup is an approximate identity on
%\(C_c^\infty(D)\); hence the limit is \(\psi_{x,r}(q)\). 
%By
%\eqref{eq:vanish} and \(\varphi_{r^2}(t)=0\), all the comparison-point
%terms vanish. Therefore
It follows from \eqref{SLND:I:eval} and \eqref{eq:vanish} that
\begin{align}\label{SLND:I:value}
    I
    =2\pi\varphi_{r^2}(0)\psi_{x,r}(x)
    =2\pi r^{-d-2}.
%    ,
%    \qquad
%    |I|^2\asymp r^{-2d-4}.
\end{align}

Applying Cauchy-Schwarz directly to \eqref{SLND:I}, with weight
\(\lambda_k^{-\alpha}(\lambda_k^2+\tau^2)^{-1}\), and using
\eqref{eq:Fourier}, gives
\begin{align}\label{SLND:I:CS}
    |I|^2
    \le
    C\,\E\Bigg[\bigg(u(t,x)-\sum_{j=1}^n a_j u(t_j,x_j)\bigg)^2\Bigg]\times J,
\end{align}
where
\[
    J
    :=
    \sum_{k=1}^\infty
    \int_\R
    \lambda_k^\alpha(\lambda_k^2+\tau^2)
    \left|\widehat\varphi_{r^2}(\tau)\right|^2
    \left|\inner{\psi_{x,r}}{f_k}_{L^2(D)}\right|^2\,\d\tau .
\]
Since \(\widehat\varphi_{r^2}(\tau)=\widehat\varphi(r^2\tau)\) and $\widehat{\varphi}$ is rapidly
decreasing, we have
\[
    \int_\R \left|\widehat\varphi_{r^2}(\tau)\right|^2\,\d\tau\lesssim r^{-2},
    \qquad
    \int_\R\tau^2\left|\widehat\varphi_{r^2}(\tau)\right|^2\,\d\tau\lesssim r^{-6}.
\]
Consequently,
\[
	J \lesssim r^{-2} \sum_{k=1}^\infty \lambda_k^{\alpha+2} \left| \langle \psi_{x,r}, f_k \rangle_{L^2(D)} \right|^2 + r^{-6} \sum_{k=1}^\infty \lambda_k^{\alpha} \left| \langle \psi_{x,r}, f_k \rangle_{L^2(D)} \right|^2.
\]
%\[
%    J
%    \lesssim
%    r^{-2}\|A^{\alpha/2+1}\psi_{x,r}\|_2^2
%    +
%    r^{-6}\|A^{\alpha/2}\psi_{x,r}\|_2^2 .
%\]
%Lemma~\ref{lem:scaling} gives
%\[
%    \|A^{\alpha/2+1}\psi_{x,r}\|_2^2\lesssim r^{-d-2\alpha-4},
%    \qquad
%    \|A^{\alpha/2}\psi_{x,r}\|_2^2\lesssim r^{-d-2\alpha}.
%\]
Next, we show that for any $\beta\ge 0$, there exists $C=C_{\beta, \psi}>0$ such that
\begin{align}\label{SLND:claim:J}
	\sum_{k=1}^\infty \lambda_k^{\beta} \left| \langle \psi_{x,r}, f_k \rangle_{L^2(D)} \right|^2 \le C r^{-d-2\beta}.
\end{align}

Case 1: $\beta = m \in \N_0$.
In this case, we may use $-\Delta f_k = \lambda_k f_k$, integration by parts, and the definition of $\psi_{x,r}$ to deduce that
\begin{align*}
	&\langle \psi_{x,r} , \lambda_k^m f_k\rangle
	= \int_D \psi_{x,r}(y) ((-\Delta)^m f_k) (y) \, \d y
	=\int_D ((-\Delta)^m \psi_{x,r}) (y) f_k(y)  \, \d y\\
	& = (-1)^m r^{-2m} \int_D r^{-d}(\Delta^m\psi)(\tfrac{y-x}{r}) f_k(y) \, \d y = (-1)^m r^{-2m} \langle (\Delta^m\psi)_{x,r} , f_k\rangle_{L^2(D)}.
\end{align*}
The preceding, Parseval's identity, and the change of variable $z=(y-x)/r$ then imply that
\begin{align*}
	&\sum_{k=1}^\infty \lambda_k^{m} \left| \langle \psi_{x,r}, f_k \rangle_{L^2(D)} \right|^2
	= \sum_{k=1}^\infty \langle \psi_{x,r}, f_k \rangle_{L^2(D)} \cdot \langle \psi_{x,r}, \lambda_k^m f_k \rangle_{L^2(D)}\\
	&= (-1)^m r^{-2m} \sum_{k=1}^\infty \langle \psi_{x,r}, f_k \rangle_{L^2(D)} \cdot \langle (\Delta^m\psi)_{x,r}, f_k \rangle_{L^2(D)}\\
	&= (-1)^m r^{-2m} \langle \psi_{x,r} , (\Delta^m \psi)_{x,r} \rangle_{L^2(D)}
	=(-1)^m r^{-2d-2m} \int_D \psi(\tfrac{y-x}{r}) (\Delta^m \psi)(\tfrac{y-x}{r}) \, \d y\\
	& \le r^{-d-2m} \int_{\R^d} |\psi(z)| |\Delta^m \psi (z)| \, \d z = C_{m,\psi}\, r^{-d-2m}.
\end{align*}

Case 2: $\beta \not\in \N_0$. In this case, we can find $m \in \N_0$ such that $m < \beta < m+1$.
Take $p = 1/(m+1-\beta)>1$ and $q = 1/(\beta-m)>1$, so that $1/p+1/q=1$.
Set $\beta_1 = m/p$ and $\beta_2 = (m+1)/q$.
Note that $\beta_1+\beta_2 = \beta$.
Hence, we may apply H\"older's inequality and Case 1 to deduce that
\begin{align*}
	&\sum_{k=1}^\infty \lambda_k^{\beta} \left| \langle \psi_{x,r}, f_k \rangle_{L^2(D)} \right|^2\\
%	=\sum_{k=1}^\infty \lambda_k^{\beta_1+\beta_2} \left| \langle \psi_{x,r}, f_k \rangle_{L^2(D)} \right|^{2/p + 2/q}\\
	&\le \left( \sum_{k=1}^\infty \lambda_k^m \left| \langle \psi_{x,r}, f_k \rangle_{L^2(D)} \right|^2\right)^{m+1-\beta} \left(\sum_{k=1}^\infty \lambda_k^{m+1} \left| \langle \psi_{x,r}, f_k \rangle_{L^2(D)} \right|^2\right)^{\beta-m}\\
	& \le \left(C_{m,\psi} \, r^{-d-2m} \right)^{m+1-\beta} \left( C_{m+1,\psi} \, r^{-d-2m-2} \right)^{\beta-m}
	= C_{\beta, \psi} r^{-d-2\beta}.
\end{align*}
This proves \eqref{SLND:claim:J}.
It follows that
\[
	J\lesssim r^{-d-2\alpha-6}.
\]
Now, we return to \eqref{SLND:I:CS} and use \eqref{SLND:I:value} to deduce that
\begin{align}\label{SLND:claim:proved}
    \E\Bigg[\bigg(u(t,x)-\sum_{j=1}^n a_j u(t_j,x_j)\bigg)^2\Bigg]
    \gtrsim
    r^{-2d-4} J^{-1}
    \gtrsim
    r^{2\alpha-d+2}
    =
    r^{2\theta},
\end{align}
uniformly for all $n \ge 1$, $(t,x),(t_1,x_1),\dots, (t_n,x_n) \in [0,T]\times D$ and $a_1,\dots, a_n \in \R$.
%Taking the infimum over the coefficients \(a_1,\dots, a_n \in \R\) yields
%\begin{align*}
%	\Var(u(t,x)\mid u(t_1,x_1),\ldots,u(t_n,x_n)) \gtrsim r^{2\theta}.
%%    &\qquad\gtrsim
%%    \left[
%%        \min_{1\le j\le n}
%%        \bigl(|t-t_j|^{1/2}\vee |x-x_j|\bigr)
%%        \wedge \sqrt t
%%        \wedge d(x)
%%        \wedge1
%%    \right]^{2\theta}.
%\end{align*}
The factors $T$ and $R_D$ in \eqref{SLND:r} can be enlarged and factored out to give a lower bound.
Moreover, since
%$(a\vee b)^{2\theta}\asymp(a^{\theta}+b^{\theta})^2$ and
\[
    \rho((t,x),(t_j,x_j))
    \asymp
    \bigl(|t-t_j|^{1/2}\vee |x-x_j|\bigr)^{\theta},
\]
the implicit constant above can also be factored and absorbed by $C$ to obtain \eqref{SLND:claim}.
%the fourth power of the right-hand side is comparable to
%\(\rho((t,x),(t_j,x_j))^2\). This uses the elementary equivalence
%\((a\vee b)^{4\theta}\asymp(a^{2\theta}+b^{2\theta})^2\).
%Finally, \(f_1(x)\asymp d(x)\) on bounded \(C^2\) domains (see \cite[Remark~1.1(b)]{Z02}). 
This completes the proof of the proposition.
%give exactly the three capped terms in the proposition.
\end{proof}

%\begin{proof}[Proof of Theorem~\ref{thm:SLND}]
%Combine Propositions~\ref{prop:metric} and \ref{prop:SLND}.
%\end{proof}

\subsection{Proof of Theorem \ref{thm:optimal}}

\begin{proof}
Under \eqref{Dalang}, the upper bound in \eqref{var:u-u} follows from Propositions \ref{prop:var:u}, \ref{prop:u-u:t}, \ref{prop:u-u:x}.
%, and the fact that $f_1(x) \asymp d(x)$ (see Lemma \ref{lem:P}).
The lower bound in \eqref{var:u-u} follows from an application of \eqref{SLND:claim:proved} with $n=1$, $a_1=1$ and $(t_1,x_1)=(s,y)$, and another application of \eqref{SLND:claim:proved} with $(t,x)$ and $(s,y)$ swapped.
The SLND property on $[\delta,T]\times D'$ follows from Proposition \ref{prop:SLND} and the variance bound \eqref{var:u-u}, which implies that $\Var(u(t,x)-u(s,y)) \asymp \rho((t,x),(s,y))^2$ uniformly for all $(t,x),(s,y)\in [\delta,T]\times D'$ since $t,s$ and $d(x),d(y)$ are all bounded above and below.
\end{proof}


\section{Proofs of Theorems \ref{thm:moc}--\ref{thm:sbp}}\label{s:proofs}

%Exact moduli of continuity for anisotropic Gaussian fields have been studied
%in \cite{MWX13,LX23}, building on classical Gaussian sample-path theory such
%as \cite{AT,MR}. Chung-type laws originate from Chung's classical theorem
%\cite{Chung}.
%small ball probabilities for Gaussian processes and fields are treated
%systematically in \cite{LS01,T95}. Related small ball estimates for stochastic
%heat equations and SPDEs appear in \cite{AJM21,Chen24,Chen25,KKM24}. 
%In the present Gaussian random-field setting we use the
%anisotropic formulation of \cite{LX23}.


\subsection{Harmonizable representation}

We use the framework of Lee and Xiao \cite{LX23} to prove Theorems \ref{thm:moc}--\ref{thm:sbp}. Their first assumption
requires an independently scattered Gaussian field indexed by Borel subsets
of \([0,\infty)\), together with quantitative estimates for its low- and
high-frequency parts. We construct this field from the temporal Fourier
transform of the heat kernel.

Let \(\{W_n^{(1)}\}_{n\ge1}\) and \(\{W_n^{(2)}\}_{n\ge1}\) be independent
sequences of real-valued centered Gaussian white noises on \(\R\), and set
\(W_n=W_n^{(1)}+iW_n^{(2)}\). Thus, for \(f,g\in L^2(\R;\mathbb C)\),
\[
    \E\left[\left(\Re\int_\R f\,\d W_n\right)
             \left(\Re\int_\R g\,\d W_m\right)\right]
    =\delta_{nm}\Re\int_\R f(\tau)\overline{g(\tau)}\,\d\tau,
\]
where $\Re$ denotes real part.
Define
\begin{align}\label{q}
    q_{n,t}(\tau)
    =\frac{\e^{-i\tau t}-\e^{-\lambda_n t}}{\lambda_n-i\tau}.
%    \qquad
%    \kappa_n(\tau)=\lambda_n^{\theta/2}\vee|\tau|^{\theta/2}.
\end{align}
%The factor
%\(\lambda_n^\theta\) is the spectral analogue of the spatial frequency
%\(|\xi|^{2\theta}\), since \(\lambda_n\) has the scaling of \(|\xi|^2\).
For any Borel set \(A\subset [0,\infty)\), define the centered Gaussian field
\begin{equation}\label{eq:harm}
    u(A,t,x)
    =\frac1{\sqrt{2\pi}}
    \sum_{n=1}^{\infty}\lambda_n^{-\alpha/2}f_n(x)
    \Re\int_{\{\tau \in \R:\,\lambda_n^{\theta/2}\vee|\tau|^{\theta/2}\in A\}}
    q_{n,t}(\tau)\,W_n(\d\tau).
\end{equation}
Since \(q_{n,t}\) is the Fourier transform of
\(r\mapsto\1_{(0,t)}(r)\e^{-\lambda_n(t-r)}\), Plancherel's theorem and Proposition \ref{prop:spec:rep} imply that, under condition \eqref{Dalang}, $\{u([0,\infty),t,x)\}_{t \ge 0, x \in D}$ has the same covariance, and hence the same law, as $\{ u(t,x) \}_{t \ge 0, x \in D}$.
Therefore, we may and will identify $u([0,\infty),\cdot) = u(\cdot)$ for the rest of the paper.
Following \cite{Balan12,DMX17,LX23}, we call \eqref{eq:harm} the harmonizable representation of $u$.

%\begin{lemma}[Harmonizable representation]\label{lem:harm}
%For any \(A\in\mathcal B([0,\infty))\), define
%\begin{equation}\label{eq:harm}
%    u(A,t,x)
%    =\frac1{\sqrt{2\pi}}
%    \sum_{n=1}^{\infty}\lambda_n^{-\alpha/2}f_n(x)
%    \Re\int_{\{\tau:\kappa_n(\tau)\in A\}}
%    q_{n,t}(\tau)\,W_n(\d\tau).
%\end{equation}
%Then, under \eqref{Dalang}, the following statements hold.
%\begin{enumerate}
%\item The series in \eqref{eq:harm} converges in
%\(L^2(\Omega)\) for every \(A,t,x\).
%\item For every \((t,x)\), the map \(A\mapsto u(A,t,x)\) is independently
%scattered. Moreover, the random fields \(u(A,\cdot)\) and \(u(B,\cdot)\) are
%independent whenever $A$ and $B$ are disjoint.
%\item The centered Gaussian random field
%\(\{u([0,\infty),t,x)\}_{t \ge 0, x \in D}\) has the same law as \(\{u(t,x)\}_{t \ge 0, x\in D}\).
%\end{enumerate}
%The field \(v\) has a separable jointly continuous version on
%\([0,T]\times\overline D\) for every \(T>0\). We henceforth realize \(u\) as
%this version of \(v\), and write \(u(A,t,x)=v(A,t,x)\).
%\end{lemma}

%\begin{proof}
%The first and third statements follow from Proposition \ref{prop:spec:rep}, Plancherel's theorem, and the identity
%\[
%    \int_0^t \e^{-\lambda(t-r)}\e^{-i\tau r}\,\d r
%    =
%    \frac{e^{-i\tau t}-e^{-\lambda t}}{\lambda-i\tau}.
%\]
%The second statement is clear from construction.
%By the white-noise isometry, restricting the frequency integral to
%\(\{\kappa_n\in A\}\) can only decrease the variance.
%For \(A=[0,\infty)\), Plancherel's theorem identifies the variance of
%\(v_M-v_N\) with the corresponding spectral tail in
%Theorem~\ref{thm:exist}. This tail converges to zero by
%Lemma~\ref{lem:Dalang}, proving the first statement.
%
%If \(A\) and \(B\) are disjoint, then, for every \(n\), the sets
%\(\{\tau:\kappa_n(\tau)\in A\}\) and
%\(\{\tau:\kappa_n(\tau)\in B\}\) are disjoint. The independently scattered
%property of each \(W_n\), together with independence over \(n\), proves the
%second statement.
%It remains to compare the laws of \(v\) and \(u\). 
%Let \(v_N(A,t,x)\) denote the sum of the first \(N\) modes in
%\eqref{eq:harm}, and write
%\[
%    v_N(t,x)=v_N([0,\infty),t,x).
%\]
%For each $N$,
%\begin{align*}
%    \E[v_N(t,x)v_N(s,y)]
%    &=
%    \frac1{2\pi}
%    \sum_{n=1}^N
%    \lambda_n^{-\alpha}f_n(x)f_n(y)                                     
%    \int_\R
%    \frac{(\e^{-i\tau t}-\e^{-\lambda_n t})
%    (\e^{i\tau s}-\e^{-\lambda_n s})}
%    {\lambda_n^2+\tau^2}\,\d\tau .
%\end{align*}
%The elementary Fourier identity
%\[
%    \int_0^t \e^{-\lambda(t-r)}\e^{-i\tau r}\,\d r
%    =
%    \frac{e^{-i\tau t}-e^{-\lambda t}}{\lambda-i\tau}
%\]
%and Plancherel's theorem give
%\[
%    \int_\R
%    \frac{(\e^{-i\tau t}-\e^{-\lambda_n t})
%    (\e^{i\tau s}-\e^{-\lambda_n s})}
%    {\lambda_n^2+\tau^2}\,\d\tau
%    =
%    2\pi\int_0^{t\wedge s}
%    \e^{-\lambda_n(t-r)}\e^{-\lambda_n(s-r)}\,\d r .
%\]
%Passing \(N\to\infty\) in \(L^2(\Omega)\) yields \eqref{eq:cov}.
%Therefore \(v\) and \(u\) have the same finite-dimensional distributions.
%\end{proof}


\begin{proposition}\label{prop:LX}
Suppose \(D\) is a bounded \(C^2\) domain and
\eqref{alpha:Holder:range} holds.
Let \(R=R_1 \times R_2\) be a
compact rectangle in $(0,\infty)\times D$. Write \(z=(z_0,\ldots,z_d)=(t,x_1,\ldots,x_d)\)
and set
\[
    \alpha_0=\theta/2,\qquad \alpha_i=\theta\ \, (1\le i\le d),
    \qquad
    \gamma_i=\alpha_i^{-1}-1\ \, (0\le i \le d).
\]
Then every \(\gamma_i\) is positive, and the Gaussian field \(u(A,z)\) from
\eqref{eq:harm} satisfies the following properties.

\begin{enumerate}
\item For every \(z\in R\), the map \(A\mapsto u(A,z)\) is an
independently scattered Gaussian noise,
\(u([0,\infty),z)=u(z)\), and the fields \(u(A,\cdot)\) and
\(u(B,\cdot)\) are independent for disjoint \(A,B\). 
Thus Assumption~2.1(a) of \cite{LX23} holds.\smallskip

\item There exists $C_R>0$ such that for every \(1\le a<b\le\infty\) and \(z,z'\in R\),
\begin{gather*}
\qquad\qquad \|u([a,b),z)-u(z)-u([a,b),z')+u(z')\|_2\le C_R\left[
   \sum_{i=0}^d a^{\gamma_i}|z_i-z'_i|+b^{-1}\right],\\
	\|u([0,1),z) - u([0,1),z')\|_2 \le C_R \sum_{i=0}^d |z_i-z_i'|,
\end{gather*}
with \(0^{\gamma_i}=0\) and \(\infty^{-1}=0\). So, Assumption~2.1(b) of \cite{LX23} holds with \(a_0=1\).\smallskip

\item Let \(z^0=(0,\dots, 0)\in\R^{d+1}\). There exists \(c_0>0\) such that
\[
    \Var\left(u(z)\mid u(z^1),\ldots,u(z^n)\right)
    \ge
    c_0\min_{0\le j\le n}\rho(z,z^j)^2
\]
for all \(n\ge1\) and \(z,z^1,\ldots,z^n\in R\). Thus
Assumption~2.2 of \cite{LX23} holds.\smallskip

\item There exist $c_1,c_2>0$ such that the metric $d_u(z,z') := \|u(z)-u(z')\|_2$ satisfies
\[
    c_1\rho(z,z')\le d_u(z,z')\le c_2\rho(z,z'),
    \qquad \forall z,z'\in R,
\]
so Assumption~2.3 of \cite{LX23} holds.
\end{enumerate}
%
%Moreover, \(u\) is continuous and separable on \(R\), and
%\[
%    N_\Delta(R,r)\le C_Rr^{-Q},\qquad 0<r<1.
%\]
%Consequently, Theorems~4.4, 5.2, and 6.1 of  \cite{LX23} apply on
%\(R\), with metric \(\rho\) and
%\(Q=\sum_{i=0}^d\alpha_i^{-1}\).
\end{proposition}

\begin{proof}
First, under condition \eqref{alpha:Holder:range}, \(0<\theta<1\), so
\(\gamma_0=(\theta/2)^{-1}-1>0\) and
\(\gamma_i=\theta^{-1}-1>0\) for \(1\le i\le d\).
Item (1) follows from the definition \eqref{eq:harm} and the fact that $W_n$ are independent Gaussian noises. 

To prove item (2), we claim that there exists $C>0$ such that
\begin{align*}
	&\|u([a,b),t,x)-u(t,x)-u([a,b),s,y)+u(s,y)\|_{2}\\
	&\hskip1.5in \le C\bigl(a^{\gamma_0}|t-s|+a^{\gamma_1}|x-y|+b^{-1}\bigr),\\
	& \|u([0,1),t,x) - u([0,1),s,y)\|_2 \le C (|t-s|+|x-y|)
\end{align*}
for all $1\le a <b \le \infty$ and $(t,x),(s,y) \in R$.

The following identities will be used throughout the proof:
\[
    \frac{2}{\theta}(1+d/2-\alpha)=2\gamma_0,
    \qquad
    \frac{2}{\theta}(d/2-\alpha)=2\gamma_1,
    \qquad
    \frac{2}{\theta}(d/2-\alpha-1)=-2 .
\]
Since
\[
\begin{aligned}
&u([a,b),t,x)-u(t,x)-u([a,b),s,y)+u(s,y)  \\
&\quad =
    -\bigl(u([0,a),t,x)-u([0,a),s,y)\bigr)
    -\bigl(u([b,\infty),t,x)-u([b,\infty),s,y)\bigr),
\end{aligned}
\]
its $\|\cdot\|_2$ norm is bounded by
\[
    I_a+J_b(t,x)+J_b(s,y),
\]
where
\[
    I_a=\|u([0,a),t,x)-u([0,a),s,y)\|_2,
    \qquad
    J_b(t,x)=\|u([b,\infty),t,x)\|_2 .
\]

We first estimate $J_b(t,x)$. If \(b=\infty\), then \(J_b=0\).
%If \(0<b\le1\), then \(J_b(t,x)\le C\le Cb^{-1}\), using boundedness of $\Var(u(t,x))$ from Lemma \ref{lem:var:u:1}. 
Let \(1<b<\infty\) and set
\[
	\Lambda_b=b^{2/\theta}.
\] 
By Wiener isometry and
\(|\e^{-i\tau r}-\e^{-\lambda_n r}|\le2\),
\[
    |J_b(t,x)|^2
    \lesssim
    \sum_{n=1}^\infty
    \lambda_n^{-\alpha}|f_n(x)|^2
    \int_{\{\tau:\lambda_n^{\theta/2}\vee|\tau|^{\theta/2}\ge b\}}
    \frac{\d\tau}{\lambda_n^2+\tau^2}.
\]
If \(\lambda_n<\Lambda_b\), then the domain of integration is contained in
\(\{\tau: |\tau|\ge\Lambda_b\}\), and
\[
\begin{aligned}
    \sum_{n \ge 1:\lambda_n<\Lambda_b}
    \lambda_n^{-\alpha}|f_n(x)|^2
    \int_{|\tau|\ge\Lambda_b}
    \frac{\d\tau}{\lambda_n^2+\tau^2}
    &\lesssim
    \Lambda_b^{-1}
    \sum_{n \ge 1: \lambda_n<\Lambda_b}\lambda_n^{-\alpha}|f_n(x)|^2  \\
    &\lesssim
    \Lambda_b^{d/2-\alpha-1}
    =
    b^{-2},
\end{aligned}
\]
thanks to Lemma~\ref{lem:spectral:1}. If \(\lambda_n\ge\Lambda_b\), then the domain of integration is $\R$ and
\[
    \int_\R\frac{\d\tau}{\lambda_n^2+\tau^2}=\frac{\pi}{\lambda_n}.
\]
Thus by Lemma \ref{lem:spectral:2},
\begin{align*}
    \sum_{n \ge 1: \lambda_n\ge\Lambda_b}
    \lambda_n^{-\alpha}|f_n(x)|^2
    \int_\R\frac{\d\tau}{\lambda_n^2+\tau^2}
    &\lesssim
    \sum_{n \ge 1:\lambda_n\ge\Lambda_b}
    \lambda_n^{-(\alpha+1)}|f_n(x)|^2\\
    &\lesssim
    \Lambda_b^{d/2-\alpha-1}
    =
    b^{-2}.
\end{align*}
Consequently, \(J_b(t,x)\lesssim b^{-1}\) for all $1<b\le \infty$ and $(t,x) \in R$.

%We estimate $I_a$ next. 
Next, we prove that $I_a \lesssim a^{\gamma_0}|t-s| + a^{\gamma_1} |x-y|$ for all $a \ge 1$ and $(t,x),(s,y) \in R$.
Note that it is enough to prove this for all $|t-s| \le \varepsilon_0$ and $|x-y| \le \varepsilon_0$, where $\varepsilon_0>0$ is a fixed number, because $\|u([0,a),t,x)\|_2 \le \|u(t,x)\|_2$ and Lemma  \ref{lem:var:u:1} implies that $\|u(t,x)\|_2$ is uniformly bounded.
%If \(a=0\), then \(I_a=0\). Assume \(a>0\) and 
Let $a \ge 1$ and set 
\[
	\Lambda_a=a^{2/\theta}.
\]
Recall \eqref{q}. 
By Wiener isometry,
\[
    I_a^2
    \lesssim
    \sum_{n \ge 1: \lambda_n<\Lambda_a}
    \lambda_n^{-\alpha}
    \int_{|\tau|<\Lambda_a}
    |f_n(x)q_{n,t}(\tau)-f_n(y)q_{n,s}(\tau)|^2\,\d\tau .
\]
Split the last expression as \(A_1+A_2\), where
\[
    A_1:=
    \sum_{n \ge 1: \lambda_n<\Lambda_a}
    \lambda_n^{-\alpha}|f_n(x)|^2
    \int_{|\tau|<\Lambda_a}|q_{n,t}(\tau)-q_{n,s}(\tau)|^2\,\d\tau,
\]
and
\[
    A_2:=
    \sum_{n \ge 1: \lambda_n<\Lambda_a}
    \lambda_n^{-\alpha}|f_n(x)-f_n(y)|^2
    \int_{|\tau|<\Lambda_a}|q_{n,s}(\tau)|^2\,\d\tau .
\]
Since $|\e^{-i\tau t}-\e^{-i\tau s}|\le|\tau||t-s|$ and $|\e^{-\lambda_n t}-\e^{-\lambda_n s}|\le\lambda_n|t-s|$,
we have 
\[
	|q_{n,t}(\tau)-q_{n,s}(\tau)|\lesssim |t-s|.
\]
Hence, by Lemma \ref{lem:spectral:1},
\[
    A_1
    \lesssim
    |t-s|^2\Lambda_a
    \sum_{n \ge 1: \lambda_n<\Lambda_a}\lambda_n^{-\alpha}|f_n(x)|^2
    \lesssim
    |t-s|^2\Lambda_a^{1+d/2-\alpha}
    =
    a^{2\gamma_0}|t-s|^2 .
\]
Also, we have
\[
	|q_{n,s}(\tau)|\le \frac2{(\lambda_n^2+\tau^2)^{1/2}}
	\quad \text{and}\quad 
	\int_\R |q_{n,s}(\tau)|^2\,\d\tau\lesssim\lambda_n^{-1}.
\]
Thus,
\[
    A_2
    \lesssim
    \sum_{n \ge 1: \lambda_n<\Lambda_a}
    \lambda_n^{-(\alpha+1)}|f_n(x)-f_n(y)|^2 .
\]
Let $\varepsilon_0 \in (0,1)$ and $C_0>0$ be the constants given by Lemma \ref{lem:curve}.
%Since bounded \(C^2\)-domains are quasiconvex, \(x\) and \(y\) can be joined
%by a rectifiable curve \(\gamma:[0,\ell]\to D\), parametrized by arclength,
%with \(\ell\le C|x-y|\); see Brudnyi and Brudnyi
%\cite[Definition~3.38 and Example~3.39(a)]{BB12}. Therefore
By that lemma, for any $x,y \in R_2$ with $|x-y|\le\varepsilon_0$, we can find a $C^2$ curve $\gamma:[0,1] \to D$ such that $\gamma(0)=x$, $\gamma(1) = y$ and 
\[
	|\gamma'(r)| \le C_0|x-y| \quad \text{for all $r \in [0,1]$.}
\]
Then, by the fundamental theorem of calculus and Cauchy-Schwarz inequality,
\begin{align*}
    |f_n(x)-f_n(y)|^2 
    &=
    \left|\int_0^1 \nabla f_n(\gamma(r))\cdot\gamma'(r)\,\d r\right|^2
    \le C_0^2 |x-y|^2 \int_0^1 |\nabla f_n(\gamma(r))|^2\,\d r.
\end{align*}
This together with Lemma~\ref{lem:spectral:3} yields
\begin{align*}
    A_2
    &\lesssim
    |x-y|^2 \int_0^1
    \sum_{n \ge 1: \lambda_n<\Lambda_a}
    \lambda_n^{-(\alpha+1)}|\nabla f_n(\gamma(r))|^2\,\d r\\
    &\lesssim
    |x-y|^2\Lambda_a^{d/2-\alpha}
    =
    a^{2\gamma_1}|x-y|^2 .
\end{align*}
Combining \(A_1\) and \(A_2\), we have
\[
    I_a
    \lesssim
    a^{\gamma_0}|t-s|+a^{\gamma_1}|x-y|.
\]
Since \(|x-y|\asymp \sum_{i=1}^d|x_i-y_i|\), the last display and the
previous estimates for $J_b$ give exactly the first bound in item (2). 
The estimate for $I_a$ above with $a=1$ gives the second bound.
%The choice \(a_0=0\) is valid because \(u([0,0),\cdot)=0\). 
This proves item (2).

%Next, we verify item (3). Set
%\[
%    K_R=\inf_{(t,x)\in R}\bigl(t^{\theta}\wedge d(x)^{2\theta}\bigr)>0
%\]
%and
%\[
%    M_R=1\vee\sup_{z,z'\in R}\rho^2(z,z')
%          \vee\sup_{z\in R}\rho^2(z,0)<\infty .
%\]
%By Proposition~\ref{prop:SLND}, there exists $c_R>0$ such that
%\[
%\begin{aligned}
%\Var\left(u(z)\mid u(z^1),\ldots,u(z^n)\right)
%&\ge
%    c_R\left(\min_{1\le j\le n}\rho(z,z^j)^2\wedge K_R\right)\\
%&\ge
%    c_R'\min_{0\le j\le n}\rho(\xi,\xi^j)^2,
%    \qquad \xi^0=0,
%\end{aligned}
%\]
%where the last inequality follows after multiplying the constant by
%\(\min\{1,b_R/M_R\}\). 
%This is precisely Assumption~2.2 of 
%\cite{LX23}.
Item (3) follows from Proposition \ref{prop:SLND} and
$\inf_{(t,x) \in R}[t^\theta \wedge (d(x))^{2\theta}] > 0$ because the compact rectangle $R \subset (0,\infty) \times D$ has a strictly positive distance from $\{0\} \times D$ and $[0,T] \times \partial D$.
Finally, item (4) follows from \eqref{var:u-u}.
%Proposition~\ref{prop:metric} 
%and
%\(\rho\asymp\Delta\). 
%Finally,
%Proposition~\ref{prop:continuous} gives continuity and separability,
%and Lemma~\ref{lem:entropy} gives the entropy bound.
\end{proof}

\subsection{Proof of Theorem~\ref{thm:moc}}

\begin{proof}
Fix $z_0 \in (0,\infty) \times D$.
Choose a compact rectangle \(R \subset (0,\infty) \times D\) with \(z_0\) in its interior. Proposition~\ref{prop:LX}
verifies Assumptions~2.1 and 2.3 of \cite{LX23}. Therefore, Theorem~5.2 of 
\cite{LX23} applies to $\{u(t,x)\}_{(t,x) \in R}$ and gives the asserted exact local
modulus with respect to the metric \(\rho\), with a constant $K_0$ that is finite and strictly positive.
\end{proof}

\subsection{Proof of Theorem~\ref{thm:umoc}}

We first recall the notion of metric entropy.
Let $d_u(z,z') = \|u(z)-u(z')\|_2$ denote the canonical metric for $u$ on $[0,\infty) \times D$.
For any set $A \subset [0,\infty)\times D$, let $N(A,r)$ be the entropy number, defined as the smallest number of $d_u$-balls of radius $r$ needed to cover $A$.

\begin{lemma}\label{lem:N}
Suppose $D\subset \R^d$ is a bounded $C^2$ domain and \eqref{alpha:Holder:range} holds.
Fix $T>0$. Then, there exists $C>0$ such that
\begin{align*}
	N([0,T]\times D, r) \le Cr^{-Q} \quad \forall r \in (0,1],
\end{align*}
where $Q$ is given by \eqref{eq:Q}.
\end{lemma}

\begin{proof}
For any $r>0$, define
\[
	D_{r} := \{ x \in D: d(x) > r^{1/\theta} \},
\]
where $\theta$ and $d(x)$ are given by \eqref{theta} and \eqref{d(x)}, respectively.
By Theorem \ref{thm:optimal}, there exist $c_0,c_1>0$ such that
\begin{align}\begin{split}\label{d_u:bd}
	&c_0 \left[ \rho((t,x),(s,y)) \wedge \left( (\sqrt{t} \wedge d(x))^\theta \vee (\sqrt{s} \wedge d(y))^\theta \right) \right]\\
	&\le d_u((t,x),(s,y)) \le c_1 \left[ \rho((t,x),(s,y)) \wedge  \left( (\sqrt{t} \wedge d(x))^\theta \vee (\sqrt{s} \wedge d(y))^\theta \right)  \right]
\end{split}\end{align}
for all $(t,x), (s,y) \in [0,T] \times D$.
For any $r \in (0,r_0)$, we write $[0,T]\times D = E_r \cup F_r$, where
\[
	E_r = [(r/c_1)^{2/\theta}, T] \times D_{r/c_1} \quad \text{and} \quad 
	F_r = ([0,T]\times D) \setminus E_r.
\]
By \eqref{d_u:bd}, $F_r$ can be covered by a single $d_u$-ball of radius $r$, thus \[
	N(F_r, r) = 1.
\]
To estimate $N(E_r,r)$, let $B_1,\dots, B_p$ be a maximal set of pairwise disjoint $d_u$-balls of radius $r/2$ with centers in $E_r$. 
Maximality implies that the balls $2B_1,\dots, 2B_p$ form a cover for $E_r$, where $2B_i$ is the $d_u$-ball with radius $r$ and the same center as $B_i$.
Hence, $N(E_r,r) \le p$.
Since $D$ is a bounded domain and $B_1,\dots, B_p$ are pairwise disjoint, their total volume $\sum_{i=1}^p \mathrm{Vol}(B_i)$ is uniformly bounded.
By \eqref{d_u:bd}, $N(E_r,r) r^Q \le p r^Q\lesssim \sum_{i=1}^p \mathrm{Vol}(B_i) = O(1)$, hence we can find $C_0>0$ such that
\[
	N(E_r,r) \le p \le C_0 r^{-Q}.
\]
Therefore, we have $N([0,T]\times D, r) \le 1 + C_0 r^{-Q} \lesssim r^{-Q}$ for all $r \in (0,1]$.
\end{proof}

\begin{proof}[Proof of Theorem~\ref{thm:umoc}]
%Define $d_u(z,z') = \|u(z)-u(z')\|_2$.
First, by \eqref{var:u-u}, we have
\[
	\lim_{\varepsilon\to0^+} \sup_{\substack{z,z'\in [0,T]\times D\\ 0<\rho(z,z') \le \varepsilon}} \frac{d_u(z,z')}{\rho(z,z') \sqrt{\log(1/\rho(z,z'))}} = 0.
\]
Hence, we may use Lemma 7.1.1 of \cite{MR} to deduce that \eqref{umoc} holds for some constant $K_1 \in [0,\infty]$.
It remains to show that $0<K_1<\infty$.

To show that $K_1<\infty$, we apply Theorem 1.3.5 of \cite{AT} to see that there exist a universal constant $K>0$ and a random variable $\eta>0$ such that a.s., for all $\varepsilon \in (0,\eta]$,
\begin{align*}
	\sup_{\substack{z,z'\in [0,T]\times D\\\rho(z,z') \le \varepsilon}} |u(z)-u(z')| \le K \int_0^\varepsilon \sqrt{\log N([0,T]\times D, r)} \, \d r.
\end{align*}
By Lemma \ref{lem:N}, $N([0,T]\times D, r) \le C r^{-Q}$ for all $r \in (0,1]$.
Hence, we may let $\varepsilon_n = e^{-n}$ and proceed as in the proof of Theorem 3.13 of \cite{HL26} to deduce that there exist constants $C_1, C_2 \in (0,\infty)$ such that a.s., for all $n$ large,
\[
	\sup_{\substack{z,z'\in [0,T]\times D\\\rho(z,z') \le \varepsilon_n}} |u(z)-u(z')| 
	\le C_1 \varepsilon_n \sqrt{\log(1/\varepsilon_n)}
\]
and
\[
	\lim_{n \to \infty} \sup_{\substack{z,z'\in [0,T]\times D\\\rho(z,z') \le \varepsilon_n}} \frac{|u(z)-u(z')|}{\rho(z,z')\sqrt{\log(1/\rho(z,z'))}} \le C_2 \quad \text{a.s.}
\]
This implies that $K_1 \le C_2 <\infty$.

To show that $K_1>0$, we choose and fix any compact rectangle $R \subset (0,T] \times D$.
By Proposition \ref{prop:LX}, $\{u(t,x)\}_{(t,x) \in R}$ satisfies Assumptions~2.1--2.3 of \cite{LX23}.
Hence, Theorem~6.1 of \cite{LX23} implies that
\[
    \lim_{\varepsilon\to0^+}
    \sup_{\substack{z,z'\in R\\0<\rho(z,z')\le\varepsilon}}
    \frac{|u(z)-u(z')|}
    {\rho(z,z')\sqrt{\log(1/\rho(z,z'))}}
    =
    C
    \qquad\text{a.s.}
\]
for some constant $C \in (0,\infty)$.
Clearly, we have $K_1 \ge C>0$.
This completes the proof of Theorem \ref{thm:umoc}.
% For the compact rectangle \(R \subset (0,\infty) \times D\) in the statement,
% Under \eqref{alpha:Holder:range}, Proposition~\ref{prop:LX} verifies 
% Assumptions~2.1--2.3 of \cite{LX23} with
% \(\alpha_0=\theta/2\) and \(\alpha_1=\cdots=\alpha_d=\theta\). 
% Hence, Theorem~6.1 of \cite{LX23} applies directly and gives the displayed exact uniform modulus of continuity.
\end{proof}

\subsection{Proof of Theorem~\ref{thm:lil}}

\begin{proof}
Fix \(z_0\in(0,\infty)\times D\) and choose a compact rectangle
\(R \subset (0,\infty) \times D\) with \(z_0\) in its interior. 
Under \eqref{alpha:Holder:range}, Proposition~\ref{prop:LX} shows that $\{u(t,x)\}_{(t,x) \in R}$ satisfies Assumptions~2.1 and 2.2 of 
\cite{LX23} with \(\alpha_0=\theta/2\) and
\(\alpha_i=\theta\) for \(1\le i\le d\).
Therefore, we may apply Theorem~4.4 of \cite{LX23} to obtain the asserted Chung's law of the iterated logarithm with a  constant in \((0,\infty)\). 
\end{proof}

\subsection{Proof of Theorem~\ref{thm:sbp}}

%The proof below uses the Talagrand entropy lemma \cite[Lemma~4.1]{LX23} and the one-dimensional Gaussian unimodality form of Anderson's inequality \cite{Anderson}. 
%Measurability of all suprema follows from Proposition~\ref{prop:continuous}.

%\begin{lemma}\label{lem:entropy}
%Let \(Q\) be defined by \eqref{eq:Q}. For every \(T>0\) and
%\(\delta>0\),
%\[
%    N_\rho([0,T]\times D_\delta,r)\le Cr^{-Q},
%    \qquad 0<r<1.
%\]
%The constant may depend on \(T\), \(D\), and \(\delta\); only the exponent
%\(Q\) is used in the regularity and small ball arguments.
%More generally, if \(B\subset\R^d\) is bounded and \(0<a<T<\infty\), then
%\[
%    N_\rho([a,T]\times B,r)
%    \le C\bigl(1+(T-a)r^{-1/\theta}\bigr)
%       \bigl(1+N_E(B,c r^{1/(2\theta)})\bigr),
%    \qquad 0<r<1,
%\]
%where \(N_E(B,s)\) denotes the covering number of \(B\) by Euclidean
%balls of radius \(s\). Placing \(B\) in a fixed bounding box gives
%\(N_\rho([a,T]\times B,r)\le C_{a,T,B}r^{-Q}\).
%Moreover, for every fixed \(z_0\) and \(0<r\le R\le1\),
%\[
%    N_\rho(B_\rho(z_0,R),r)\le C\left(\frac Rr\right)^Q .
%\]
%\end{lemma}
%
%\begin{proof}
%Cover time by intervals of length \(r^{1/\theta}\) and space by Euclidean
%balls of radius \(r^{1/(2\theta)}\). This requires at most
%\[
%    Cr^{-1/\theta}r^{-d/(2\theta)}=Cr^{-Q}
%\]
%sets. The local estimate is the same argument applied to a time interval of
%length \(R^{1/\theta}\) and a spatial ball of radius \(R^{1/(2\theta)}\).
%The same covering argument applies to any bounded spatial set \(B\), after
%placing \(B\) in a fixed Euclidean box.
%\end{proof}


\begin{proof}[Proof of Theorem \ref{thm:sbp}]
By Lemma 2.2 of Talagrand \cite{Talagrand} (see also \cite[Lemma 3.4]{DLMX} for a more precise statement) and Lemma \ref{lem:N}, there exists $K>0$ such that 
\[
	\P\left\{ \sup_{z,z' \in [0,T]\times D}|u(z)-u(z')| \le \varepsilon \right\} \ge e^{-K/\varepsilon^Q}\qquad
	\forall \varepsilon \in (0,1].
\]
Fix $t_0 \in (0,T]$ and $x_0 \in D$.
Thanks to the Gaussian correlation inequality and the preceding, we can find $K_3>0$ such that
\begin{align*}
	\P\left\{ \sup_{z \in [0,T]\times D}|u(z)| \le \varepsilon \right\}
	&\ge \P\left\{ |u(t_0,x_0)| \le \frac{\varepsilon}{2} \right\} \, \P\left\{ \sup_{z,z' \in [0,T]\times D}|u(z)-u(z')| \le \frac{\varepsilon}{2} \right\} \\
	& \ge \P\left\{ |Z| \le \frac{\varepsilon}{2\|u(t_0,x_0)\|_2} \right\} e^{-K2^Q/\varepsilon^Q}
	\ge e^{-K_3/\varepsilon^Q}
\end{align*}
for all $\varepsilon \in (0,1]$, where $Z$ has a standard normal distribution, and we have used $\Var(u(t,x)) \asymp t^{\theta} \wedge (d(x))^{2\theta}$, which follows from  \eqref{var:u-u}, and $\P\{|Z| \le a \} \ge C_A a$ for $|a| \le A$ to obtain the last inequality. This proves the lower bound in Theorem \ref{thm:sbp}.

To show the upper bound, we fix a compact rectangle $R \subset (0,T] \times D$.
By Theorem \ref{thm:optimal}, $\|u(z)-u(z')\|_2 \asymp \rho(z,z')$ uniformly for all $z,z' \in R$ and $\{u(z)\}_{z \in R}$ satisfies the SLND property.
Hence, we may apply Theorem 5.1 of \cite{X09} to find $K_4>0$ such that
\[
	\P\left\{ \sup_{z \in R} |u(z)| \le \varepsilon \right\} \le e^{-K_4 \varepsilon^{-Q}}
\]
uniformly for all $\varepsilon \in (0,1]$.
This implies the upper bound in Theorem \ref{thm:sbp}.
%
%
%Let \(I=[a,T]\times A\). The global upper canonical metric bound gives
%\[
%    d_u(z,z')\le C\rho(z,z'),\qquad z,z'\in I.
%\]
%Place \(A\) in a bounding box containing \(D\); the entropy estimate
%Lemma~\ref{lem:entropy} gives
%\[
%    N_{d_u}(I,r)
%    \le C_{a,T,A}r^{-Q},
%    \qquad 0<r<1.
%\]
%
%\noindent\textit{Lower bound.}
%Augment the parameter set by a fixed reference point \(*\) and define a
%centered Gaussian process \(X\) by \(X(*)=0\) and \(X(z)=u(z)\) for
%\(z\in I\). The augmented canonical metric satisfies
%\(N(S,d_X,r)\le 1+N_{d_u}(I,r)\le C_{a,T,A}r^{-Q}\) for
%\(S=I\cup\{*\}\). Applying Talagrand's entropy lower bound for Gaussian
%processes, in the form of \cite[Lemma~4.1]{LX23} with
%\(\psi(r)=C_{a,T,A}r^{-Q}\), yields
%\[
%    \P\left\{\sup_{s,t\in S}|X(s)-X(t)|\le\varepsilon\right\}
%    \ge
%    \exp\left(-c_1\varepsilon^{-Q}\right).
%\]
%Since the event on the left is contained in
%\(\{\sup_{z\in I}|u(z)|\le\varepsilon\}\), the lower bound follows.
%
%\noindent\textit{Upper bound.}
%By inner regularity of Lebesgue measure, choose a compact set \(K\subset A\)
%such that \(|K|\ge |A|/2\). Since \(K\) is compactly contained in \(D\),
%there exists \(\delta_K>0\) such that \(K\subset D_{\delta_K}\).
%Fix small constants \(b_t,b_x>0\). Choose a time grid in \([a,T]\) with
%spacing \(b_t\varepsilon^{1/\theta}\), and a maximal Euclidean
%\(b_x\varepsilon^{1/(2\theta)}\)-separated set in \(K\). For all sufficiently
%small \(\varepsilon\), the time grid has at least
%\(c(T-a)\varepsilon^{-1/\theta}\) points. Maximality of the spatial set makes
%the balls of radius \(b_x\varepsilon^{1/(2\theta)}\) cover \(K\), so volume
%comparison gives at least
%\(c|K|\varepsilon^{-d/(2\theta)}\) spatial points. Their Cartesian product is
%a \(\rho\)-packing
%\(\{z_1,\ldots,z_N\}\subset [a,T]\times K\) with mutual distances at least
%\(c\varepsilon\) and
%\[
%    N\ge c_K\varepsilon^{-1/\theta-d/(2\theta)}
%      =c_K\varepsilon^{-Q}.
%\]
%
%Order the points arbitrarily. Continuity and positivity of the variance give
%\[
%    v_K:=\inf_{z\in[a,T]\times K}\Var(u(z))>0.
%\]
%Thus the first point, whose conditional variance is its ordinary variance, is
%handled separately. If \(\varepsilon\le\sqrt{v_K}\), then for a standard
%normal random variable \(Z\),
%\[
%    \P\{|u(z_1)|\le\varepsilon\}
%    \le \P\{|Z|\le1\}=:p_0<1.
%\]
%For \(k\ge2\), Proposition~\ref{prop:SLND} and the
%packing condition give a uniform conditional variance lower bound
%\[
%    \Var(u(z_k)\mid u(z_1),\ldots,u(z_{k-1}))
%    \ge c\varepsilon^2,
%\]
%since \(t_k\ge a\) and \(f_1(x_k)^{4\theta}\ge c_K>0\) on \(K\).
%
%Let \(\mathcal F_{k-1}=\sigma(u(z_1),\ldots,u(z_{k-1}))\). Conditionally on
%\(\mathcal F_{k-1}\), \(u(z_k)\) is Gaussian with mean
%\(\mu_k=\E[u(z_k)\mid u(z_1),\ldots,u(z_{k-1})]\) and variance
%\(\sigma_k^2\ge c\varepsilon^2\). By Anderson's inequality (or unimodality),
%for a standard normal \(Z\),
%\[
%\begin{aligned}
%    \P\{|u(z_k)|\le\varepsilon\mid\mathcal F_{k-1}\}
%    &\le \sup_{\mu\in\R}\P\{|\mu+\sigma_kZ|\le\varepsilon\}\\
%    &=\P\{|Z|\le\varepsilon/\sigma_k\}
%      \le\P\{|Z|\le C\}=:p<1,
%\end{aligned}
%\]
%uniformly in \(k\). Iterating these conditional probabilities yields
%\[
%    \P\Bigl\{\max_{1\le k\le N}|u(z_k)|\le\varepsilon\Bigr\}
%    \le p_0p^{N-1}
%    \le p_*^N
%    \le
%    \exp(-c_2\varepsilon^{-Q}).
%\]
%Here \(p_*=\max\{p_0,p\}<1\).
%Since
%\[
%    \bigl\{\sup\nolimits_{z\in I}|u(z)|\le\varepsilon\bigr\}
%    \subset
%    \bigl\{\max\nolimits_{1\le k\le N}|u(z_k)|\le\varepsilon\bigr\},
%\]
%the upper bound follows. Choose \(\varepsilon_0\) as the minimum of
%the small-radius thresholds in the entropy lemma, the two grid counts, the
%standard-deviation scales \(a^\theta\) and
%\(\inf_K f_1^{2\theta}\) from the SLND caps, and the first-point scale
%\(\sqrt{v_K}\), with the fixed packing constants absorbed into the
%threshold. This completes the proof.
\end{proof}

{\bf Acknowledgments.}
C.Y. Lee was supported in part by the Shenzhen Peacock grant
2025TC0013.

%\appendix
%\section{Auxiliary estimates}

%\begin{lemma}[Heat kernel estimates]\label{lem:heat}
%Let \(D\subset\R^d\) be a bounded \(C^2\)-domain. For every \(T>0\) there are
%constants \(C,c>0\) such that for all \(t\in(0,T]\) and \(x,y\in D\),
%\[
%    0\le P_t(x,y)\le
%    Ct^{-d/2}\exp\left(-\frac{|x-y|^2}{ct}\right),
%\]
%\[
%    P_t(x,y)\le
%    C\left(1\wedge\frac{d(x)}{\sqrt t}\right)
%    \left(1\wedge\frac{d(y)}{\sqrt t}\right)
%    t^{-d/2}\exp\left(-\frac{|x-y|^2}{ct}\right),
%\]
%and
%\[
%    |\nabla_xP_t(x,y)|
%    \le
%    Ct^{-(d+1)/2}\exp\left(-\frac{|x-y|^2}{ct}\right).
%\]
%For large \(t\), the same estimates hold after multiplying the right-hand
%side by \(\e^{-\lambda_1t/2}\), with adjusted constants.
%\end{lemma}
%
%\begin{proof}
%By the spectral theorem,
%\[
%    P_t(x,y)=\sum_{n=1}^\infty \e^{-\lambda_n t}f_n(x)f_n(y),
%\]
%which follows from the spectral theorem for \(A=-\Delta_D\). It gives
%symmetry \(P_t(x,y)=P_t(y,x)\) and, by orthonormality of the eigenfunctions,
%the semigroup identity
%\[
%    \int_D P_t(x,z)P_s(z,y)\,\d z=P_{t+s}(x,y).
%\]
%The maximum principle gives \(P_t(x,y)\ge0\).
%
%For \(0<t\le T\), domination by the Euclidean heat kernel gives the Gaussian
%upper bound
%\[
%    P_t(x,y)
%    \le
%    Ct^{-d/2}\exp\left(-\frac{|x-y|^2}{ct}\right),
%\]
%see also Davies \cite{Davies}. The boundary-sensitive version is
%\[
%    P_t(x,y)
%    \le
%    C
%    \left(1\wedge\frac{d(x)}{\sqrt t}\right)
%    \left(1\wedge\frac{d(y)}{\sqrt t}\right)
%    t^{-d/2}
%    \exp\left(-\frac{|x-y|^2}{ct}\right),
%\]
%which is the upper half of the sharp Dirichlet heat-kernel estimate; see
%\cite[Theorem~3]{Davies87} and, for the full bounded \(C^{1,1}\)-domain
%comparison, \cite[Theorem~1.1 and Remark~1.1]{Z02}. The large-time part of the same estimate, or equivalently boundary regularity and Hopf's lemma for the positive first eigenfunction, gives
%\[
%    f_1(x)\asymp d(x),\qquad x\in D,
%\]
%for the positive \(L^2\)-normalized \(f_1\); see
%\cite[Remark~1.1(b)]{Z02}. This is the boundary scale that later appears in
%the SLND and boundary-layer variance estimates.
%
%It remains to explain the gradient estimate. We use the logarithmic gradient
%bound for the Dirichlet heat kernel on \(C^2\)-domains: for \(0<t\le T\),
%\[
%    |\nabla_x\log P_t(x,y)|
%    \le
%    C\left(\frac1{d(x)}+\frac1{\sqrt t}+\frac{|x-y|}{t}\right),
%\]
%whenever \(P_t(x,y)>0\); see \cite[Theorem~2.1]{Z06}. Hence
%\[
%    |\nabla_xP_t(x,y)|
%    \le
%    P_t(x,y)
%    C\left(\frac1{d(x)}+\frac1{\sqrt t}+\frac{|x-y|}{t}\right).
%\]
%If \(d(x)\le\sqrt t\), we insert the boundary-weighted heat-kernel bound and
%drop the \(y\)-boundary factor:
%\[
%    P_t(x,y)
%    \le
%    C\frac{d(x)}{\sqrt t}t^{-d/2}
%    \exp\left(-\frac{|x-y|^2}{ct}\right).
%\]
%The terms \(d(x)^{-1}\) and \(t^{-1/2}\) then give the factor
%\(t^{-(d+1)/2}\). The term \(|x-y|/t\) is handled by the elementary inequality
%\[
%    \frac{|x-y|}{t}
%    \exp\left(-\frac{|x-y|^2}{ct}\right)
%    \le
%    Ct^{-1/2}\exp\left(-\frac{|x-y|^2}{c't}\right).
%\]
%Thus the desired gradient bound holds in the boundary case. If
%\(d(x)>\sqrt t\), then \(d(x)^{-1}\le t^{-1/2}\); using the unweighted
%Gaussian upper bound and the same absorption inequality for \(|x-y|/t\) gives
%\[
%    |\nabla_xP_t(x,y)|
%    \le
%    Ct^{-(d+1)/2}
%    \exp\left(-\frac{|x-y|^2}{c't}\right).
%\]
%This proves the displayed gradient estimate on every finite time interval.
%
%We finally record the large-time exponential factor. By the spectral
%expansion and Cauchy's inequality, for \(s\ge1\),
%\[
%\begin{aligned}
%    |P_s(z,w)|
%    &\le
%    \e^{-\lambda_1(s-1)}
%    \left(\sum_{n\ge1}\e^{-\lambda_n}f_n(z)^2\right)^{1/2}
%    \left(\sum_{n\ge1}\e^{-\lambda_n}f_n(w)^2\right)^{1/2}  \\
%    &=
%    \e^{-\lambda_1(s-1)}P_1(z,z)^{1/2}P_1(w,w)^{1/2}
%    \le
%    C\e^{-\lambda_1s}.
%\end{aligned}
%\]
%The last step uses the finite-time diagonal bound \(P_1(z,z)\le C\). For every
%fixed \(m\),
%\[
%    \sup_{s\ge1}s^m\e^{-\lambda_1s/2}<\infty .
%\]
%Thus the Gaussian upper bound holds with an additional factor
%\(\e^{-\lambda_1s/2}\) for large \(s\).
%
%For the large-time boundary-weighted estimate, write, for \(s\ge2\),
%\[
%    P_s(x,y)
%    =
%    \int_D\int_D
%    P_{1/2}(x,z)P_{s-1}(z,w)P_{1/2}(w,y)\,\d z\,\d w .
%\]
%The preceding spectral-gap estimate bounds \(P_{s-1}(z,w)\) by
%\(C\e^{-\lambda_1s}\). The boundary-weighted estimate at time \(1/2\) gives
%\[
%    \int_D P_{1/2}(x,z)\,\d z\le C(1\wedge d(x)),
%    \qquad
%    \int_D P_{1/2}(y,w)\,\d w\le C(1\wedge d(y)).
%\]
%Consequently
%\[
%    P_s(x,y)
%    \le
%    C\e^{-\lambda_1s}(1\wedge d(x))(1\wedge d(y)).
%\]
%For \(s\ge1\),
%\[
%    1\wedge d(x)
%    \le
%    Cs^{1/2}\left(1\wedge\frac{d(x)}{\sqrt s}\right),
%\]
%and similarly for \(y\). The two boundary factors and the desired
%\(s^{-d/2}\) prefactor require
%\(s^{d/2+1}\e^{-\lambda_1s}\le C\e^{-\lambda_1s/2}\), so all polynomial
%factors are absorbed into the exponential decay.
%This gives the boundary-weighted large-time bound.
%The remaining bounded interval \(1\le s\le2\) is already covered by the
%finite-time estimate after adjusting constants.
%
%For the large-time gradient estimate, the semigroup identity gives, for
%\(s\ge 3/2\),
%\[
%    \nabla_xP_s(x,y)
%    =
%    \int_D \nabla_xP_{1/2}(x,z)P_{s-1/2}(z,y)\,\d z .
%\]
%The finite-time gradient estimate makes
%\(\int_D|\nabla_xP_{1/2}(x,z)|\,\d z\) uniformly bounded, and the spectral-gap
%bound gives \(P_{s-1/2}(z,y)\le C\e^{-\lambda_1s}\). Hence
%\[
%    |\nabla_xP_s(x,y)|\le C\e^{-\lambda_1s}.
%\]
%Because \(D\) is bounded,
%\(\exp(-|x-y|^2/(cs))\) is bounded from below for \(s\ge1\), and the
%polynomial factor \(s^{(d+1)/2}\) is absorbed by the exponential decay:
%\[
%    \e^{-\lambda_1s}
%    \le
%    Cs^{-(d+1)/2}\e^{-\lambda_1s/2}.
%\]
%This proves the large-time gradient estimate with the extra factor
%\(\e^{-\lambda_1s/2}\); the interval \(1\le s\le3/2\) is again absorbed into
%the finite-time constants.
%\end{proof}

%\begin{lemma}[Quadratic form bound]\label{lem:qform}
%For every open \(D\subset\R^d\), \(F\in L^1(D)\cap L^2(D)\), and \(r>0\),
%\[
%    \inner{F}{P_rF}_{L^2(D)}
%    \le
%    C\left(\|F\|_2^2\wedge r^{-d/2}\|F\|_1^2\right).
%\]
%\end{lemma}
%
%\begin{proof}
%Positivity and self-adjointness give
%\(\inner{F}{P_rF}=\|P_{r/2}F\|_2^2\ge0\), and the \(L^2\)-bound follows from
%semigroup contraction. The \(L^1\)-bound follows for every open domain from
%domain domination by the Euclidean heat kernel:
%\[
%    |P_rF(x)|
%    \le
%    \int_D P_r(x,y)|F(y)|\,\d y
%    \le
%    Cr^{-d/2}\|F\|_1,
%\]
%Hence
%\[
%    |\inner{F}{P_rF}|
%    \le
%    \|F\|_1\|P_rF\|_\infty
%    \le
%    Cr^{-d/2}\|F\|_1^2.
%\]
%Taking the minimum of this estimate and the \(L^2\)-contraction estimate gives
%the stated bound.
%\end{proof}

%\begin{lemma}[Test-function scaling]\label{lem:scaling}
%Let \(A=-\Delta_D\). Let \(\psi\in C_c^\infty(B(0,1/2))\), and put
%\[
%    \psi_{x,r}(y)=r^{-d}\psi\left(\frac{y-x}{r}\right),
%    \qquad B(x,r/2)\subset D.
%\]
%For every \(s\ge0\),
%\[
%    \|A^s\psi_{x,r}\|_{L^2(D)}^2\le C r^{-d-4s}.
%\]
%\end{lemma}
%
%\begin{proof}
%First suppose \(s=m\) is a nonnegative integer. Since \(\psi_{x,r}\) is
%smooth and compactly supported in \(D\), it belongs to \(\operatorname{Dom}(A^m)\)
%for every \(m\) for the Friedrichs Dirichlet realization. The boundary
%condition therefore does not affect the local computation, and
%\[
%    A^m\psi_{x,r}(y)
%    =
%    (-\Delta)^m\psi_{x,r}(y)
%    =
%    (-1)^m r^{-d-2m}
%    (\Delta^m\psi)\left(\frac{y-x}{r}\right).
%\]
%Changing variables \(z=(y-x)/r\), we get
%\[
%\begin{aligned}
%    \|A^m\psi_{x,r}\|_2^2
%    &=
%    r^{-2d-4m}
%    \int_D
%    \left|
%        (\Delta^m\psi)\left(\frac{y-x}{r}\right)
%    \right|^2\,\d y       \\
%    &=
%    r^{-d-4m}
%    \int_{B(0,1/2)}|\Delta^m\psi(z)|^2\,\d z .
%\end{aligned}
%\]
%For a general \(s\ge0\), choose an integer \(k\ge1\) such that
%\(k-1\le s\le k\), and write \(s=(1-\eta)(k-1)+\eta k\). The spectral
%interpolation inequality for powers of \(A\) gives
%\[
%    \|A^s\psi_{x,r}\|_2^2
%    \le
%    \|A^{k-1}\psi_{x,r}\|_2^{2(1-\eta)}
%    \|A^k\psi_{x,r}\|_2^{2\eta}.
%\]
%Indeed, if \(g=\sum_n g_nf_n\), then
%\[
%    \|A^s g\|_2^2
%    =
%    \sum_n \lambda_n^{2s}|g_n|^2
%    =
%    \sum_n
%    \left(\lambda_n^{2(k-1)}|g_n|^2\right)^{1-\eta}
%    \left(\lambda_n^{2k}|g_n|^2\right)^\eta,
%\]
%and Hölder's inequality gives the displayed log-convexity bound. Applying
%this with \(g=\psi_{x,r}\) and using the integer estimates above yields
%\[
%    \|A^s\psi_{x,r}\|_2^2
%    \le
%    C(r^{-d-4(k-1)})^{1-\eta}(r^{-d-4k})^\eta
%    =
%    Cr^{-d-4s}.
%\]
%The interpolation constant is independent of \(x\) and \(r\).
%\end{proof}

%\begin{lemma}[Uniform spectral estimates]\label{lem:spectral}
%Assume \(D\subset\R^d\) is bounded and \(C^2\), and
%\eqref{alpha} holds. For \(\Lambda>0\) and \(z\in D\),
%\[
%    \sum_{\lambda_n\le\Lambda}\lambda_n^{-\alpha}f_n(z)^2
%    \le C\Lambda^{d/2-\alpha},
%\]
%\[
%    \sum_{\lambda_n>\Lambda}\lambda_n^{-(\alpha+1)}f_n(z)^2
%    \le C\Lambda^{d/2-\alpha-1},
%\]
%and
%\[
%    \sum_{\lambda_n\le\Lambda}\lambda_n^{-(\alpha+1)}|\nabla f_n(z)|^2
%    \le C\Lambda^{d/2-\alpha}.
%\]
%If \(\Lambda<\lambda_1\), the first and third sums are empty. The second is
%bounded by its value at \(\lambda_1\), while
%\(\Lambda^{d/2-\alpha-1}\ge\lambda_1^{d/2-\alpha-1}\) because the exponent is
%negative. Enlarging \(C\) therefore makes all displays valid for every
%\(\Lambda>0\).
%\end{lemma}
%
%\begin{proof}
%The Gamma-integral identities
%\[
%    \lambda^{-\alpha}
%    =
%    \frac1{\Gamma(\alpha)}
%    \int_0^\infty t^{\alpha-1}\e^{-\lambda t}\,\d t,
%    \qquad
%    \lambda^{-(\alpha+1)}
%    =
%    \frac1{\Gamma(\alpha+1)}
%    \int_0^\infty t^\alpha\e^{-\lambda t}\,\d t
%\]
%show why heat-kernel diagonal estimates control the weighted spectral sums.
%For the truncated pointwise estimates below, it is convenient to use the
%equivalent dyadic form of this principle.
%
%The diagonal heat-kernel estimate gives the local spectral counting bound.
%Indeed, for \(M>0\),
%\[
%    \sum_{\lambda_n\le M}f_n(z)^2
%    \le
%    e\sum_n\e^{-\lambda_n/M}f_n(z)^2
%    =
%    eP_{1/M}(z,z)
%    \le
%    CM^{d/2}.
%\]
%For the first weighted sum, decompose \((0,\Lambda]\) into the dyadic shells
%\[
%    (2^{-j-1}\Lambda,2^{-j}\Lambda],
%    \qquad j\ge0 .
%\]
%On the \(j\)-th shell,
%\(\lambda_n^{-\alpha}\le C(2^j\Lambda^{-1})^\alpha\), and the spectral
%counting bound gives
%\[
%\begin{aligned}
%    \sum_{\lambda_n\le\Lambda}\lambda_n^{-\alpha}f_n(z)^2
%    &\le
%    C\sum_{j=0}^\infty
%    (2^j\Lambda^{-1})^\alpha
%    (2^{-j}\Lambda)^{d/2}       \\
%    &=
%    C\Lambda^{d/2-\alpha}
%    \sum_{j=0}^\infty 2^{-j(d/2-\alpha)}
%    \le C\Lambda^{d/2-\alpha},
%\end{aligned}
%\]
%because \(\alpha<d/2\).
%
%For the tail, use dyadic shells \((2^j\Lambda,2^{j+1}\Lambda]\), \(j\ge0\).
%Then
%\[
%\begin{aligned}
%    \sum_{\lambda_n>\Lambda}\lambda_n^{-(\alpha+1)}f_n(z)^2
%    &\le
%    C\sum_{j=0}^\infty
%    (2^j\Lambda)^{-(\alpha+1)}
%    (2^{j+1}\Lambda)^{d/2}       \\
%    &\le
%    C\Lambda^{d/2-\alpha-1}
%    \sum_{j=0}^\infty 2^{j(d/2-\alpha-1)}
%    \le C\Lambda^{d/2-\alpha-1},
%\end{aligned}
%\]
%because \(\alpha>d/2-1\).
%
%It remains to prove the gradient-weighted sum. The heat-kernel expansion and
%orthonormality imply
%\[
%\begin{aligned}
%    \sum_n\e^{-2\lambda_n/\Lambda}|\nabla f_n(z)|^2
%    &=
%    \int_D|\nabla_zP_{\Lambda^{-1}}(z,w)|^2\,\d w        \\
%    &\le
%    C\Lambda^{d+1}
%    \int_D
%    \exp\left(-\frac{2|z-w|^2\Lambda}{c}\right)\,\d w    \\
%    &\le C\Lambda^{d/2+1},
%\end{aligned}
%\]
%where the gradient heat-kernel estimate from
%Lemma~\ref{lem:heat} was used in the last line. Therefore
%\[
%    \sum_{\lambda_n\le M}|\nabla f_n(z)|^2
%    \le
%    \e^2\sum_n\e^{-2\lambda_n/M}|\nabla f_n(z)|^2
%    \le C M^{d/2+1},
%\]
%because \(\e^{-2\lambda_n/M}\ge\e^{-2}\) on \(\{\lambda_n\le M\}\).
%Applying the same dyadic decomposition as in the first estimate, but with
%weight \(\lambda_n^{-(\alpha+1)}\) and gradient counting exponent
%\(d/2+1\), gives
%\[
%\begin{aligned}
%    \sum_{\lambda_n\le\Lambda}
%    \lambda_n^{-(\alpha+1)}|\nabla f_n(z)|^2
%    &\le
%    C\sum_{j=0}^\infty
%    (2^j\Lambda^{-1})^{\alpha+1}
%    (2^{-j}\Lambda)^{d/2+1}   \\
%    &=
%    C\Lambda^{d/2-\alpha}
%    \sum_{j=0}^\infty 2^{-j(d/2-\alpha)}
%    \le C\Lambda^{d/2-\alpha}.
%\end{aligned}
%\]
%This completes the proof of all three estimates.
%\end{proof}

%\begin{lemma}[Metric entropy for the parabolic metric]\label{lem:entropy}
%Let \(Q\) be defined by \eqref{eq:Q}. For every \(T>0\) and
%\(\delta>0\),
%\[
%    N_\rho([0,T]\times D_\delta,r)\le Cr^{-Q},
%    \qquad 0<r<1.
%\]
%The constant may depend on \(T\), \(D\), and \(\delta\); only the exponent
%\(Q\) is used in the regularity and small ball arguments.
%More generally, if \(B\subset\R^d\) is bounded and \(0<a<T<\infty\), then
%\[
%    N_\rho([a,T]\times B,r)
%    \le C\bigl(1+(T-a)r^{-1/\theta}\bigr)
%       \bigl(1+N_E(B,c r^{1/(2\theta)})\bigr),
%    \qquad 0<r<1,
%\]
%where \(N_E(B,s)\) denotes the covering number of \(B\) by Euclidean
%balls of radius \(s\). Placing \(B\) in a fixed bounding box gives
%\(N_\rho([a,T]\times B,r)\le C_{a,T,B}r^{-Q}\).
%Moreover, for every fixed \(z_0\) and \(0<r\le R\le1\),
%\[
%    N_\rho(B_\rho(z_0,R),r)\le C\left(\frac Rr\right)^Q .
%\]
%\end{lemma}
%
%\begin{proof}
%Cover time by intervals of length \(r^{1/\theta}\) and space by Euclidean
%balls of radius \(r^{1/(2\theta)}\). This requires at most
%\[
%    Cr^{-1/\theta}r^{-d/(2\theta)}=Cr^{-Q}
%\]
%sets. The local estimate is the same argument applied to a time interval of
%length \(R^{1/\theta}\) and a spatial ball of radius \(R^{1/(2\theta)}\).
%The same covering argument applies to any bounded spatial set \(B\), after
%placing \(B\) in a fixed Euclidean box.
%\end{proof}

\bibliographystyle{plain}
\bibliography{SHE}


\end{document}
